% !TeX root = ../Book3.tex
\BookPart{整函数与亚纯函数}{b3:part:04}
% 本编导读（章级标题，不计入编号）
\BookPartGuide

整函数的局部 Taylor 级数可以在全平面继续，但全局行为仍有巨大差别。多项式、指数函数与带稀疏零点的无穷乘积，都需要比局部收敛半径更细的尺度。本编用最大模、增长阶和典范乘积描述这种差别，再以 Gamma 函数检验整套方法。
\[
\begin{tikzcd}[row sep=large]
\text{第十九章：增长理论}\arrow[d]\\
\text{第二十章：零点与因子分解}\arrow[d]\\
\text{第二十一章：Gamma 与 Mellin 变换}
\end{tikzcd}
\]
第十八章的 Jensen 公式直接进入增长与零点的比较；第二十章的正弦乘积又进入 Gamma 反射公式。Gamma 的 Euler 积分和对数凸性可以先独立阅读，随后再看其乘积表示。Mellin 变换则是通往第六编的另一出口。
\ProseParagraphBreak
本编的证明常要在两个方向之间切换。一方面，给定函数的增长限制了零点的密度；另一方面，给定零点序列可以构造整函数，但还留有无零点指数因子的自由。Hadamard 分解只有在有限增长假设下才把这份自由压缩成一个多项式。读到每个乘积时，应同时辨认收敛方式、零点重数与正规化常数。
\ProseParagraphBreak
第二十一章中特别区分实正轴的 Stirling 渐近与一般复扇形渐近。积分变换的反演也保留明确的可积条件。公式的形式虽然简短，其使用范围来自前面实际证明的估计。

% !TeX root = ../../Book3.tex
\chapter{整函数的增长}
\label{b3:ch:19}
\BookChapterNumber{b3:ch:19}
\begin{chapterabstract}
整函数在无穷远的大小可以从圆周最大模、Taylor 系数和零点计数三个方面观察。本章先用三圆定理建立对数半径上的凸性，并比较最大模与最大单项；随后定义增长阶与型，证明系数判据及导数、坐标变化和多项式复合下的增长规律。Jensen 公式把增长界转成零点收敛指数与密度的限制。最后在扇形和条带中用增长条件控制无穷远边界，分别建立严格次临界和临界最小型的 Phragmén--Lindelöf 比较，以及允许竖向增长的三线估计。
\end{chapterabstract}

\begin{learninggoals}
\begin{itemize}
\item 区分圆周最大模与平均对数，用三圆不等式及其等号情形识别径向增长的限制。
\item 比较最大模与 Taylor 最大项，追踪主导次数随半径的变化，并由全局幂增长判别多项式。
\item 正确使用阶与型中的上极限，通过系数衰减计算或构造给定增长，保留零阶、最小型和最大型的边界情形。
\item 说明导数、仿射变换与多项式复合如何改变阶型，识别加法与乘法中可能出现的抵消。
\item 按重数计算零点收敛指数，区分指数、临界幂和与密度常数提供的不同资料。
\item 在无界区域比较中逐一控制有限边界、顶点与远处边界，辨明严格指数条件、临界最小型条件和条带宽度的作用。
\end{itemize}
\end{learninggoals}

多项式的次数可以由 \(M(r)\) 的幂增长读出，而 \(e^{z^2}\) 的最大模是 \(e^{r^2}\)，必须先取对数才能看见幂次。另一方面，\(e^{\mathrm iz^2}\) 在实轴和虚轴上都具有模一值，却在两轴之间的某些方向快速增长。一个径向尺度如何从所有方向中取得，几条射线上的信息又在什么条件下足以约束整个平面，是本章始终需要分清的两个问题。

\ProseParagraphBreak
\secref{b3:sec:1901}先建立最大模的基本结构。全局幂界结合 Cauchy 系数估计使高次系数消失；三圆定理则用环域上的调和函数插值两条圆周的模界。等号情形还要利用复函数的刚性，才能从径向直线恢复单项式。Taylor 最大项提供另一个可算的径向量，其对数是一族整数斜率直线的上包络，但比较它与整个级数时仍须控制其余各项。

\ProseParagraphBreak
\secref{b3:sec:1902}把这些估计组织成阶与型。阶确定 \(\log M(r)\) 的幂尺度，型比较这个尺度下的系数；零阶仍容许超越整函数，给定正阶也容许零型和最大型。系数判据的两个方向分别来自最优选择 Cauchy 半径和控制模级数尾部。读者可以先跟算基本例子，再用这些判据处理稀疏系数；指标中有大量零系数，并不会使非零子列的上极限失去作用。

\ProseParagraphBreak
\secref{b3:sec:1903}回到零点。收敛指数可以由幂和或整数计数的增长读出，临界指数处的幂和则还取决于更细的修正。Jensen 公式使最大模限制零点数量，并把给定密度转成型的必要下界。无零点因子能够改变增长，所以这些比较的方向和前提应当保留；仅知道一个计数上界，不能反推出函数的最大模上界。

\ProseParagraphBreak
\secref{b3:sec:1904}处理无界区域的边界比较。严格次临界增长可以由更高次的衰减指数压住；达到临界指数时，这个余量消失，需要最小型条件及另一种调和屏障。有限多条射线把平面分割后，每个扇形都须满足自己的角宽条件。条带中的三线比较同样需要控制横向截断边界，允许的竖向增长与条带宽度相联系。边界模的比较与连续零边界的局部唯一性承担不同任务，不应把增长条件不加区分地附到所有唯一性问题上。

\ProseParagraphBreak
本章使用 Cauchy 公式与系数估计、次调和比较，以及\chapref{b3:ch:18}的 Jensen 公式和零点计数。Levin 的《Lectures on Entire Functions》\cite{Levin1996Lectures}可作为继续学习增长方向与零点分布的参考。本章先建立后续整函数乘积理论所需的增长尺度；一般指标理论和给定零点的乘积构造不作为本章证明的前提。

% !TeX root = ../../Book3.tex
\section{最大模}
\label{b3:sec:1901}

多项式与指数函数都在全平面全纯，却有不同的无穷远尺度。最大模把方向各异的增长压成一个径向量，三圆定理则限制这个量怎样随半径变化。

\subsection{\texorpdfstring{\(M(r,f)\)}{M(r,f)}}
\label{b3:subsec:190101}

对整函数 \(f\)，定义最大模 \(M(r,f)=\max_{|z|=r}|f(z)|\)。最大模原理使它等于闭圆盘上的最大模，并随 \(r\) 非减。若 \(f\) 非常值，则 \(M\) 严格递增：两半径相同最大值会在大圆盘内部取得最大模，迫使函数常数。

多项式 \(a_dz^d+\cdots+a_0\) 有 \(M(r)=|a_d|r^d(1+o(1))\)；\(e^z\) 有 \(M(r)=e^r\)。沿负实轴的指数趋零，与它的最大模指数增长并不矛盾。

\subsection{最大模}
\label{b3:subsec:190102}

\begin{proposition}\label{b3:prop:19-polynomial-growth}
若整函数满足 \(M(r,f)\le C(1+r)^N\)，则为次数不超过 \(\lfloor N\rfloor\) 的多项式。
\end{proposition}
\begin{proof}
Cauchy 估计给 \(|f^{(k)}(0)|\le k!C(1+r)^N/r^k\)。当整数 \(k>N\) 时令 \(r\to\infty\)，得到导数零；Taylor 展开因而有限。反向显然由有限项三角不等式成立。
\end{proof}
这里控制的是所有方向，而不是一条射线。函数 \(e^{-z}\) 在正实轴有界，却不是多项式。

\subsection{Hadamard 三圆定理}
\label{b3:subsec:190103}

\term[hadamard-sanyuan]{Hadamard 三圆定理}[Hadamard three-circles theorem]可由次调和比较直接得到。
\begin{theorem}[Hadamard 三圆]\label{b3:thm:19-three-circles}
若 \(f\) 在环域 \(r_1\le|z|\le r_3\) 邻域全纯且不恒零，则 \(r_1<r_2<r_3\) 时
\[
\log M(r_2)\le
\frac{\log(r_3/r_2)}{\log(r_3/r_1)}\log M(r_1)
+\frac{\log(r_2/r_1)}{\log(r_3/r_1)}\log M(r_3).
\]
\end{theorem}
\begin{proof}
选择 \(A,B\) 使调和函数 \(A\log|z|+B\) 在两边界等于相应最大对数。函数 \(\log|f|-A\log|z|-B\) 次调和，两边界均非正，最大值原理给内部非正。在半径 \(r_2\) 取最大值，线性插值就是题式。
\end{proof}

\subsection{对数凸性}
\label{b3:subsec:190104}

令 \(t=\log r\)，三圆定理说 \(t\mapsto\log M(e^t)\) 凸。其割线斜率随半径向外不减；这是增长约束，而不是通常变量 \(r\) 下的凸性。

例如单项式 \(z^d\) 给严格直线 \(dt\)；指数给 \(e^t\)。如果某非零整函数在所有大半径的对数斜率都不超过整数 \(d\)，则 \(M(r)\le Cr^d\)，上一命题将它识别为多项式。这使一个径向几何性质决定整个函数的代数结构。


\ProseParagraphBreak

三圆定理也能在同一半径上比较几何平均与最大值。Jensen 的圆周平均描述零点总影响，最大模则取最有利增长的方向；前者可能恒定，后者却增长很快。对 \(e^z\)，圆周对数平均零，而 \(\log M(r)=r\)。因此由平均对数的凸性不能直接推出最大模的精确增长，三圆定理另需对边界最大值作调和比较。这两种凸性虽然都以 \(\log r\) 为变量，控制的对象仍不同。

\ProseParagraphBreak

最大模关于半径连续。若 \(r_j\to r>0\)，在共同紧环带上 \(f\) 一致连续，将 \(r_je^{\mathrm{i}t}\) 与 \(re^{\mathrm{i}t}\) 比较，函数值差一致趋零，故两圆周最大值之差趋零。对非恒零整函数 \(M(r)>0\)，对数也连续。三圆不等式因此给一个通常意义的连续凸函数，而非只在排除零点的半径上定义的对象。圆周即使碰到函数零点，最大模仍由其他位置取得，不需要移开半径。

\ProseParagraphBreak

若三圆不等式在一个中间半径取得等号，比较函数 \(\log|f|-A\log|z|-B\) 在环域内部取得最大值零，故恒为零。于是 \(|f(z)|=e^B|z|^A\)，\(f\) 在环域无零点。局部对数微分给 \(f'/f=A/z\)，沿绕原点的圆周积分，辐角原理使 \(A\) 必为整数；继而 \(f(z)=cz^A\) 于环域。若 \(f\) 还在原点全纯，整数 \(A\) 必非负。三圆等号的刚性来自次调和强最大值原理，不能只由一个实凸函数的一次等号就辨认出复函数的形式。

\ProseParagraphBreak

最大模还可与 Taylor 展开的最大单项比较。对 \(f(z)=\sum c_nz^n\)，定义 \(\mu(r)=\max_n|c_n|r^n\)。整性使项随 \(n\to\infty\) 趋零，所以最大值确实达到。Cauchy 系数界给 \(\mu(r)\le M(r)\)；另一方面，对 \(0<r<R\)，
\[
 M(r)\le\sum_n|c_n|r^n
       \le\mu(R)\sum_n(r/R)^n
       =\frac{\mu(R)}{1-r/R}.
\]
这两边估计说明，模级数中的相消可能影响某个方向的函数值，却不能使整个圆周上的最大模比所有单项都小。取 \(R=2r\)，两种径向量具有相同的增长阶；当讨论有限正阶的型时，再取 \(R=(1+\varepsilon)r\) 并让 \(\varepsilon\downarrow0\)，它们的型也相同。
\ProseParagraphBreak
从对数坐标看，\(\log\mu(e^t)=\max_n(\log|c_n|+nt)\)，是斜率为非负整数的一族直线的上包络，因而凸。取得最大值的指标随 \(t\) 非减：若在 \(t_1<t_2\) 分别由 \(n_1,n_2\) 取到最大值，把两次比较不等式相加便得 \((n_2-n_1)(t_2-t_1)\ge0\)。这使 Taylor 展开的主导项随半径逐渐移向高次幂。若主导指标最终有共同上界，则所有足够高次系数都必须为零，否则该项的斜率更大，最终会超过已有的有限条直线。多项式与超越整函数的差别，由此再次表现为无穷远是否不断出现新的主导次数。

% !TeX root = ../../Book3.tex
\section{增长阶与型}
\label{b3:sec:1902}

多项式的次数可以从 \(\log M(r)/\log r\) 读出，但这个量对每个超越整函数都趋于无穷。若 \(M(r)\) 大致像 \(e^{r^s}\)，应当转而比较 \(\log M(r)\) 与 \(r^s\)。阶用来确定这里的幂指数，型则在同一幂尺度上保留前面的系数。先把这两层比较的量词说清楚，再研究它们怎样由导数、复合和 Taylor 系数决定。

\subsection{阶}
\label{b3:subsec:190201}

\Needspace{6\baselineskip}
\begin{definition}\label{b3:def:19-order}
设 \(f\) 为非常值整函数，\(M(r,f)=\max_{|z|=r}|f(z)|\)。其\term[zengzhang-jie]{增长阶}[order of growth]定义为
\[
 \rho(f)\coloneqq\limsup_{r\to\infty}
 \frac{\log\log M(r,f)}{\log r}\in[0,\infty].
\]
常数函数，包括零函数，约定增长阶为零。增长阶有限的整函数称为\term[youxianjie-zhenghanshu]{有限阶整函数}[entire function of finite order]。
\end{definition}
\symidx{\(\rho(f)\)：整函数的增长阶}
非常值时 \(M(r,f)\to\infty\)，所以定义中的双对数只需在充分大半径使用。其上极限不会为负：取某个 \(c_k\ne0\)、\(k\ge1\)，便有 \(M(r,f)\ge|c_k|r^k\)，从而双对数的商的下极限不小于零。采用上极限，是为了容许增长在不同半径上出现变化，不预先假定存在一个精确的渐近式。

\ProseParagraphBreak

阶最常用的含义是可用哪些指数上界。严格大于阶的指数总能给出上界；严格小于阶的指数则无论乘上多大的常数都不够。

\Needspace{6\baselineskip}
\begin{proposition}[阶与指数增长尺度]\label{b3:prop:19-growth-scales}
设 \(f\) 为整函数，其增长阶为 \(\rho\in[0,\infty]\)。对每个实数 \(s>\rho\)，存在 \(C_s>0\)，使
\[
 M(r,f)\le C_s e^{r^s},\quad r\ge0.
\]
对每个 \(0<s<\rho\) 及任意 \(A,C>0\)，存在任意大的半径 \(r\) 使 \(M(r,f)>C e^{Ar^s}\)。因此
\[
 \rho=\inf\bigl\{\,s>0\mid
 \exists A,C>0\ \forall r\ge0,\ M(r,f)\le C e^{Ar^s}\,\bigr\},
\]
其中空集的下确界约定为无穷。
\end{proposition}
\begin{proof}
常数函数直接满足结论，以下设 \(f\) 非常值。若 \(s>\rho\)，上极限的定义给充分大半径上的
\[
 \log\log M(r,f)\le s\log r,
 \quad\text{即}\quad \log M(r,f)\le r^s.
\]
再用 \(M\) 在有限半径区间的连续性扩大 \(C_s\)，得到全半径的估计。

若 \(0<s<\rho\)，选 \(s'\) 使 \(s<s'<\rho\)；\(\rho=\infty\) 时也可如此选择。存在趋于无穷的半径序列 \(r_j\)，使 \(\log\log M(r_j,f)>s'\log r_j\)。故
\[
 \frac{\log M(r_j,f)}{r_j^s}>r_j^{s'-s}\longrightarrow\infty.
\]
对给定的 \(A,C\)，充分大的 \(j\) 便满足 \(\log M(r_j,f)>A r_j^s+\log C\)。最后，任何上界 \(M(r,f)\le C e^{Ar^s}\) 都使双对数的上极限不超过 \(s\)；结合已经证明的严格阈值，得到下确界公式。
\end{proof}
下确界不一定在 \(s=\rho\) 处取得。知道 \(\rho<\infty\)，只保证每个 \(\varepsilon>0\) 都能给出 \(M(r)\le C_\varepsilon e^{r^{\rho+\varepsilon}}\)。能否在临界指数 \(\rho\) 上取得有限常数，需要再看型。

\subsection{型}
\label{b3:subsec:190202}

\Needspace{6\baselineskip}
\begin{definition}\label{b3:def:19-type}
设整函数 \(f\) 的增长阶为 \(\rho\in(0,\infty)\)，\(M(r,f)=\max_{|z|=r}|f(z)|\)。其在自身增长阶上的\term[zengzhang-xing]{型}[type]定义为
\[
 \tau(f)\coloneqq\limsup_{r\to\infty}
 \frac{\log M(r,f)}{r^\rho}\in[0,\infty].
\]
\(\tau=0\)、\(0<\tau<\infty\)、\(\tau=\infty\) 分别称为\term[zuixiaoxing]{最小型}[minimal type]、\term[zhengchangxing]{正常型}[normal type]、\term[zuidaxing]{最大型}[maximal type]。
\end{definition}
\symidx{\(\tau(f)\)：整函数在自身有限正阶上的型}
这里的尺度必须是函数自身的阶。在更大指数上，这个上极限自动为零；在更小的正指数上，它自动为无穷，均不能细分同阶函数。零阶不按上述公式作统一型分类：非常值多项式与超越零阶函数都有 \(\log M(r)\to\infty\)，除以 \(r^0=1\) 后无法区分它们。无穷阶也不属于此定义的范围。

\ProseParagraphBreak

与阶的阈值类似，若 \(\tau<\infty\)，则对每个 \(b>\tau\)，存在 \(C_b>0\) 使
\[
 M(r,f)\le C_b e^{b r^\rho},\quad r\ge0.
\]
若 \(0<b<\tau\)，则 \(M(r,f)e^{-b r^\rho}\) 沿某个趋于无穷的半径序列趋于无穷。前者来自上极限的最终上界，后者可选 \(b'\in(b,\tau)\)，在满足 \(\log M(r)>b'r^\rho\) 的序列上验证。因而 \(\tau\) 是所有能给出上述全半径上界的正数 \(b\) 的下确界；\(\tau=\infty\) 时这一集合为空。

\ProseParagraphBreak

临界常数本身仍未必可用。例如 \(f(z)=(1+z)e^z\) 满足
\[
 M(r,f)=(1+r)e^r.
\]
上界来自 \(|1+z|\le1+r\)、\(\operatorname{Re}z\le r\)，而在正实点 \(z=r\) 同时取得等号。于是 \(\rho=1\)、\(\tau=1\)，却不存在统一的 \(C\) 使 \(M(r,f)\le C e^r\)。因此从型的定义得到的 \(\tau+\varepsilon\) 通常不能直接换成 \(\tau\)。

\ProseParagraphBreak

前节的最大项也能读出这两个量。比较型时必须让两个半径的比可以任意接近一；只比较 \(r\) 和 \(2r\)，会留下一个无法消去的常数因子。

\Needspace{6\baselineskip}
\begin{proposition}[最大项给出相同的阶与型]\label{b3:prop:19-maximum-term-growth}
设 \(f(z)=\sum_{n=0}^{\infty}c_nz^n\) 为非常值整函数，\(\mu(r,f)=\max_{n\ge0}|c_n|r^n\)。则
\[
 \rho(f)=\limsup_{r\to\infty}
 \frac{\log\log\mu(r,f)}{\log r}.
\]
若 \(0<\rho(f)<\infty\)，则其型还满足
\[
 \tau(f)=\limsup_{r\to\infty}
 \frac{\log\mu(r,f)}{r^{\rho(f)}}.
\]
两个公式均容许其所定义的量为零或无穷。
\end{proposition}
\begin{proof}
某个正次系数非零，所以 \(\mu(r,f)\to\infty\)，两个公式在大半径有意义。固定 \(\kappa>1\)，由\cref{b3:prop:19-maximum-term}，
\[
 \mu(r,f)\le M(r,f)
 \le\frac{\kappa}{\kappa-1}\mu(\kappa r,f).
\]
乘以一个正常数只使双对数改变 \(o(1)\)，把半径 \(r\) 换成 \(\kappa r\) 则有 \(\log(\kappa r)/\log r\to1\)。两边取双对数的上极限，便得到阶相同，包括零阶和无穷阶。

设共同的阶为 \(\rho\in(0,\infty)\)，并把最大项公式右端记为 \(\tau_\mu\)。取一次对数并除以 \(r^\rho\)，上述不等式给
\[
 \tau_\mu\le\tau(f)\le\kappa^\rho\tau_\mu.
\]
若 \(\tau_\mu\) 有限，让 \(\kappa\downarrow1\) 得到相等；若 \(\tau_\mu=\infty\)，左侧已经迫使 \(\tau(f)=\infty\)。
\end{proof}

\subsection{基本例子}
\label{b3:subsec:190203}

先看 \(f(z)=e^{az^m}\)，其中 \(a\ne0\)、\(m\ge1\) 为整数。在 \(|z|=r\) 上有 \(\operatorname{Re}(az^m)\le|a|r^m\)，又能选取一个角度使 \(az^m\) 为正实数，因而
\[
 M(r,f)=e^{|a|r^m},\quad \rho(f)=m,\quad \tau(f)=|a|.
\]
这里必须选到使实部最大的位置；只看实轴，在 \(a\) 为纯虚数时便会遗漏全部增长。

\ProseParagraphBreak

三角函数也属于一阶正常型。由指数表达式，
\[
 |\sin z|\le\frac{e^{\operatorname{Im}z}+e^{-\operatorname{Im}z}}2\le e^r,
 \quad
 |\cos z|\le e^r,\quad |z|=r.
\]
在虚轴上则有 \(|\sin(\mathrm{i}r)|=\sinh r\)、\(|\cos(\mathrm{i}r)|=\cosh r\)。两者均为 \(\tfrac12e^r(1+o(1))\)，所以两种最大模都满足 \(\log M(r)=r+O(1)\)，得到阶一、型一。它们在实轴有界与此并不矛盾。

\ProseParagraphBreak

多项式和迭代指数给出两个端点。若 \(P\) 为 \(d\ge1\) 次多项式，\(\log M(r,P)=d\log r+O(1)\)，从而 \(\rho(P)=0\)，无论次数多大。另一方面，\(f(z)=e^{e^z}\) 满足
\[
 |f(z)|=e^{\operatorname{Re}(e^z)}\le e^{e^r},\quad |z|=r,
\]
在 \(z=r\) 取得等号。因此 \(\log\log M(r,f)=r\)，阶为无穷。

\ProseParagraphBreak

零阶还包含超越整函数。考虑
\[
 F(z)=\sum_{n=0}^{\infty}e^{-n^2}z^n.
\]
系数的 \(n\) 次根为 \(e^{-n}\)，趋于零，所以收敛半径无穷。所有系数为正，三角不等式的上界在 \(z=r\) 取得，于是令 \(t=\log r\) 后有
\[
 M(r,F)=F(r)
 =\sum_{n=0}^{\infty}e^{-n^2+nt}
 =e^{t^2/4}\sum_{n=0}^{\infty}e^{-(n-t/2)^2}.
\]
当 \(t\ge1\) 时，选离 \(t/2\) 最近的非负整数 \(n_0\)，其距离不超过 \(1/2\)，便得到最后一和的下界 \(e^{-1/4}\)。对其他指标按 \(|n-n_0|\) 分组，则该和至多为
\[
 1+2\sum_{k=1}^{\infty}e^{-(k-1/2)^2},
\]
这是与 \(t\) 无关的有限常数。因此
\[
 \log M(r,F)=\frac14(\log r)^2+O(1).
\]
它超过每个固定的 \(d\log r\)，却小于每个 \(r^\varepsilon\) 的最终尺度；所以 \(F\) 超越而阶为零。最大模的对数斜率仍趋于无穷，符合前节的凸斜率判别。

\ProseParagraphBreak

一般的代数运算先给上界。若 \(f,g\) 为整函数，则
\[
 M(r,f+g)\le M(r,f)+M(r,g),\quad
 M(r,fg)\le M(r,f)\cdot M(r,g).
\]
结合指数尺度判据可得
\[
 \rho(f+g)\le\max\{\rho(f),\rho(g)\},\quad
 \rho(fg)\le\max\{\rho(f),\rho(g)\}.
\]
例如当两阶有限时，对每个大于二者的 \(s\)，和、积都满足 \(C e^{A r^s}\) 型上界，再让 \(s\) 下降即可；若最大阶无穷，上界本身自动成立。若两函数阶不同，和的阶必等于较大者，否则把高阶函数写成和减去低阶函数，将与刚才的上界矛盾。同阶相消则可能使阶下降：\(e^z+(-e^z+1)=1\)，而乘积 \(e^ze^{-z}=1\) 也说明乘法一般没有对应的等式。

\ProseParagraphBreak

平移和缩放改变圆心与半径，却不会任意改变增长尺度。用闭圆盘最大值表示 \(M\)，可以把这种变化写成精确的包含关系。

\Needspace{6\baselineskip}
\begin{proposition}[仿射变量变换]\label{b3:prop:19-affine-growth}
设 \(f\) 为整函数，\(a,b\in\mathbb C\)、\(a\ne0\)，并令 \(g(z)=f(az+b)\)。则
\[
 \rho(g)=\rho(f).
\]
若 \(\rho(f)=\rho\in(0,\infty)\)，则
\[
 \tau(g)=|a|^\rho\tau(f),
\]
其中型容许为零或无穷。
\end{proposition}
\begin{proof}
仿射映射把 \(|z|\le r\) 送到以 \(b\) 为中心、半径 \(|a|r\) 的闭圆盘。当 \(|a|r>|b|\) 时，这个圆盘包含以原点为中心、半径 \(|a|r-|b|\) 的闭圆盘，并包含于半径 \(|a|r+|b|\) 的闭圆盘。因此
\[
 M(|a|r-|b|,f)\le M(r,g)\le M(|a|r+|b|,f).
\]
常数函数的结论直接成立。若 \(f\) 非常值，令 \(s=|a|r\pm|b|\)，则每种代换都遍历一个充分大的半径区间，且 \(\log s/\log r\to1\)。对夹逼取双对数的上极限，得到阶不变，包括零阶和无穷阶。

若共同的阶为 \(\rho\in(0,\infty)\)，则
\[
 \frac{\log M(s,f)}{r^\rho}
 =\left(\frac{s}{r}\right)^\rho
   \frac{\log M(s,f)}{s^\rho},
 \quad \frac{s}{r}\longrightarrow|a|.
\]
夹逼的两端因此有同一上极限 \(|a|^\rho\tau(f)\)，得到型的公式。
\end{proof}
特别地，平移不改变阶与型。缩放的影响则保留在型的因子 \(|a|^\rho\) 中。

\ProseParagraphBreak

把仿射映射换成 \(d\) 次多项式，半径 \(r\) 的主尺度就变为 \(r^d\)。上界只需三角不等式；为了得到反向估计，还要保证多项式的像覆盖一个足够大的圆盘。

\Needspace{6\baselineskip}
\begin{proposition}[多项式复合的增长]\label{b3:prop:19-polynomial-composition}
设 \(f\) 为整函数，\(P(z)=a_dz^d+\cdots+a_0\) 为 \(d\ge1\) 次多项式，\(a_d\ne0\)。则
\[
 \rho(f\circ P)=d\rho(f),
\]
其中零阶和无穷阶按扩展数值解释。若 \(\rho(f)=\rho\in(0,\infty)\)，则在 \(f\circ P\) 自身的阶 \(d\rho\) 上，
\[
 \tau(f\circ P)=|a_d|^\rho\tau(f).
\]
\end{proposition}
\begin{proof}
固定 \(0<q<|a_d|<Q\)。当 \(r\) 充分大时，三角不等式给
\[
 |P(z)|\le Qr^d,\quad |z|\le r.
\]
另一方面，对任意 \(|w|\le q r^d\)，在 \(|z|=r\) 上都有
\[
 |P(z)-w-a_dz^d|
 \le\sum_{k<d}|a_k|r^k+q r^d
 <|a_d|r^d.
\]
\(P-w\) 与 \(a_dz^d\) 在闭圆盘邻域全纯，且满足严格的边界比较，故\cref{b3:thm:03-rouche}使二者在盘内有相同的零点数 \(d\)。特别地，每个这样的 \(w\) 都是某个 \(|z|<r\) 的像。于是
\[
 M(qr^d,f)\le M(r,f\circ P)\le M(Qr^d,f).
\]

若 \(f\) 为常数，结论直接成立。其余情形将 \(s=q r^d\) 或 \(s=Q r^d\) 代入，有 \(\log s/\log r\to d\)，而 \(s\) 遍历所有充分大的半径。夹逼两端的双对数商都以上极限 \(d\rho(f)\) 为界，故得到阶的等式。

设 \(0<\rho<\infty\)。同一夹逼除以 \(r^{d\rho}\) 后给
\[
 q^\rho\tau(f)\le\tau(f\circ P)\le Q^\rho\tau(f).
\]
令 \(q\uparrow|a_d|\)、\(Q\downarrow|a_d|\)，即得有限型的公式；若 \(\tau(f)=\infty\)，左边已经迫使复合函数也为无穷型，零型则由右边直接得到。
\end{proof}
因此，多项式的次数决定复合后阶的倍数，首项系数决定型的缩放，低次项只影响较小的半径误差。这一论证同时使用了像的上界与覆盖性质，单凭 \(|P(z)|\le C r^d\) 只能得到单向增长估计。

\subsection{导数增长与系数判据}
\label{b3:subsec:190204}

Cauchy 估计把较大圆周上的函数值转成较小圆周上的导数界；沿线段积分又能从导数恢复函数。二者合起来说明微分不会改变阶，在有限正阶上也不会改变型。

\Needspace{6\baselineskip}
\begin{proposition}[导数的阶与型]\label{b3:prop:19-derivative-order}
设 \(f\) 为整函数，常数函数的阶按零计。则
\[
 \rho(f')=\rho(f).
\]
若 \(\rho(f)=\rho\in(0,\infty)\)，则 \(\tau(f')=\tau(f)\)，包括最小型与最大型。
\end{proposition}
\begin{proof}
常数和一次多项式的结论直接成立，以下设 \(f'\) 非常值。对每个 \(|z|\le r\)，以 \(z\) 为圆心、\(r\) 为半径使用 Cauchy 导数估计，所用圆盘包含于 \(|\zeta|\le2r\)，故
\[
 M(r,f')\le\frac{M(2r,f)}r.
\]
取 \(r>1\)，双对数与固定半径倍增给 \(\rho(f')\le\rho(f)\)。反过来，沿从 \(0\) 到 \(z\) 的线段积分，得到
\[
 |f(z)|\le|f(0)|+|z|M(|z|,f'),
 \quad
 M(r,f)\le|f(0)|+rM(r,f').
\]
因 \(f'\) 非常值，某个正次 Taylor 系数非零，Cauchy 估计使 \(M(r,f')\ge c r^k\)，其中 \(c>0\)、\(k\ge1\)。于是 \(\log r=O(\log M(r,f'))\)，积分估计便给
\[
 \log\log M(r,f)\le\log\log M(r,f')+O(1).
\]
除以 \(\log r\) 并取上极限，得反向阶不等式。整个论证同样适用于零阶与无穷阶。

为比较型，固定 \(\varepsilon>0\)，在以 \(z\) 为中心、半径 \(\varepsilon r\) 的圆周上使用 Cauchy 估计，得到
\[
 M(r,f')\le\frac{M((1+\varepsilon)r,f)}{\varepsilon r}.
\]
取对数、除以 \(r^\rho\) 并取上极限，给出
\[
 \tau(f')\le(1+\varepsilon)^\rho\tau(f).
\]
积分估计在大半径可写为 \(M(r,f)\le2rM(r,f')\)，故 \(\tau(f)\le\tau(f')\)，因为 \(\log(2r)/r^\rho\to0\)。若 \(\tau(f)\) 有限，让 \(\varepsilon\downarrow0\) 得到相等；若 \(\tau(f)=\infty\)，反向不等式已经迫使 \(\tau(f')=\infty\)。
\end{proof}
反复使用命题可知，任何固定阶导数都保留原函数的阶；在有限正阶上还保留型。对整函数的任意全纯原函数应用同一命题，可知原函数的阶也相同，在有限正阶上型也相同。积分常数只改变零次项。

\ProseParagraphBreak

Taylor 系数则给出更直接的判据。Cauchy 估计允许为每个系数选择不同的半径：当最大模近似 \(e^{r^s}\) 时，最有利的半径由指数增长与分母 \(r^n\) 的竞争决定。反过来，从系数估计恢复最大模时，除了控制最大的单项，还必须控制有效项的数量和剩余尾部。

\Needspace{6\baselineskip}
\begin{theorem}[增长的系数判据]\label{b3:thm:19-coefficient-growth}
设 \(f(z)=\sum_{n=0}^{\infty}c_nz^n\) 为超越整函数。则
\[
 \rho(f)=\limsup_{\substack{n\to\infty\\c_n\ne0}}
 \frac{n\log n}{\log(1/|c_n|)}.
\]
非零系数的分母最终为正，公式容许零阶与无穷阶。若 \(\rho=\rho(f)\in(0,\infty)\)，其型还满足
\[
 \limsup_{n\to\infty}n^{1/\rho}|c_n|^{1/n}
 =(e\rho\tau(f))^{1/\rho},
\]
等价地，
\[
 e\rho\tau(f)=\limsup_{n\to\infty}n|c_n|^{\rho/n}.
\]
型的两个公式中，零系数对应的项取零，等式按零值和无穷值的意义解释。
\end{theorem}
\begin{proof}
先证明阶公式。整性使 \(|c_n|^{1/n}\to0\)，所以非零系数最终满足 \(|c_n|<1\)，分母为正。超越性又保证非零系数有无穷多个，因而指定子序列上的上极限有意义。记这个上极限为 \(a\in[0,\infty]\)。

若 \(\rho(f)<\infty\)，取任意 \(s>\rho(f)\)。指数尺度判据给 \(M(r,f)\le C e^{r^s}\)，Cauchy 系数界于是为
\[
 |c_n|\le C e^{r^s}r^{-n},\quad r>0.
\]
为了使右边尽量小，对 \(r^s-n\log r\) 求极小值。导数为零时 \(sr^s=n\)，代入 \(r=(n/s)^{1/s}\) 得
\[
 \log(1/|c_n|)
 \ge\frac ns\log n-\frac ns(1+\log s)-\log C.
\]
右边的主项为 \((n/s)\log n\)，其余为 \(O(n)\)，故 \(a\le s\)。让 \(s\downarrow\rho(f)\)，得到 \(a\le\rho(f)\)；若 \(\rho(f)=\infty\)，这一不等式自动成立。

反过来，若 \(a<\infty\)，取任意 \(s>a\)。上极限的定义使充分大的非零系数满足
\[
 \frac{n\log n}{\log(1/|c_n|)}<s,
 \quad\text{从而}\quad |c_n|\le n^{-n/s}.
\]
零系数也满足最后的界。将不满足此界的有限头部记为多项式 \(P\)，并只考虑 \(r\ge1\)。对尾部的单项，把指标看成正实数 \(x\)，其对数由
\[
 \phi(x)=x\log r-\frac{x}{s}\log x
\]
控制。\(\phi'(x)=0\) 给 \(x=r^s/e\)，且 \(\phi''(x)=-1/(sx)<0\)，故最大值为 \(r^s/(es)\)。在 \(n\le(2r)^s\) 的部分，项数至多 \(1+(2r)^s\)；在 \(n>(2r)^s\) 的部分，\(r n^{-1/s}<1/2\)，其和至多为 \(\sum_{n\ge1}2^{-n}=1\)。于是
\[
 M(r,f)\le M(r,P)+\bigl(1+(2r)^s\bigr)e^{r^s/(es)}+1.
\]
取对数后，有限多项式与项数至多贡献 \(O(\log r)\)，故 \(\log M(r,f)=O(r^s)\)，从而 \(\rho(f)\le s\)。让 \(s\downarrow a\) 得 \(\rho(f)\le a\)。当 \(a=0\) 时，此论证对每个 \(s>0\) 成立，给 \(\rho(f)=0\)；当 \(a=\infty\) 时，已证的 \(a\le\rho(f)\) 迫使 \(\rho(f)=\infty\)。阶公式的所有情形至此闭合。

下面固定 \(0<\rho=\rho(f)<\infty\)，并记
\[
 L=\limsup_{n\to\infty}n^{1/\rho}|c_n|^{1/n}.
\]
若 \(\tau(f)<\infty\)，任取 \(t>\tau(f)\)，有 \(M(r,f)\le C_t e^{t r^\rho}\)。最小化 \(t r^\rho-n\log r\)，应取 \(r^\rho=n/(\rho t)\)，于是
\[
 |c_n|\le C_t\left(\frac{e\rho t}{n}\right)^{n/\rho}.
\]
取 \(n\) 次根、乘 \(n^{1/\rho}\)，再取上极限，\(C_t^{1/n}\to1\)，所以 \(L\le(e\rho t)^{1/\rho}\)。令 \(t\downarrow\tau(f)\)，得到
\[
 L^\rho\le e\rho\tau(f).
\]
\(\tau(f)=0\) 时仍可从任意正 \(t\) 向下取极限；\(\tau(f)=\infty\) 时这个方向自动成立。

为证反向，若 \(L<\infty\)，任取 \(l>L\)。除有限项外，
\[
 |c_n|\le\bigl(l n^{-1/\rho}\bigr)^n.
\]
连续指标函数 \(x\log(lr)-(x/\rho)\log x\) 在 \(x=(lr)^\rho/e\) 取最大值 \((lr)^\rho/(e\rho)\)。将级数分在 \(n\le(2lr)^\rho\) 与 \(n>(2lr)^\rho\) 两部分：前一部分至多有 \(1+(2lr)^\rho\) 项，后一部分仍由几何级数 \(\sum2^{-n}\) 控制。连同有限头部，得到
\[
 \log M(r,f)\le\frac{l^\rho}{e\rho}r^\rho+O(\log r).
\]
除以 \(r^\rho\) 并取上极限，再令 \(l\downarrow L\)，便有 \(e\rho\tau(f)\le L^\rho\)。这也覆盖 \(L=0\)，因为可对每个 \(l>0\) 使用该估计。若 \(L=\infty\) 而 \(\tau(f)<\infty\)，与前一方向矛盾；若 \(\tau(f)=\infty\) 而 \(L<\infty\)，与后一方向矛盾。故两个端点同样相等。最后，正幂函数保持上极限，得到两种型公式的等价性。
\end{proof}
阶公式使用的是系数衰减的主尺度 \(n\log n\)，型则还保留其下一层指数常数。稀疏级数也可应用：零系数在阶公式中略去，在型公式中贡献零；但非零系数的下标必须始终使用它在原 Taylor 级数中的次数。

\ProseParagraphBreak

这些公式给出任意正实数阶的具体函数。设 \(s>0\)，令
\[
 F_s(z)=\sum_{n=0}^{\infty}\frac{z^n}{(n!)^{1/s}}.
\]
对递增函数 \(\log t\) 作积分比较，有
\[
 \int_1^n\log t\,\mathrm{d}t
 \le\log n!
 \le\int_1^n\log t\,\mathrm{d}t+\log n,
\]
因而 \(\log n!=n\log n-n+O(\log n)\)。特别地，\((n!)^{1/n}\sim n/e\)，系数根趋于零，故 \(F_s\) 为超越整函数。代入阶公式，
\[
 \frac{n\log n}{\log\bigl((n!)^{1/s}\bigr)}\longrightarrow s.
\]
再用型公式的第二种写法，
\[
 e s\tau(F_s)
 =\lim_{n\to\infty}\frac{n}{(n!)^{1/n}}=e,
\]
于是 \(\rho(F_s)=s\)、\(\tau(F_s)=1/s\)。例如 \(\sum z^n/(n!)^2\) 的阶为 \(1/2\)、型为 \(2\)。增长阶可以取非整数值，并不受多项式次数的整数性限制。

% !TeX root = ../../Book3.tex
\section{零点序列}
\label{b3:sec:1903}

Jensen 公式限制零点的数量，但还可以把整个零点序列压成一个收敛指数。下一章构造典范乘积时，这个指数决定必须扣掉多少低阶项。

\subsection{收敛指数}
\label{b3:subsec:190301}

按重数列出趋于无穷的非零零点 \(a_n\)。定义\term[lingdian-shoulian-zhishu]{零点收敛指数}[exponent of convergence]为 \(\lambda=\inf\{s>0:\sum|a_n|^{-s}<\infty\}\)。有限零点序列取零。对于无穷序列，若 \(n(r)\) 为计数，则
\[
\lambda=\limsup_{r\to\infty}\frac{\log n(r)}{\log r}.
\]
证明可用分部求和：\(\sum_{|a_n|>1}|a_n|^{-s}=s\int_1^\infty n(t)t^{-s-1}\,\mathrm{d}t\) 加有限边界项。若计数为 \(O(t^{s-\varepsilon})\) 积分收敛；若序列幂和收敛，则 \(n(r)\le r^s\sum|a_n|^{-s}\)，给反向。

\subsection{零点密度}
\label{b3:subsec:190302}

\begin{proposition}\label{b3:prop:19-zero-density}
有限阶 \(\rho\) 的非恒零整函数具有 \(n(r)=O_\varepsilon(r^{\rho+\varepsilon})\)，零点收敛指数不超过 \(\rho\)。
\end{proposition}
\begin{proof}
先除去原点零点。Jensen 估计取外半径 \(2r\)，得到 \(n(r)\log2\le\log M(2r)-\log|f(0)|\)。阶的定义给计数界；取稍小的 \(\varepsilon\) 后，上一积分判据使每个 \(s>\rho\) 的零点幂和收敛。
\end{proof}
原点的有限重数只增加常数项，不改变收敛指数。

\subsection{增长限制}
\label{b3:subsec:190303}

反向并不唯一：\(e^{z^m}\) 无零点而阶 \(m\)，所以零点只能约束增长的一部分。\(\sin\pi z\) 的零点为整数，收敛指数一，和阶一相同。

若指定 \(a_n=n^\alpha\)，其收敛指数为 \(1/\alpha\)。任何具有这些零点的整函数必须有阶至少 \(1/\alpha\)。下一章将证明离散零点总能实现，但若还要求有限阶，就要检查这个收敛指数。


\ProseParagraphBreak

零点计数的幂界可用二进尺度再证明收敛指数判据。若 \(n(r)\le Cr^a\)，在 \(2^j<|a_n|\le2^{j+1}\) 内，\(s>a\) 时幂和至多 \(C2^{(j+1)a-js}\)，对 \(j\) 求和为收敛几何级数。若 \(\sum|a_n|^{-s}<\infty\)，则半径 \(r\) 内每项至少 \(r^{-s}\)，故 \(n(r)r^{-s}\) 有界。这个分组证明不要求计数函数可微，也自然保留重数；分部求和则给更精确的常数与临界积分。

\ProseParagraphBreak

在临界指数上，下确界不决定幂和是否收敛。零点模长为 \(n\) 时收敛指数一，临界一次幂和发散；模长为 \(n\log^2(n+1)\) 时指数仍一，却临界一次幂和收敛。若模长为 \(n\log(n+1)\)，临界再次发散。对数修正不会改变 \(\log n(r)/\log r\) 的上极限，却决定下一章的最小初等因子次数。阶只决定粗尺度，属还要检查临界幂和，二者不可互换。

\ProseParagraphBreak

零点重数也能提高收敛指数。若每个正整数 \(n\) 的重数为 \(n\)，幂和变成 \(\sum n^{1-s}\)，收敛指数为二；计数约为 \(r^2/2\)，因此任何实现函数阶至少二。若仅按不同零点个数，则会得到错误下界一。对整函数增长、典范乘积和 Jensen 的贡献，统一按重数计数是必要条件，并非一个可选的记账方式。

\ProseParagraphBreak

增长下界只来自零点径向分布，角位置可以有进一步作用。例如整数零点在两条相反射线上排列，正弦给出阶一实现；把正整数零点全部放在同一射线，最小属仍是一，但沿某些方向的精确增长会不同。下一章的乘积修正会显式出现一次项，说明方向排列影响低次指数自由度。当前的 \(\lambda\le\rho\) 是普适必要条件，还没有声称它给所有方向的精确渐近。

\ProseParagraphBreak

若不仅知道阶，还知道正常型，Jensen 估计可给更具体的密度常数。假设 \(0<\rho<\infty\)、\(\tau<\infty\)，先去掉原点零点。对任意固定 \(u>1\)，有
\[
 n(r)\log u\le\log M(ur)-\log|f(0)|.
\]
除以 \(r^\rho\) 后取上极限，右侧不超过 \(\tau u^\rho\)。函数 \(u^\rho/\log u\) 在 \(\log u=1/\rho\) 处取最小值 \(e\rho\)，因此
\[
 \limsup_{r\to\infty}\frac{n(r)}{r^\rho}\le e\rho\tau.
\]
这仍是一般上界，不声称常数对每类函数都最优。它至少说明最小型 \(\tau=0\) 时 \(n(r)=o(r^\rho)\)，比只给 \(O(r^{\rho+\varepsilon})\) 更细。有限个原点零点除以 \(r^\rho\) 后消失，不改变结果。
\ProseParagraphBreak
对正弦函数，\(n(r)\) 与 \(r\) 同阶；对指数函数，计数恒为零，二者却有相同的阶与型。这再次说明增长和零点不是互相决定的两组资料。若只给定零点序列，下一章构造的典范乘积会固定其中一部分增长，而允许的无零点因子仍可改变最大模。使用密度估计时，正确方向是“增长上界限制零点”，不能不加额外假设把它倒置为“零点上界限制函数”。

% !TeX root = ../../Book3.tex
\section{Phragmén–Lindelöf 原理}
\label{b3:sec:1904}

最大值原理在无界扇形中不能直接应用。Phragmén--Lindelöf 方法先乘一个衰减指数，在大圆弧压住增长，再把有限区域的最大值原理送到无穷远。

\subsection{扇形域}
\label{b3:subsec:190401}

扇形 \(S=\{r e^{\mathrm{i}\theta}:r>0,\ |\theta|<\beta/2\}\)，\(0<\beta<2\pi\)，具有可选定的幂分支 \(z^p\)。若 \(p\beta/2<\pi/2\)，则 \(\operatorname{Re}z^p\ge c|z|^p\)，其中 \(c=\cos(p\beta/2)>0\)。这给整个闭扇形同方向的衰减，不能只检查中心射线。

\subsection{增长条件}
\label{b3:subsec:190402}

\term[phragmen-lindelof]{Phragmén--Lindelöf 原理}[Phragmén--Lindelöf principle]的一个常用版本如下。
\begin{theorem}[Phragmén--Lindelöf：扇形]\label{b3:thm:19-pl}
设 \(f\) 在上述扇形全纯并连续到有限边界与顶点。若两边界射线上 \(|f|\le1\)，且 \(|f(z)|\le C e^{A|z|^q}\)，其中 \(q<\pi/\beta\)，则 \(|f|\le1\) 于整个扇形。
\end{theorem}
\begin{proof}
选 \(q<p<\pi/\beta\)，令 \(f_\varepsilon=f e^{-\varepsilon z^p}\)。两射线上模不超过一；大圆弧上模不超过 \(C e^{Ar^q-\varepsilon c r^p}\)，充分大时亦不超过一。顶点值由连续性不超过一。对截断扇形用最大模原理得 \(|f_\varepsilon|\le1\)，再令大半径趋无穷，最后 \(\varepsilon\downarrow0\)。
\end{proof}
临界指数必须另设条件：\(e^{z^{\pi/\beta}}\) 在两边界模一，在中心射线却发散，给出严格不等式的必要性。

\subsection{整函数应用}
\label{b3:subsec:190403}

整函数若阶 \(\rho<1\)，并在实轴有界，则为常数。选择 \(\rho<q<1\)，阶的定义给 \(Ce^{Ar^q}\) 增长；分别对上、下半平面（角宽 \(\pi\)）应用扇形原理，得到全平面有界，再用 Liouville。

阶一的 \(e^{\mathrm{i}z}\) 在实轴模一而非常值，显示阈值不能只把严格不等号改成不大于。若函数仅在一个半平面全纯，结论只保证内部有界，不会自动成为常数。

\subsection{唯一性定理}
\label{b3:subsec:190404}

同一方法给唯一性。若 \(f,g\) 在扇形满足相同的次临界增长，且在两边界相等，则差 \(h\) 的边界零。对每个 \(\varepsilon>0\)，缩放 \(h/\varepsilon\) 满足边界模不超过一，扇形原理给 \(|h|\le\varepsilon\)，故 \(f=g\)。

这里边界相等不能直接当作恒等定理的内部聚点条件。本证明通过增长界得到全域比较；若边界为可作 Schwarz 反射的直线，还能用反射后的恒等定理得到更强的局部唯一性。临界例子中的零边值实部 \(\operatorname{Re}z^{\pi/\beta}\) 也说明无穷远会引入额外自由度。


\ProseParagraphBreak

扇形的旋转不改变临界阈值。若中轴辐角为 \(\theta_0\)，应在衰减因子中使用 \((e^{-\mathrm{i}\theta_0}z)^p\)，使实部在整个扇形有正下界；直接保留 \(z^p\) 可能在旋转后变成增长因子。幂的分支固定在该扇形中，连续延到两边界射线；顶点只有连续值，不必声称非整数幂在顶点全纯。截断扇形应用最大模原理时，顶点属于边界，内域不包含那个可能的分支点。

\ProseParagraphBreak

若 \(f\) 在顶点仅有界而未声明连续，证明可以先删去小圆。小圆上的模上界可能大于一，不能直接丢弃；若顶点为可去奇点，补上后再使用原定理即可。若允许顶点有极点，则边界射线在接近顶点仍受模一控制已经强烈限制它；但对于只控制远处边界的变式，小圆会留下额外常数。定理的有限边界与顶点假设用来控制截断域的全部边界，不能只关注大圆弧的增长。

\ProseParagraphBreak

一个多射线应用可以把阶阈值提高。若整函数阶 \(\rho\) 小于 \(\pi/\beta\)，并在若干从原点发出的射线上有共同模界 \(M\)，且相邻射线最大角宽为 \(\beta\)，则每个夹角扇形都可选增长指数 \(q\) 位于 \(\rho\) 与 \(\pi/\beta\) 之间。分别使用扇形定理，得到全平面 \(|f|\le M\)，故常数。例如四条坐标半轴把平面分成角宽 \(\pi/2\) 的扇形，于是阶小于二且在实轴、虚轴都有界 的整函数为常数。临界阶二的 \(e^{z^2}\) 在虚轴有界，在实轴却无界；\(e^{\mathrm{i}z^2}\) 在两轴均模一而非常值，给出正确的临界反例。

\ProseParagraphBreak

这里有限多个扇形的常数可以统一，因为只取有限次最大值；若射线无限多而各自界不同，不应未经估计就宣称有一个全平面共同界。把每条射线上的界与一个统一界区别开，也是使用 Liouville 结论前必须核查的量词。

\ProseParagraphBreak

边界零的唯一性应用需要同时假设两函数的增长。差函数的界由两者相加获得，常数可能加倍但指数仍是同一个次临界 \(q\)。然后对任意 \(\delta>0\) 使用 \(h/\delta\)，增长常数会随 \(\delta\) 改变；定理允许每次使用不同有限常数，因此最终 \(|h|\le\delta\) 合法。不能要求所有 \(\delta\) 共用一个增长常数，那会把证明中的无害缩放误认为失败。

\ProseParagraphBreak

无界区域比较还有一个常用的条带版本，它让边界模界不同的情形呈现出插值形式。
\begin{theorem}[有界函数的三线不等式]\label{b3:thm:19-three-lines}
设 \(f\) 在 \(0<\operatorname{Re}z<1\) 全纯，在闭条带连续有界，并满足 \(|f(\mathrm{i}t)|\le M_0\)、\(|f(1+\mathrm{i}t)|\le M_1\)，其中 \(M_0,M_1>0\)。则
\[
 |f(x+\mathrm{i}y)|\le M_0^{1-x}M_1^x\qquad(0\le x\le1).
\]
\end{theorem}
\begin{proof}
令 \(g(z)=f(z)\exp(-(1-z)\log M_0-z\log M_1)\)。实对数无分支问题，\(g\) 在闭条带有界，两竖边模不超过一。对 \(\varepsilon>0\)，令 \(g_\varepsilon(z)=g(z)e^{\varepsilon z^2-\varepsilon}\)。在两竖边有 \(\operatorname{Re}(\varepsilon z^2-\varepsilon)=\varepsilon(x^2-y^2-1)\le0\)。在横边 \(y=\pm T\)，其模不超过 \(\|g\|_\infty e^{-\varepsilon T^2}\)，充分大 \(T\) 时也不超过一。矩形最大模原理给 \(|g_\varepsilon|\le1\)，令 \(T\to\infty\)、\(\varepsilon\downarrow0\) 得 \(|g|\le1\)，还原即得结论。
\end{proof}
有界性在这条证明中控制横边，指数修正把横边送到无穷远后消去。若希望允许更快的竖向增长，必须选择新的修正并重新估计，不能直接引用当前版本。三线不等式的权重由观察点在两边界之间的相对位置决定，正如三圆定理的权重由对数半径决定；两个公式背后都是求出合适的调和比较函数。
\ProseParagraphBreak
如果某一边界模界为零，可先把该 \(M_j\) 替换为任意小正数，应用已经证明的正界版本再取极限；于是整个条带函数为零。这一处理不需要对零取对数，也不会留下未定义的指数因子。一般的边界参数极限应在不等式已经建立之后执行，而不是在构造辅助全纯函数时直接代入奇异参数。


\begin{chaptersummary}
三圆定理使 \(\log M(e^t)\) 成为凸函数，等号通过次调和强最大值原理限制全纯函数的形式。最大项的对数则是由 Taylor 系数决定的直线上包络。Cauchy 估计与模级数控制把这两个径向量联系起来，既给出多项式判别，也说明何时主导次数必须不断向高次移动。

\ProseParagraphBreak
阶与型记录不同层次的增长。阶为零仍可能是超越整函数，有限正阶也未必有有限正型。系数判据必须同时证明必要性和充分性，不能只由 Cauchy 估计得到单向限制。精确型的比较需要让半径比趋于一；固定倍增半径通常只足以确定阶。加法与乘法的上界允许抵消，具体等式还须检查函数之间的关系。

\ProseParagraphBreak
零点收敛指数由计入重数的幂和及计数增长共同刻画，下确界本身并不决定临界幂和是否收敛。Jensen 公式把零点的对数权重放在最大模之下，所以增长上界限制零点密度，既有密度又给增长提供必要下界。无零点的指数因子没有改变这些零点资料，却能改变型乃至阶。

\ProseParagraphBreak
无界区域的比较需要控制送往无穷远的那部分边界。扇形的临界指数由角宽决定，严格次临界情形保留衰减幂次的选择余量，临界最小型情形则用更精细的屏障。三线比较将两条边界的常数对数作调和插值，竖向增长条件保证截断边界的影响消失。把射线界或条带界用于全局结论时，应逐项检查这些增长条件及有限边界假设。
\end{chaptersummary}

\BookExercises

% 自拟：计算有限指数和的逐方向上极限；对并列主导指数用长区间平方平均排除持续相消，不调用一般指标理论。
\Needspace{6\baselineskip}
\begin{exercise}[有限指数和的方向增长]\label{b3:ex:19-draft-01}
设 \(\lambda_1,\ldots,\lambda_k\) 为互异复数，\(c_1,\ldots,c_k\ne0\)，且至少一个 \(\lambda_j\ne0\)。令
\[
 F(z)=\sum_{j=1}^kc_je^{\lambda_jz}.
\]
证明对每个固定 \(\theta\in\mathbb R\)，有
\[
 \limsup_{r\to\infty}\frac{\log|F(re^{\mathrm i\theta})|}{r}
 =\max_j\operatorname{Re}(\lambda_je^{\mathrm i\theta}).
\]
允许射线上出现零点，届时对数取 \(-\infty\)。特别要处理两个或更多指数具有相同最大实部的情形。由此求 \(F\) 的阶与型。
\end{exercise}
\begin{solution}
固定 \(\theta\)，记
\[
 H=\max_j\operatorname{Re}(\lambda_je^{\mathrm i\theta}),\quad
 J=\{j:\operatorname{Re}(\lambda_je^{\mathrm i\theta})=H\}.
\]
三角不等式给 \(|F(re^{\mathrm i\theta})|\le e^{Hr}\sum_j|c_j|\)，所以题中上极限不超过 \(H\)。为得反向估计，写
\[
 e^{-Hr}F(re^{\mathrm i\theta})=p(r)+E(r),\quad
 p(r)=\sum_{j\in J}c_je^{\mathrm i\omega_jr},\quad
 \omega_j=\operatorname{Im}(\lambda_je^{\mathrm i\theta}).
\]
不在 \(J\) 中的项具有严格较小实部，故 \(E(r)\to0\)；若没有这些项，取 \(E=0\)。\(J\) 中的实频率 \(\omega_j\) 互异，因为相同的实部和虚部会使两个 \(\lambda_j\) 相等。

将平方展开，在长区间 \([T,2T]\) 上平均，得到
\[
 \frac1T\int_T^{2T}|p(r)|^2\,\mathrm dr
 =\sum_{j\in J}|c_j|^2
 +\sum_{\substack{j,\ell\in J\\j\ne\ell}}
 c_j\overline{c_\ell}\frac1T\int_T^{2T}e^{\mathrm i(\omega_j-\omega_\ell)r}\,\mathrm dr.
\]
每个非对角积分除以 \(T\) 后的模不超过 \(2/(T|\omega_j-\omega_\ell|)\)，故平均趋于
\(S=\sum_{j\in J}|c_j|^2>0\)。因此对充分大的 \(T\)，区间内有一点 \(r_T\)，使 \(|p(r_T)|\ge\sqrt{S/2}\)。又因 \(E\) 在 \([T,\infty)\) 一致趋小，最终有
\[
 |F(r_Te^{\mathrm i\theta})|\ge\frac{\sqrt S}{4}e^{Hr_T}.
\]
沿 \(r_T\to\infty\) 取对数并除以 \(r_T\)，得到下界 \(H\)。这只需要找到不断远离原点的一列非相消位置，并未声称射线上存在普通极限。

令 \(L=\max_j|\lambda_j|>0\)。全局上界为 \(M(r,F)\le(\sum|c_j|)e^{Lr}\)。选一个满足 \(|\lambda_j|=L\) 的指数，再取 \(\theta=-\arg\lambda_j\)，则该方向的 \(H=L\)。前面的下界遂给
\[
 \limsup_{r\to\infty}\frac{\log M(r,F)}r=L.
\]
所以 \(F\) 的阶为一，型为 \(L\)。
\end{solution}

% 自拟：反向设计任意正阶和扩展型，允许任意稀疏非零系数支集并保留有限 Taylor 数据。
\Needspace{6\baselineskip}
\begin{exercise}[稀疏系数的阶型构造]\label{b3:ex:19-draft-02}
给定 \(\rho>0\)、\(\tau\in[0,\infty]\)、整数 \(N\ge0\)、次数不超过 \(N\) 的多项式 \(P\)，以及任意严格递增整数列
\[
 N<n_1<n_2<\cdots.
\]
构造超越整函数 \(f\)，使其阶为 \(\rho\)、型为 \(\tau\)，并且 \(f-P\) 的非零 Taylor 系数指标恰为 \(n_1,n_2,\ldots\)。构造应同时覆盖最小型、有限正型和最大型，并说明为什么对指标列不需要额外的密度条件。
\end{exercise}
\begin{solution}
按所需型取正数
\[
 s_k=
 \begin{cases}
 (\log(n_k+e))^{-1},&\tau=0,\\
 \tau,&0<\tau<\infty,\\
 \log(n_k+e),&\tau=\infty,
 \end{cases}
 \quad
 c_{n_k}=\left(\frac{e\rho s_k}{n_k}\right)^{n_k/\rho},
\]
并定义
\[
 f(z)=P(z)+\sum_{k=1}^\infty c_{n_k}z^{n_k}.
\]
其余高次系数都取零。三种情况下均有 \(s_k/n_k\to0\)，所以
\[
 |c_{n_k}|^{1/n_k}=\left(\frac{e\rho s_k}{n_k}\right)^{1/\rho}\longrightarrow0.
\]
全体系数的根也趋于零，级数在每个闭圆盘绝对一致收敛，定义整函数。所列高次系数全部为正且无穷多个非零，所以函数不是多项式。

又有 \(\log s_k=o(\log n_k)\)，从而
\[
 \frac{n_k\log n_k}{\log(1/|c_{n_k}|)}
 =\frac{\rho\log n_k}{\log n_k-\log(e\rho s_k)}
 \longrightarrow\rho.
\]
系数判据给出所需阶；有限多项式头部与零系数不改变这个上极限。型的判据则给
\[
 n_k|c_{n_k}|^{\rho/n_k}=e\rho s_k,\quad
 \tau(f)=\frac1{e\rho}\limsup_k e\rho s_k=\tau.
\]
当目标型为零或无穷时，最后的等式分别按趋零或趋于无穷解释。

整个计算只用到了 \(n_k\to\infty\)，没有用相邻指标的间距。因此无论这些指标多么稀疏，均可在每个非零系数上安排适当大小，保持指定的阶与型；任意给定的有限 Taylor 数据也保持不变。
\end{solution}

% 自拟：以双指数衰减系数求主导项的确切转折，再用最大项比较导出不同于正文高斯系数例子的零阶增长。
\Needspace{6\baselineskip}
\begin{exercise}[最大项的转折与零阶增长]\label{b3:ex:19-draft-03}
令
\[
 F(z)=\sum_{n=0}^\infty e^{-2^n}z^n,\quad
 \mu(r)=\max_{n\ge0}e^{-2^n}r^n.
\]
证明 \(F\) 为整函数，并精确求出最大项指标发生变化的半径。当 \(t=\log r>1\) 时，确定哪些整数 \(n\) 能使 \(nt-2^n=\log\mu(r)\)，包括并列最大情形。最后证明
\[
 \log M(r,F)=\frac{\log r\,\log\log r}{\log2}+O(\log r)
 \quad(r\to\infty),
\]
并据此说明 \(F\) 是零阶超越整函数，而且其最大模超过每个固定幂 \(r^d\)。
\end{exercise}
\begin{solution}
系数根为 \(e^{-2^n/n}\to0\)，所以收敛半径无穷。由于系数非负，三角不等式在正实轴取等号，故 \(M(r,F)=F(r)\)。

令 \(\phi_n(t)=nt-2^n\)。相邻项之差为
\[
 \phi_{n+1}(t)-\phi_n(t)=t-2^n,
\]
它随 \(n\) 严格递减。因此各项先增加后减少，\(n\ge1\) 为最大指标的充要条件是
\[
 2^{n-1}\le t\le2^n.
\]
若两个不等式均严格，最大指标唯一；若 \(t=2^n\)，恰有 \(n,n+1\) 两个并列指标。\(t<1\) 时指标为零，\(t=1\) 时指标零与一并列。所有转折半径因而正是
\(r=\exp(2^n)\)、\(n\ge0\)。

当 \(t>1\) 时，选择任一最大指标 \(n\)，有
\[
 n=\frac{\log t}{\log2}+O(1),\quad t\le2^n\le2t.
\]
所以
\[
 \log\mu(e^t)=nt-2^n=\frac{t\log t}{\log2}+O(t).
\]
最大项比较给
\(\mu(r)\le M(r,F)\le2\mu(2r)\)。将后一项的对数半径写成 \(t+\log2\)，上面的渐近式仍具有主项 \(t\log t/\log2\)，误差为 \(O(t)\)。上下夹逼便得所求最大模渐近。

于是
\[
 \frac{\log\log M(r,F)}{\log r}\longrightarrow0,
 \quad
 \frac{\log M(r,F)}{\log r}
 =\frac{\log\log r}{\log2}+O(1)\longrightarrow\infty.
\]
第一式给零阶，第二式说明对每个固定 \(d\)，最终 \(M(r,F)>r^d\)，并排除多项式。主导项虽然移动得很慢，仍会依次经过任意高的次数。
\end{solution}

% 自拟：将离散有限差分还原为插值多项式，再以整数零点和最小指数增长排除非零余项。
\Needspace{6\baselineskip}
\begin{exercise}[整数数据与多项式唯一性]\label{b3:ex:19-draft-04}
设整函数 \(f\) 满足：对每个 \(\varepsilon>0\)，存在 \(C_\varepsilon<\infty\)，使
\[
 |f(z)|\le C_\varepsilon e^{\varepsilon|z|}\quad(z\in\mathbb C).
\]
又设某个整数 \(m\ge1\) 使
\[
 \sum_{j=0}^m(-1)^{m-j}\binom mj f(n+j)=0
 \quad(n=0,1,2,\ldots).
\]
证明 \(f\) 是次数小于 \(m\) 的多项式，并用 \(f(0),\ldots,f(m-1)\) 写出它。构造一个阶一、有限正型的非多项式，满足相同的离散条件，从而说明增长假设的作用。
\end{exercise}
\begin{solution}
对整数数据记 \(\Delta f(n)=f(n+1)-f(n)\)，题设即 \(\Delta^mf(n)=0\)。定义多项式
\[
 P(z)=\sum_{k=0}^{m-1}(\Delta^kf)(0)\binom zk,
 \quad
 \binom z0=1,\quad
 \binom zk=\frac{z(z-1)\cdots(z-k+1)}{k!}\quad(k\ge1).
\]
其中每个 \((\Delta^kf)(0)\) 只使用 \(f(0),\ldots,f(k)\)。移位算子满足 \(E=I+\Delta\)，故对 \(0\le j<m\)，有限二项式展开给
\[
 f(j)=(E^jf)(0)=\sum_{k=0}^j\binom jk(\Delta^kf)(0)=P(j).
\]
多项式 \(P\) 的次数小于 \(m\)，因此 \(\Delta^mP=0\)。两列数据 \(f(n)\)、\(P(n)\) 满足同一个 \(m\) 阶递推式，最后一项的系数为一，且前 \(m\) 项相同。逐项递推便得 \(f(n)=P(n)\) 对所有非负整数成立。

令 \(h=f-P\)。多项式也满足任意正指数的增长界，因此若 \(h\not\equiv0\)，则
\(\log M(r,h)=o(r)\)。设原点重数为 \(\ell\)，写 \(h(z)=z^\ell g(z)\)、\(g(0)\ne0\)。所有正整数仍为 \(g\) 的零点。Jensen 的两半径估计于是给
\[
 \lfloor R\rfloor\log2
 \le\log M(2R,g)-\log|g(0)|
 =\log M(2R,h)-\ell\log(2R)-\log|g(0)|.
\]
右端是 \(o(R)\)，左端与 \(R\) 成正比，令 \(R\to\infty\) 即矛盾。所以 \(h=0\)，所求函数就是上面显式给出的 \(P\)。

函数 \(\sin\pi z\) 在所有整数上为零，故对任意 \(m\ge1\) 都满足题中的离散条件，却不是多项式。它的阶一、型为 \(\pi\)，不满足对每个 \(\varepsilon>0\) 的原增长界。可见把最小指数增长放宽为任意有限正常型，会失去该唯一性结论。
\end{solution}

% 对原自拟题实质改编：对同一组显式零点计算计数与增长，并在给定指数修正族内求精确最小型值，不使用分解定理。
\Needspace{6\baselineskip}
\begin{exercise}[同一组零点与可变的型]\label{b3:ex:19-draft-05}
对整数 \(m\ge1\) 和 \(a\in\mathbb C\)，令
\[
 F_a(z)=e^{az^m}(e^{z^m}-1).
\]
求其全部零点与重数、闭圆盘零点计数及收敛指数，并判定临界幂和是否收敛。证明 \(F_a\) 的阶为 \(m\)，型为
\[
 \tau(F_a)=\max\{|a|,|a+1|\}.
\]
确定这族函数的最小型值及达到该值的参数，并说明其圆周对数平均为什么不依赖 \(a\)。这里型的最小值只在给定的 \(F_a\) 族内比较。
\end{exercise}
\begin{solution}
指数前因子没有零点。方程 \(e^{z^m}=1\) 等价于 \(z^m=2\pi\mathrm i k\)、\(k\in\mathbb Z\)。\(k=0\) 给原点的 \(m\) 重零点，因为 \(e^{z^m}-1=z^m+O(z^{2m})\)。对 \(k\ne0\)，方程有 \(m\) 个互异根，且导数 \(mz^{m-1}e^{z^m}\) 在各根不为零，所以它们都是单零点。

于是对 \(r>0\)，闭圆盘计数为
\[
 n_{F_a}(r)=m+2m\left\lfloor\frac{r^m}{2\pi}\right\rfloor.
\]
全部非零零点的 \(s\) 次幂和为
\[
 2m\sum_{k=1}^\infty(2\pi k)^{-s/m}.
\]
它收敛当且仅当 \(s>m\)，故收敛指数是 \(m\)，临界 \(s=m\) 时发散。

记 \(H_a(w)=e^{(a+1)w}-e^{aw}\)，并令 \(L=\max\{|a|,|a+1|\}>0\)。上界为 \(|H_a(w)|\le2e^{L|w|}\)。在 \(a,a+1\) 中选一个模为 \(L\) 的数 \(\lambda\)，另一个记为 \(\nu\)。取 \(w=Re^{-\mathrm i\arg\lambda}\)，则
\[
 \operatorname{Re}(\lambda w)=LR,\quad
 \operatorname{Re}(\nu e^{-\mathrm i\arg\lambda})<L.
\]
后一不等式严格，因为 \(|\nu|\le L\)，而实部达到 \(L\) 会迫使 \(\nu=\lambda\)，与两指数相差一矛盾。因此存在 \(\delta>0\)，使
\[
 |H_a(w)|\ge e^{LR}(1-e^{-\delta R}).
\]
由 \(z\mapsto z^m\) 把半径 \(r\) 的圆周满射到半径 \(r^m\) 的圆周，得到
\[
 \log M(r,F_a)=Lr^m+O(1)
\]
于充分大的 \(r\)。所以阶为 \(m\)，型为 \(L\)。

三角不等式给
\[
 1=|(a+1)-a|\le|a+1|+|a|\le2L,
\]
故型的最小可能值为 \(1/2\)，在 \(a=-1/2\) 时达到，此时 \(F_a(z)=2\sinh(z^m/2)\)。若 \(L=1/2\)，上面两个不等式均取等号，\(a\) 与 \(a+1\) 必须位于同一直线的相反方向且模都为 \(1/2\)，故仍有 \(a=-1/2\)。此外取实数 \(a=L-1\)，可实现任意 \(L\ge1/2\)，这族函数的型恰遍历 \([1/2,\infty)\)。

最后，在每个圆周上有几乎处处的等式
\[
 \log|F_a(re^{\mathrm it})|
 =\operatorname{Re}(ar^me^{\mathrm imt})+\log|e^{r^me^{\mathrm imt}}-1|.
\]
第二项的零点奇性可积，第一项的平均为零，所以圆周对数平均完全不依赖 \(a\)。零点和这一平均保持不变，最大模的型却可以在上述范围内变化。
\end{solution}

% 对原自拟题实质改编：由导数界建立外域 Laurent 生成函数与积分重建，并用轮廓变形识别精确指数增长率。
\Needspace{6\baselineskip}
\begin{exercise}[导数生成函数与指数增长]\label{b3:ex:19-draft-06}
设 \(f\) 是非恒零整函数，且
\[
 T=\limsup_{r\to\infty}\frac{\log M(r,f)}r<\infty.
\]
证明
\[
 T=\limsup_{n\to\infty}|f^{(n)}(0)|^{1/n}.
\]
定义 Laurent 级数
\[
 \mathcal B_f(\zeta)=\sum_{n=0}^\infty\frac{f^{(n)}(0)}{\zeta^{n+1}}.
\]
证明它在 \(|\zeta|>T\) 全纯，且对任意 \(R>T\)、\(z\in\mathbb C\)，有
\[
 f(z)=\frac1{2\pi\mathrm i}\int_{|\zeta|=R}e^{z\zeta}\mathcal B_f(\zeta)\,\mathrm d\zeta.
\]
若 \(T>0\)，证明 \(\mathcal B_f\) 不可能全纯延拓到任何外域 \(|\zeta|>s\)、\(0\le s<T\)。
\end{exercise}
\begin{solution}
非恒零整函数的最大模在大圆上有一个正下界，因此 \(T\ge0\)。给定 \(t>T\)，有 \(M(r,f)\le C_t e^{tr}\)。对 \(n\ge1\)，Cauchy 估计并取 \(r=n/t\)，得到
\[
 |f^{(n)}(0)|\le n!C_t e^{tr}r^{-n}
 =C_t n!\left(\frac{et}{n}\right)^n.
\]
由 \(\log n!=n\log n-n+O(\log n)\)，有 \((n!)^{1/n}/n\to e^{-1}\)。所以若记右侧所求上极限为 \(L\)，便得 \(L\le t\)。令 \(t\downarrow T\)，得到 \(L\le T\)，也包括 \(T=0\) 的情形。

反向，给定 \(b>L\)，所有导数可用某个常数 \(A_b\) 统一估计为 \(|f^{(n)}(0)|\le A_b b^n\)，有限个初始指标由扩大常数吸收。Taylor 级数于是给
\[
 |f(z)|\le A_b\sum_{n=0}^\infty\frac{(b|z|)^n}{n!}=A_b e^{b|z|}.
\]
所以 \(T\le b\)，再令 \(b\downarrow L\) 即得 \(T=L\)。

这一系数根上极限说明 \(\mathcal B_f\) 在 \(|\zeta|>T\) 的每个紧集上绝对一致收敛，因而全纯。固定 \(R>T\)。Laurent 级数和 \(e^{z\zeta}\) 的幂级数在积分圆周上都绝对一致收敛，可以逐项相乘积分。只有指数为 \(\zeta^{-1}\) 的项留下，故
\[
 \frac1{2\pi\mathrm i}\int_{|\zeta|=R}e^{z\zeta}\mathcal B_f(\zeta)\,\mathrm d\zeta
 =\sum_{n=0}^\infty f^{(n)}(0)\frac{z^n}{n!}=f(z).
\]
这里 \(z\) 可以是任意固定复数，不需要限制在某个小圆盘。

若 \(T>0\) 且该函数全纯延拓到 \(|\zeta|>s\)、\(s<T\)，选择 \(s<R_1<T<R\)。被积函数在两圆之间的闭环域邻域全纯，Cauchy 定理允许将重建积分缩到 \(|\zeta|=R_1\)。于是对所有 \(z\)，有
\[
 |f(z)|\le R_1\max_{|\zeta|=R_1}|\mathcal B_f(\zeta)|\,e^{R_1|z|}.
\]
这给 \(T\le R_1<T\)，矛盾。因此精确指数增长率也限制了这个外域生成函数能够向内延拓的范围。整个论证容许 \(T=0\)；此时生成函数在任意非零 \(\zeta\) 处有定义，最后的严格缩小半径命题则没有需要讨论的非负 \(s\)。
\end{solution}

% 自拟：将多射线的多项式边界界转为全局幂界；分扇形选外点分母，保留临界最小型及非整数次数限制。
\Needspace{6\baselineskip}
\begin{exercise}[射线上的多项式控制]\label{b3:ex:19-draft-07}
取有限条互异射线，辐角依次为 \(\theta_1<\cdots<\theta_k<\theta_1+2\pi\)，其中 \(k\ge2\)，并令 \(\theta_{k+1}=\theta_1+2\pi\)、
\[
 \beta=\max_j(\theta_{j+1}-\theta_j),\quad \kappa=\frac\pi\beta.
\]
设整函数 \(f\) 在这些射线上满足
\[
 |f(re^{\mathrm i\theta_j})|\le C(1+r)^N
 \quad(r\ge0,\ 1\le j\le k),
\]
其中 \(N\ge0\) 为实数。若 \(\rho(f)<\kappa\)，或 \(\rho(f)=\kappa\) 且在该阶为最小型，证明 \(f\) 是次数不超过 \(\lfloor N\rfloor\) 的多项式。说明临界有限正型不能直接纳入此结论。
\end{exercise}
\begin{solution}
令 \(d=\lceil N\rceil\)。对相邻射线夹成的每个闭扇形 \(\overline S_j\)，取一个单位模的外点 \(a_j\notin\overline S_j\)。每个角宽都小于 \(2\pi\)，所以这样的外点存在。还有常数 \(c_j>0\)，使
\[
 |z-a_j|\ge c_j(1+|z|)\quad(z\in\overline S_j).
\]
事实上，在 \(|z|\ge2\) 时有 \(|z-a_j|\ge|z|-1\ge(1+|z|)/3\)；在剩下的紧集上，距离比为连续正函数，取其正最小值即可。

在 \(S_j\) 上定义
\[
 g_j(z)=\frac{f(z)}{(z-a_j)^d}.
\]
极点位于闭扇形之外，所以 \(g_j\) 在扇形全纯并连续到有限边界和顶点。两条边界射线上有
\[
 |g_j(z)|\le Cc_j^{-d}(1+|z|)^{N-d}\le Cc_j^{-d}.
\]
分母的下界还说明它保留了所需的增长上界。

若 \(\rho(f)<\kappa\)，取 \(\rho(f)<q<\kappa\)。每个扇形的临界指数 \(\pi/(\theta_{j+1}-\theta_j)\) 都不小于 \(\kappa\)，故可对 \(g_j\) 使用严格次临界 Phragmén--Lindelöf 原理。若 \(f\) 为阶 \(\kappa\) 的最小型，则在角宽等于 \(\beta\) 的扇形中使用临界最小型版本；在较窄扇形中，\(\kappa\) 已严格小于该扇形的临界指数，仍可使用严格版本。两种情况下均得 \(g_j\) 在其扇形有界。

因此存在 \(A_j<\infty\)，使 \(|f(z)|\le A_j(1+|z|)^d\) 于 \(\overline S_j\)。扇形只有有限个，取这些常数的最大值便得到全平面的同类幂界。多项式增长判别说明 \(f\) 是次数不超过 \(d\) 的多项式。若其实际次数 \(\ell>N\)，沿任一固定射线都有
\[
 |f(re^{\mathrm i\theta_j})|=|c_\ell|r^\ell(1+o(1)),
\]
与原来的 \((1+r)^N\) 界矛盾，所以实际次数至多 \(\lfloor N\rfloor\)。零多项式自然包括在结论中。这里每个商只用于自己的扇形，不要求把带外部极点的各个商合成一个整函数。

为看清临界限制，取整数 \(m\ge1\) 及射线 \(\theta_j=j\pi/m\)、\(j=0,\ldots,2m-1\)。它们的最大角宽为 \(\pi/m\)，临界指数为 \(m\)。整函数 \(e^{\mathrm iz^m}\) 在每条射线上模恒为一，阶 \(m\)、型一，却不是多项式。因此临界处仅有有限正型不足以替代最小型。
\end{solution}

% 自拟：为两侧高斯模界构造精确的三次调和修正，给出尖锐内部估计及机械同高插值的反例。
\Needspace{6\baselineskip}
\begin{exercise}[高斯边界权的三线比较]\label{b3:ex:19-draft-08}
设 \(a_0,a_1\ge0\)，\(f\) 在条带 \(0<\operatorname{Re}z<1\) 全纯，并连续到闭条带。假设存在 \(C>0\)、\(0<q<\pi\)，使
\[
 \log^+|f(x+\mathrm iy)|\le Ce^{q|y|}\quad(0\le x\le1),
\]
并且
\[
 |f(\mathrm iy)|\le e^{a_0y^2},\quad
 |f(1+\mathrm iy)|\le e^{a_1y^2}\quad(y\in\mathbb R).
\]
构造次数不超过三的全纯多项式 \(Q\)，使其在两条边界上的实部分别为 \(a_0y^2\)、\(a_1y^2\)，由此求出逐点最优的内部模界，并确定在一个内部点取等号的全部函数。最后用 \(a_0=a_1>0\) 说明，直接把同一高度的两个边界界作几何插值一般不正确。
\end{exercise}
\begin{solution}
令 \(\delta=a_1-a_0\)，取
\[
 Q(z)=-\frac\delta3z^3-a_0z^2+\frac{2a_0+a_1}{3}z.
\]
由 \(\operatorname{Re}z^2=x^2-y^2\) 和 \(\operatorname{Re}z^3=x^3-3xy^2\)，可算得
\[
 \begin{aligned}
 H(x,y)\coloneqq\operatorname{Re}Q(x+\mathrm iy)
 &=\bigl((1-x)a_0+xa_1\bigr)y^2\\
 &\quad+\frac{x(1-x)}3\bigl((2-x)a_0+(1+x)a_1\bigr).
 \end{aligned}
\]
这给 \(H(0,y)=a_0y^2\)、\(H(1,y)=a_1y^2\)。额外的只依赖 \(x\) 的项使整个表达式调和；只对两侧 \(y^2\) 系数作线性插值，通常没有这个性质。

令 \(g=f e^{-Q}\)。它在两竖边的模不超过一。由于 \(x\in[0,1]\)，上式还给 \(|\operatorname{Re}Q(x+\mathrm iy)|\le C_1(1+y^2)\)，所以
\[
 \log^+|g(x+\mathrm iy)|
 \le Ce^{q|y|}+C_1(1+y^2)
 \le C_2e^{q|y|}.
\]
最后一步使用 \(q>0\)，常数可以依赖于 \(q,a_0,a_1\)。因此适用本节允许竖向增长的三线比较，两条边界常数都为一，得到 \(|g|\le1\)。所求内部界为
\[
 \boxed{\ |f(x+\mathrm iy)|\le
 \exp\left(\bigl((1-x)a_0+xa_1\bigr)y^2
 +\frac{x(1-x)}3\bigl((2-x)a_0+(1+x)a_1\bigr)\right).\ }
\]

对任意实数 \(\alpha\)，函数 \(f(z)=e^{\mathrm i\alpha}e^{Q(z)}\) 在两侧和内部每一点都取等号。它的 \(\log^+\) 在闭条带上为 \(O(1+y^2)\)，满足题设增长条件，所以这个界逐点最优。反之，若任意可行 \(f\) 在一个内部点取等号，则 \(g\) 在内部取得最大模一，最大模原理给 \(g\equiv e^{\mathrm i\alpha}\)。因此上面的函数也就是全部等号情形。

若 \(a_0=a_1=a>0\)，则 \(Q(z)=az(1-z)\)，比较函数的实部简化为
\[
 H(x,y)=ay^2+ax(1-x).
\]
特别地，\(f(z)=e^{az(1-z)}\) 在两竖边均有模 \(e^{ay^2}\)，在 \(z=1/2\) 却有模 \(e^{a/4}>1\)。同一高度 \(y=0\) 的两个边界界都是一，其几何插值仍为一，遗漏了内部必要的 \(ax(1-x)\) 项。三线定理中的边界常数必须对整条边界成立，不能逐个高度代入不同的常数后仍沿用原公式。
\end{solution}

% !TeX root = ../../Book3.tex
\chapter{Weierstrass 乘积、Mittag–Leffler 与 Hadamard 分解}
\label{b3:ch:20}
\BookChapterNumber{b3:ch:20}
% 本章摘要（不计入编号）
\begin{chapterabstract}
无穷乘积通过扣除对数的低阶项实现指定零点，Mittag--Leffler 加法构造则实现指定主要部分。两种构造分别留下无零点乘法因子与整函数加法项。有限阶增长进一步将乘法因子压成指数多项式，得到 Hadamard 分解及正弦的典范乘积。
\end{chapterabstract}

% 本章目标（不计入编号）
\begin{learninggoals}
能检验正规乘积和零点阶数；能按离散序列选择初等因子；能构造整平面亚纯函数及整插值；能说明 Hadamard 证明中选圆周、乘积下界和实部控制的作用。
\end{learninggoals}

如果要在正整数处各放一个零点，直接相乘 \(1-z/n\) 会留下发散的一次对数项。把它改成 \((1-z/n)e^{z/n}\)，二次尾项就可以求和。一个看似小的指数修正使零点构造有了稳定的极限；而有限阶函数又反过来限制这些修正能增长多少。

\ProseParagraphBreak

本章使用\chapref{b3:ch:19}的零点计数与 Cauchy 估计。进一步阅读为 Conway 的《Functions of One Complex Variable I》\cite{Conway1978Functions}和 Levin 的《Lectures on Entire Functions》\cite{Levin1996Lectures}，可据零点构造与有限阶分解的主题继续阅读。

% !TeX root = ../../Book3.tex
\section{无穷乘积}
\label{b3:sec:2001}

有限个零点可以用多项式乘积实现。无穷个零点要求控制乘积尾部，且必须确认极限除了指定因子之外不会产生额外零点。

\subsection{收敛判据}
\label{b3:subsec:200101}

若 \(\sum_n|u_n|<\infty\)，则 \(\prod_n(1+u_n)\) 收敛；当所有因子非零时极限也非零。证明取尾部 \(|u_n|\le1/2\)，使用 \(\log(1+u_n)=\sum_{k\ge1}(-1)^{k+1}u_n^k/k\) 的固定小圆分支，得到 \(|\log(1+u_n)|\le2|u_n|\)。尾部乘积是收敛对数和的指数，故非零；有限头部另行相乘。

不控制对数和，仅有 \(u_n\to0\) 不足以保证收敛：\(\prod(1+1/n)\) 发散。也不能把“收敛”一律解释为非零；\(\prod(1-1/n)\) 从第二项起趋零。

\subsection{绝对收敛}
\label{b3:subsec:200102}

绝对收敛的乘积允许重排，其原因是小因子的对数绝对可和。无穷乘积的“绝对收敛”通常指 \(\sum|u_n|<\infty\)，不是对乘积模再作一次级数。

若有某个因子为零，整乘积为零，此时对数表述只用于去掉该因子后的尾部。后面在函数论中，零点附近总会把产生零点的有限因子与无零点的尾部拆开。

\subsection{正规收敛}
\label{b3:subsec:200103}

\begin{proposition}\label{b3:prop:20-normal-product}
若全纯函数 \(u_n\) 在区域 \(D\) 上满足对每个紧集 \(K\Subset D\)，\(\sum_n\sup_K|u_n|<\infty\)，则 \(\prod(1+u_n)\) 局部一致收敛到全纯函数；在无因子零点的地方，极限非零。
\end{proposition}
\begin{proof}
固定紧集附近稍大的紧集，尾部因子一致落在 \(|u|\le1/2\)，对数级数由上述估计正规收敛，其和全纯，指数给全纯非零尾部。乘上有限头部，得到局部一致极限。不同紧集的尾部表示来自同一有限乘积序列，故彼此一致。
\end{proof}
局部一致性保证可逐项微分对数：无零点的紧子集上 \(P'/P=\sum u_n'/(1+u_n)\)，导数收敛由全纯局部一致收敛定理给出。

\subsection{零点与乘积}
\label{b3:subsec:200104}

在固定点附近，若只有有限个因子为零，就把它们提出；余下尾部由正规对数和给非零函数。因此极限的零点阶数恰为有限零因子的阶数和，没有额外零点。

如果零因子在内部无限聚集，极限可能恒零，而指定离散零点的任务本身就与恒等定理冲突。本章的零点序列始终要求无有限聚点。


对乘积尾部还可给出误差界。若 \(S_N(K)=\sum_{n>N}\sup_K|u_n|\)，则由有限乘积展开
\[
\left|\prod_{n=N+1}^{M}(1+u_n)-1\right|\le
\prod_{n=N+1}^{M}(1+|u_n|)-1\le e^{S_N(K)}-1.
\]
令 \(M\to\infty\)，这个界仍成立。它直接控制尾部相对误差，而不仅是存在一个极限。有限头部在紧集的最大模有界，所以全乘积的绝对误差亦受同一尾和控制。若需要数值计算，先选 \(N\) 使尾和小，再单独计算头部；没有尾部界的有限截断不能证明函数恒等式。

\ProseParagraphBreak

以 \(u_n(z)=z/n^2\) 为例，\(|z|\le R\) 时 \(S_N\le R/N\)。只要没有碰到有限头部的零点，截断相对误差不超过 \(e^{R/N}-1\)。若零点是当前计算目标，必须提出相应因子后再谈相对误差，因为在零值处除以极限没有意义。相同估计还说明局部一致性不能改成全平面一致性：随着 \(R\) 增大，所需 \(N\) 也增大。

\ProseParagraphBreak

条件收敛的乘积可能依赖顺序。例如实数列 \(u_n=(-1)^n/n\) 从 \(n=2\) 起满足 \(\sum u_n\) 条件收敛、\(\sum u_n^2<\infty\)。由
\(\log(1+u_n)=u_n+O(u_n^2)\)，原顺序乘积收敛为正数，但正负项重排能改变一次和，从而改变乘积。这里所有因子甚至都正；失败来自绝对收敛不足，而不是复杂的复数分支。典范因子就是有目的地消除这些可能依赖顺序的低阶和。


\ProseParagraphBreak

正规收敛还允许检验导数，但须区分乘积与对数的适用位置。若紧集避开所有零点，取稍大的无零点邻域，并在那里对尾部统一选择对数，则 Cauchy 公式给对数级数各阶导数的局部一致收敛。若紧集包含一个零点，先分离该点的有限零因子，然后只对剩下的非零部分取对数。若直接写 \(P'/P=\sum u_n'/(1+u_n)\)，左右都是亚纯函数；在零点处谈全纯一致收敛并无意义。

\ProseParagraphBreak

例如 \(P(z)=\prod_{n\ge1}(1+z/n^2)\) 在原点附近的对数为
\[
\log P(z)=\sum_{k\ge1}\frac{(-1)^{k+1}z^k}{k}\sum_{n\ge1}n^{-2k},\qquad |z|<1.
\]
先取 \(|z|\le r<1\)，双重级数的绝对值由 \(\sum_n\sum_k r^k/n^{2k}\) 控制，后者不超过 \(r(1-r)^{-1}\sum_n n^{-2}\)。于是取原点一次导数得到 \(P'(0)=\sum n^{-2}\)，二次导数则为 \((\sum n^{-2})^2-\sum n^{-4}\)。对数系数记录单因子的贡献，函数系数还包含因子之间的交叉项；把这两类系数混在一起会漏掉平方项。

\ProseParagraphBreak

无穷乘积定理不能仅由各因子趋于一替代。常数因子 \(1+1/n\) 虽趋于一，其前 \(N\) 项的乘积为 \(N+1\)，仍发散。常数因子 \(1-1/n\) 从第二项开始乘积为 \(1/N\)，极限恒为零，不能用来断言“除有限个零因子以外尾部非零”。这两个计算说明，可和的对数尾项同时承担有限性与非消失性两项责任。判别时应检查它控制哪一项，不能只检查逐因子的极限。

% !TeX root = ../../Book3.tex
\section{Weierstrass 典范因子}
\label{b3:sec:2002}

直接相乘 \(1-z/a_n\) 往往失败，因为它的对数从一次项开始。乘上一个指数，把若干低阶项扣掉，尾部就可以变成更高次的可和项。

\subsection{初等因子}
\label{b3:subsec:200201}

定义\term[weierstrass-chudeng-yinzi]{Weierstrass 初等因子}[Weierstrass elementary factor]
\[
E_0(w)=1-w,\quad E_p(w)=(1-w)\exp\left(w+\frac{w^2}2+\cdots+\frac{w^p}p\right),\ p\ge1.
\]
它只有一个单零点 \(w=1\)，且 \(E_p(0)=1\)。指数部分不添加零点。

\subsection{典范因子}
\label{b3:subsec:200202}

当 \(|w|<1\)，固定对数分支给
\[
\log E_p(w)=-\sum_{k=p+1}^{\infty}\frac{w^k}k.
\]
所以 \(|w|\le1/2\) 时 \(|\log E_p(w)|\le2|w|^{p+1}\)，并可由 \(|e^v-1|\le e^{|v|}|v|\) 控制因子偏离一。若 \(\sum|a_n|^{-p-1}<\infty\)，则 \(\prod E_p(z/a_n)\) 在每个紧集正规收敛。这里 \(p\) 固定，而紧集半径只进入有限常数。

\subsection{属}
\label{b3:subsec:200203}

能使上述幂和收敛的最小非负整数 \(p\)，称该典范乘积的\term[dianfan-chengji-shu]{属}[genus]。当收敛指数 \(\lambda\) 不是整数时，\(p=\lfloor\lambda\rfloor\)；整数指数时，还要检查临界幂和是否收敛，不能只根据 \(\lambda\) 的值决定。

例如 \(a_n=n\) 的指数一，\(\sum1/n\) 发散而 \(\sum1/n^2\) 收敛，须用属一因子。若把正负整数成对相乘，一次指数正负抵消，可写成 \(\prod(1-z^2/n^2)\)；这依赖明确的配对，而不是许可随意重排一个条件收敛乘积。


初等因子的指数修正还控制对数导数。在避开 \(a_n\) 的地方，直接微分得到
\[
\frac{\mathrm{d}}{\mathrm{d}z}\log E_p(z/a_n)
=-\frac1{a_n-z}+\sum_{k=1}^{p}\frac{z^{k-1}}{a_n^k}
=-\frac{z^p}{a_n^p(a_n-z)}.
\]
当 \(|z|\le R\)、\(|a_n|>2R\) 时，最后一式的模不超过 \(2R^p/|a_n|^{p+1}\)。所以同一幂和不仅保证乘积存在，也保证尾部对数导数正规收敛。零点附近的留数由有限头部贡献，值为一；这将乘法的零点数据重新转成加法的主要部分数据。

\ProseParagraphBreak

对 \(a_n=n\)，属一乘积
\[
P(z)=\prod_{n\ge1}(1-z/n)e^{z/n}
\]
在原点满足 \(P(0)=1\)、\(P'(0)=0\)，且对数展开为
\[
\log P(z)=-\sum_{k\ge2}\frac{z^k}{k}\sum_{n\ge1}n^{-k},\quad |z|<1.
\]
换序由 \(\sum_{n,k}|z|^k/(kn^k)<\infty\) 保证。尤其二阶系数取决于 \(\sum n^{-2}\)，一次系数却因修正消失。这是初等因子“扣除低阶项”的具体可计算效果。

\ProseParagraphBreak

属在临界指数处的差别也能构造出来。点列 \(a_n=n\log^2(n+1)\) 的计数指数仍为一，但 \(\sum1/|a_n|\) 收敛，所以属零已够。点列 \(a_n=n\log(n+1)\) 的指数同样为一，而一次幂和发散，须用属一。二者的收敛指数相同，最小可用属不同。积分判别中出现的对数因子因此不能在临界问题中随意忽略。


\ProseParagraphBreak

若某个 \(a\) 应有重数 \(m\)，可以将它在点列中重复 \(m\) 次，也可以直接取因子 \(E_p(z/a)^m\)。对正规收敛的检验，后一写法须带上权重 \(m\)：充分条件成为 \(\sum m_n|a_n|^{-p-1}<\infty\)。遗漏重数会改变计数函数和收敛指数。例如在每个正整数 \(n\) 放重数 \(n\) 的零点时，带重数的幂和是 \(\sum n^{1-s}\)，收敛指数为二，属零或属一均不足；属二给 \(\sum n^{-2}\) 的尾和。

\ProseParagraphBreak

关于最小属还有一个容易混淆的区分。绝对收敛判据针对逐个初等因子的乘积；把具有对称性的因子成组后，低次项可能彼此抵消。例如正负整数的属一因子满足
\[
E_1(z/n)E_1(-z/n)=1-z^2/n^2.
\]
因此成对乘积有一个看似属零的二次因子表达，但它并没有使 \(\sum_{n\ne0}|n|^{-1}\) 收敛。解释“属”时必须说清楚使用单个零点的标准初等因子，还是已经经过配对的表达。后一表达式有利于计算，却不能反推前一幂和判据失效。

% !TeX root = ../../Book3.tex
\section{Weierstrass 乘积定理}
\label{b3:sec:2003}

离散性是指定零点的唯一存在性障碍。即使没有任何有限幂和收敛，逐项增加初等因子的阶数仍可使尾部收敛。

\subsection{指定零点}
\label{b3:subsec:200301}

\begin{theorem}[Weierstrass：整平面版本]\label{b3:thm:20-weierstrass}
任意无有限聚点的复平面点列与指定正整数重数，可以实现为一个非恒零整函数的全部零点及重数。
\end{theorem}
\begin{proof}
把非零点按重数重复记为 \(a_n\)，则 \(|a_n|\to\infty\)。取
\(P(z)=z^m\prod_{n\ge1}E_n(z/a_n)\)。对固定 \(|z|\le R\)，充分大 \(n\) 时 \(|z/a_n|\le1/2\)，故对数尾项不超过 \(2\cdot2^{-n-1}\)，一致可和。正规乘积命题给整函数；有限头部与非零尾部分离，保证零点及重数正好指定。
\end{proof}
所得函数可能增长非常快，定理没有给阶或型的最优控制。

\subsection{典范乘积}
\label{b3:subsec:200302}

若零点的收敛指数有限，可取固定 \(p\) 的典范乘积，而不用变动的 \(E_n\)。这种选择减少指数修正并保留增长信息。

给定同一零点列，所有实现函数都形如 \(P e^g\)，其中 \(g\) 整。因为商在每个零点处可去且无零点，商的对数导数有整平面原函数，指数还原该商。加法常数由一个点的商值决定。

\subsection{零点序列}
\label{b3:subsec:200303}

一个点列若在有限点聚集，任何包含它的整函数都由恒等定理恒零；这说明“无有限聚点”既必要又充分。有限列给多项式，无穷列给真正的超越整函数。

指定重数并不指定导数系数。若还要在离散点给出函数值或有限阶导数，需要插值；下一节用主要部分给出这种额外自由度的构造。


若指定的零点极为缓慢地趋于无穷，例如 \(a_n=\log(n+1)\)，任何固定 \(p\) 的幂和都发散。Weierstrass 定理仍适用：对紧集 \(|z|\le R\)，从足够大的 \(n\) 起有 \(|z/a_n|\le1/2\)，变动阶数 \(p_n=n\) 使误差指数衰减。这个例子区分了“离散性保证存在”与“有限收敛指数保证有限属”。对应零点计数约为 \(e^r\)，所以任何实现它们的非恒零整函数都必须是无穷阶；这由上一章的零点密度限制直接得到。

\ProseParagraphBreak

相反，零点 \(a_n=n^2\) 的属零乘积具有较慢增长。考虑直接由幂级数定义的整函数
\[
F(z)=\sum_{k\ge0}\frac{(-1)^k\pi^{2k}z^k}{(2k+1)!}
=\frac{\sin(\pi\sqrt z)}{\pi\sqrt z}
\]
的后一表达式按前面级数解释，不需要在全平面选择平方根。更换平方根的符号不改变商，所以局部表达也相容。它的零点恰为正整数平方，每个为单零点；在负实轴 \(z=-r\) 上由 \(\sinh(\pi\sqrt r)\) 得增长下界，复平面上由指数表示得同阶上界，故阶为二分之一。一个非整数阶整函数因而不必靠抽象存在定理展示，可以直接从偶幂级数得到。

\ProseParagraphBreak

对于有限零点集，分解 \(f=P e^g\) 中 \(P\) 是普通多项式。若 \(f\) 又有多项式增长，\(g\) 的自由度实际上只能给常数因子；若只要求有限阶，则 \(g\) 可为多项式。若没有任何增长条件，\(g=e^z\) 等无穷阶例子仍可出现。零点集合本身无法区分这些函数，必须把无穷远条件作为独立数据。


\ProseParagraphBreak

构造的局部性可以用圆盘上的有限头部说明。给定 \(R\)，离散性使 \(|a_n|\le2R\) 的零点只有有限多个。它们的因子可能很大，但全是已固定的整函数；收敛只由其余尾部负责。对 \(|a_n|>2R\)，取 \(p_n=n\) 后
\[
\sup_{|z|\le R}|\log E_n(z/a_n)|\le2(R/|a_n|)^{n+1}\le2^{-n}.
\]
这给出了无需根据每个紧集重新排序的选择。所有半径使用同一个乘积，只是各自的有限头部长度不同。零点集合若在平面内有聚点，上述有限性马上失去；恒等定理同时说明非恒零整函数也不可能实现这类数据。

\ProseParagraphBreak

存在性没有给出规范化的唯一性。指定 \(P(0)=1\) 只消除了常数乘法自由度，仍可乘上 \(e^{az}\)；再指定 \(P'(0)=0\) 也仍可乘上 \(e^{bz^2}\)。只有增长约束或更多结构条件才会继续减少自由度。在有限阶一的情形，两个同零点函数分别有指数一次多项式的分解，因此它们的商为 \(e^{az+b}\)；这时指定原点值及一次对数导数才足够。若原点是零点，应先除去共同的原点零因子，再施加这些条件。

% !TeX root = ../../Book3.tex
\section{Mittag–Leffler 定理}
\label{b3:sec:2004}

乘法构造指定零点，加法构造则指定极点附近的负幂项。把每个远处极点的低阶 Taylor 项扣掉，可以沿用同样的尾部控制思想。

\subsection{主要部分}
\label{b3:subsec:200401}

亚纯函数在极点 \(a\) 的 Laurent 负幂有限和 \(Q_a(z)=\sum_{k=1}^{m_a}c_{a,k}(z-a)^{-k}\) 称为\term[zhuyoubufen]{主要部分}[principal part]。剩余部分在 \(a\) 全纯。指定主要部分比只指定极点阶数强，因为它固定了每个负幂系数及留数。

\subsection{亚纯插值}
\label{b3:subsec:200402}

设 \(a_n\to\infty\)，给定主要部分 \(Q_n\)。选 \(R_n<|a_n|\)，且 \(R_n\to\infty\)，在圆盘 \(|z|\le R_n\) 上用 \(Q_n\) 关于零的 Taylor 多项式 \(T_n\) 逼近到误差 \(2^{-n}\)。这可行，因为 \(Q_n\) 在半径 \(|a_n|\) 内全纯，Taylor 级数在更小闭圆盘一致收敛。

于是 \(\sum(Q_n-T_n)\) 在避开极点的每个紧集上正规收敛。每个指定极点附近，除该项之外的尾部全纯，故其主要部分不被更改。

\subsection{Mittag–Leffler 定理}
\label{b3:subsec:200403}

\begin{theorem}[Mittag--Leffler：整平面版本]\label{b3:thm:20-mittag-leffler}
任意无有限聚点的极点集及其有限主要部分，都由一个复平面亚纯函数实现。两个实现函数之差为整函数。
\end{theorem}
\begin{proof}
除去可能在原点的有限主要部分后，用上一小节的 Taylor 扣除构造并相加。局部正规收敛给亚纯性，极点邻域中的有限项分离给准确主要部分。两解的差在每个极点负幂全部抵消，因此所有极点可去，成为整函数。
\end{proof}
一般平面区域上的相应定理需要把多项式扣除改成适应补集的有理逼近；本节给出的完整构造范围是整平面，不把这一证明无说明地扩展到任意区域。

\subsection{与 Weierstrass 定理的对偶性}
\label{b3:subsec:200404}

Weierstrass 定理留下无零点的乘法自由度 \(e^g\)，Mittag--Leffler 定理留下整函数的加法自由度。两者还能组合给整插值：取简单零点为 \(a_n\) 的整函数 \(P\)，给主要部分 \(b_n/(P'(a_n)(z-a_n))\)，构造亚纯函数 \(G\)，则 \(F=PG\) 整且 \(F(a_n)=b_n\)。局部乘法直接核验值，任意增长的插值数据也不破坏离散点上的存在性。


主要部分构造也有误差与有效范围的区别。对每个 \(Q_n\) 减去有限 Taylor 多项式，减去的项不引入新极点，并且改变的是整函数自由度。若只减常数项，在近原点的紧集上尾部可能仍不收敛；所需次数应根据该项大小选择，而不是统一宣称一次修正足够。

\ProseParagraphBreak

整插值还可以指定有限阶导数。设在 \(a_n\) 希望 \(F\) 与给定多项式芽 \(J_n\) 的差至少有 \(m_n\) 阶零点。取 \(P\) 在 \(a_n\) 有恰好 \(m_n\) 阶零点，令 \(Q_n\) 是 \(J_n/P\) 的主要部分。由 Mittag--Leffler 构造 \(G\)，则 \(G-J_n/P\) 在 \(a_n\) 全纯，故 \(F=PG\) 满足 \(F-J_n\) 被 \((z-a_n)^{m_n}\) 整除。负幂计算准确保留了全部所需导数；不需要把每阶导数当作独立极点重新构造。

\ProseParagraphBreak

这种插值没有控制全平面增长。若要求在整数上取值 \(e^{n^2}\) 且整函数阶小于二，就会与最大模的下界矛盾；若指定的值只是有界，也不能自动保证存在有界整函数，因为 Liouville 使有界整函数只能恒常。存在性定理容许任意数据，正是因为它没有许诺这些额外估计。


\ProseParagraphBreak

下面把 Taylor 扣除步骤写成一个可检查的例子。在正整数 \(n\) 指定二阶主要部分 \(c_n/(z-n)^2\)，其中 \(c_n\) 可以任意迅速增长。对 \(|z|<n\)，有
\[
\frac{c_n}{(z-n)^2}=\frac{c_n}{n^2}\sum_{k\ge0}(k+1)(z/n)^k.
\]
取 \(R_n=n/2\)，选择整数 \(d_n\) 使
\[
\frac{|c_n|}{n^2}\sum_{k>d_n}(k+1)2^{-k}<2^{-n}.
\]
有限截断 \(T_n\) 为整多项式，故 \(\sum_n(c_n/(z-n)^2-T_n(z))\) 在任意避开正整数的紧集上正规收敛。对任意固定 \(R\)，仅有限个 \(R_n<R\)；剩余项由 \(2^{-n}\) 一致控制。由于没有一次负幂项，构造出的函数在每个极点的留数仍为零，不能因为修正次数高而凭空出现留数。

\ProseParagraphBreak

若 \(c_n=1\)，可以直接用 \(\sum_{n\ge1}(z-n)^{-2}\)，因为远处尾部已经是 \(O_K(n^{-2})\)；若 \(c_n=n^2\)，直接尾项趋于一，减去常数项后尾项仍有 \(O_K(n^{-1})\)，还须减一次项。具体地，
\[
\frac{n^2}{(z-n)^2}-1-\frac{2z}{n}=O_K(n^{-2}).
\]
由此可见，扣除次数应从具体尾项的展开决定。一般定理提供可用算法，特殊数据则可用更低次数的修正。

\ProseParagraphBreak

Mittag--Leffler 构造的唯一性比乘积构造更容易表达：具有相同极点主要部分的两个亚纯函数之差在每个极点可去，因而为整函数。反过来，加上任意整函数都不改变主要部分。因此“主要部分决定函数”只有在另有增长、周期性或规范化时才可能成立。对有限个极点，所有主要部分之和已经给出有理函数，但一般解仍是这个有理函数加任意整函数；极点数有限并不使函数必为有理。

\ProseParagraphBreak

插值点必须离散。例如若在 \(1/n\) 指定值零，同时在原点指定值一，任何全纯函数的连续性已产生矛盾；即使原点也指定零，恒等定理仍迫使函数恒为零，不能再在另一点任意指定非零值。有限阶导数插值也没有绕过这个限制。所谓“任意数据”准确地指在没有有限聚点的点列上，每一点只提出有限个导数条件。

% !TeX root = ../../Book3.tex
\section{Hadamard 分解}
\label{b3:sec:2005}

一般 Weierstrass 分解中的 \(g\) 可以任意整。若原函数有限阶，增长约束会把它压成多项式；这是 Hadamard 分解比单纯零点存在性更强的部分。

\subsection{有限阶整函数}
\label{b3:subsec:200501}

设整函数 \(f\) 阶为 \(\rho<\infty\)，非零零点 \(a_n\)，原点重数 \(m\)，令 \(p=\lfloor\rho\rfloor\)。零点计数 \(n(r)=O_\varepsilon(r^{\rho+\varepsilon})\)，故 \(\sum|a_n|^{-p-1}<\infty\)（取 \(\varepsilon<p+1-\rho\)）。因此 \(P=z^m\prod E_p(z/a_n)\) 存在，且 \(f/P=e^g\)，\(g\) 整。问题归结为证明 \(g\) 的导数高阶消失。

\subsection{典范乘积}
\label{b3:subsec:200502}

先说明所需的乘积估计。对每个小 \(\varepsilon>0\)，可以在每个充分大的区间 \([r,2r]\) 选一个半径 \(R\)，使
\[
\log|P(z)|\ge-C_\varepsilon R^{\rho+\varepsilon},\quad |z|=R.
\]
选半径时避开所有 \(|a_n|\le4r\) 的半径 \(r^{-\rho-2}\) 邻域；这些禁区总长为 \(O(r^{-2+\varepsilon/4})\)，小于 \(r\)。对内层零点，\(|1-z/a_n|\) 至少为所选间隔除 \(|a_n|\)，故对数下界为 \(-C\log r\)；有限的极小零点也满足同类下界。指数修正的 \(k\) 次项由分部求和控制，\(k\le p\le\rho\)，总量为 \(O(r^{\rho+\varepsilon/4}\log r)\)。外层 \(|a_n|>4r\) 用 \(|\log E_p(z/a_n)|\le2|z/a_n|^{p+1}\)，分部求和同样给 \(O(r^{\rho+\varepsilon/4})\)。最后用 \(\log r=O(r^{\varepsilon/2})\) 吸收对数。上界也由相同的分层估计得到。

关键不是逐点避开零点，而是选出整条无过近零点的圆周。否则 \(\log|P|\) 在零点处为负无穷，不能估计商的增长。

\subsection{Hadamard 分解定理}
\label{b3:subsec:200503}

\begin{theorem}[Hadamard 分解]\label{b3:thm:20-hadamard}
非恒零有限阶整函数 \(f\) 满足
\[
f(z)=z^m e^{Q(z)}\prod_nE_p(z/a_n),\quad p=\lfloor\rho(f)\rfloor,
\]
其中 \(Q\) 为次数不超过 \(\rho(f)\) 的多项式。
\end{theorem}
\begin{proof}
在上一小节选出的圆周，\(\operatorname{Re}g=\log|f|-\log|P|\le C_\varepsilon R^{\rho+\varepsilon}\)。调和最大值原理把界传到圆盘内部。
为把实部界变成模界，令 \(F=g-g(0)\)、\(A=\max_{|z|\le R}\operatorname{Re}F\)。若 \(A=0\)，最大值原理使 \(g\) 常数。若 \(A>0\)，函数 \(F/(2A-F)\) 在圆盘模不超过一，中心为零；Schwarz 引理给半径 \(R/2\) 上模不超过 \(1/2\)，解出 \(|F|\le2A\)。所以 \(\max_{|z|\le R/2}|g(z)|\le C_\varepsilon R^{\rho+\varepsilon}+|g(0)|\)。
Cauchy 估计对整数 \(k>\rho\) 给 \(|g^{(k)}(0)|\le k!C_\varepsilon R^{\rho+\varepsilon-k}\)。选 \(0<\varepsilon<k-\rho\)，沿选定半径趋无穷，得到导数零。于是 \(g=Q\) 为要求的多项式。
\end{proof}

\subsection{属与阶}
\label{b3:subsec:200504}

整函数的属通常取典范乘积属与 \(Q\) 次数的较大者。上述定理给属不超过有限阶；零点指数与指数多项式共同决定增长，不能只数零点。

作为应用，\(\sin\pi z\) 阶一、零点为全体整数。将正负因子配对，Hadamard 给
\[
\sin\pi z=e^{az+b}z\prod_{n\ge1}\left(1-\frac{z^2}{n^2}\right).
\]
两侧的奇性使 \(e^{az}\) 为偶函数，从 \(e^{az}=e^{-az}\) 对零点求导得到 \(a=0\)；再比较原点导数得 \(e^b=\pi\)。因此经典正弦乘积得到完整证明。


进一步展开乘积估计中的分部求和。若 \(n(t)\le Ct^a\)、\(a=\rho+\delta\)，则对 \(k\le p\le\rho\)，
\[
\sum_{1\le|a_n|\le4r}|a_n|^{-k}
\le (4r)^{-k}n(4r)+k\int_1^{4r}n(t)t^{-k-1}\,\mathrm{d}t
\le C_{k,\delta}r^{a-k}.
\]
乘上 \(r^k\) 后每阶修正均为 \(O(r^a)\)。当 \(p=0\) 没有修正项，内层只剩对数因子。对外层，取 \(a<p+1\)，同样得到
\[
\sum_{|a_n|>4r}|a_n|^{-p-1}\le
(p+1)\int_{4r}^{\infty}n(t)t^{-p-2}\,\mathrm{d}t
\le C r^{a-p-1}.
\]
这解释了为什么一定要先选 \(\delta<p+1-\rho\)：尾积分的指数只有在这个范围才可积。整数阶 \(\rho=p\) 仍有一整次的余量，而不是临界发散。

\ProseParagraphBreak

另一个可能的误用是先估计 \(|f/P|\)，再把分子分母的阶简单相减。整函数的阶并不这样运算；商在零点附近会遇到很小的分母，甚至两个同阶函数相减也可能造成完全抵消。本证明专门选择圆周上的乘积下界，再利用最大值原理，正是为了把分母风险变成可核验的估计。有限阶条件没有直接赋予商同样的阶。

\ProseParagraphBreak

实部转模的步骤是 \term[borel-caratheodory]{Borel--Carathéodory 估计}[Borel--Carathéodory estimate]的特例。更一般地，\(g\) 在闭半径 \(R\) 的圆盘全纯、\(A=\max\operatorname{Re}(g-g(0))\) 时，\(0<r<R\) 上有
\[
|g(z)-g(0)|\le\frac{2Ar}{R-r}.
\]
由同一分式 \(F/(2A-F)\) 与 Schwarz 估计 \(r/R\) 解出即可。此估计不是由 \(|\operatorname{Re}g|\) 控制全部函数；一个纯虚常数仍需由 \(g(0)\) 另外记录。 Hadamard 证明中取 \(r=R/2\)，使这一步的常数不随大半径改变。

\ProseParagraphBreak

有限阶且只有有限零点的整函数现在具有完全明确的形式 \(P e^Q\)。例如 \((z-1)^2e^{3z^2+z}\) 的多项式部分固定两个零点，指数二次项决定阶二、型三，一次项不改变阶和型。无零点整函数 \(e^{z^3}\) 的阶三全部来自指数因子。对 \(\sin\pi z\)，阶一由零点密度和双向指数增长同时确认；它没有附加线性指数因子，原因不是零点本身，而是奇偶对称。

\ProseParagraphBreak

在整数极点的对称例子中，配对能减少扣除。由正弦乘积的对数导数得
\[
\pi\cot\pi z=\frac1z+\sum_{n\ge1}\left(\frac1{z-n}+\frac1{z+n}\right)
=\frac1z+\sum_{n\ge1}\frac{2z}{z^2-n^2}.
\]
最后级数在避开整数的紧集正规收敛，因为尾部为 \(O_K(n^{-2})\)。把 \(z-n\) 与 \(z+n\) 两个极点单独按不对称顺序相加，却只得到 \(O(n^{-1})\) 的尾部。明确的配对是收敛证明的一部分，不能藏在省略号中。这个公式给所有整数留数一，同时把无穷多个局部极点资料汇总成一个周期函数。

\ProseParagraphBreak

最后检验两个同零点有限阶函数之间的关系。若 \(f_1,f_2\) 均非恒零、阶不超过 \(\rho\)，按同一零点列及同一 \(p=\lfloor\rho\rfloor\) 取乘积，分别有 \(f_j=P e^{Q_j}\)，\(\deg Q_j\le\rho\)。所以 \(f_1/f_2=e^{Q_1-Q_2}\)。这里并非先假设商有限阶，而是由两份分解共同推出。特别地，当 \(\rho<1\) 时商只能为常数；当 \(\rho=1\) 时仍可有 \(e^{az}\) 自由度。

\ProseParagraphBreak

阶零也包含非多项式整函数，不能把它误写成“阶零当且仅当多项式”。例如零点 \(a_n=e^{n^2}\) 的属零乘积满足
\[
\log M(r)\le\sum_{n\ge1}\log(1+r/e^{n^2})=O((\log r)^{3/2}),\quad r\ge e.
\]
取 \(N=\lfloor\sqrt{\log r}\rfloor\)，前 \(N+1\) 项各不超过 \(\log(1+r)\)，尾部由 \(r\sum_{n>N+1}e^{-n^2}\) 控制为有界量，故得到该估计。于是阶为零，但函数有无穷多个零点而不可能是多项式。Hadamard 在此只断言指数因子为常数；无穷乘积本身仍保留全部零点。有限阶与有限个零点是两项独立条件。


\ProseParagraphBreak

Hadamard 定理中的 \(p=\lfloor\rho\rfloor\) 是统一可用的选择，未承诺最小属。若零点临界幂和已经收敛，某些因子可降阶；降阶会改变指数多项式的低次项。例如正整数的属一乘积与一个事先配对的正负整数乘积拥有不同的指数修正，虽然最终零点可相同。比较两种表示时，应先把初等因子的指数差逐项展开，证明相关系数和收敛，再把它吸收进指数多项式。没有可和性就不能把一串指数形式地合并成一个有限系数。

\ProseParagraphBreak

正弦乘积还有一个实变量误差例子。对固定 \(|z|\le r\)，当 \(N>2r\) 时，尾因子 \(1-z^2/n^2\) 的对数绝对值总和不超过 \(2r^2\sum_{n>N}n^{-2}\le2r^2/N\)。于是截断乘积与完整乘积的尾部相对误差不超过 \(e^{2r^2/N}-1\)。离原点和整数零点很近时，函数自身可能小，相对全函数的计算误差还需考虑有限头部；在零点上，只能比较提出零因子后的函数。这给出了从解析恒等式到实际近似之间需要保留的条件。

\ProseParagraphBreak

由正弦乘积与倍角恒等式还可得到余弦乘积。将 \(\sin(2\pi z)/(2\sin\pi z)\) 按可去极点解释为 \(\cos\pi z\)，并在正规收敛的乘积中约去偶数零点因子，得
\[
\cos\pi z=\prod_{n\ge1}\left(1-\frac{4z^2}{(2n-1)^2}\right).
\]
约因子的合法性先在避开整数的紧集用非零尾部完成，再由全纯恒等定理扩展；不能在每个零点处直接把两个零乘积相除。原点值一确定常数因子，偶性排除一次指数自由度。该例同时展示零点删选、对称性和规范化怎样从一般分解确定一个具体函数。

% 本章小结（不计入编号）
\begin{chaptersummary}
对数可和性保证乘积尾部非零；初等因子使任意离散零点都能实现。扣除 Taylor 多项式把同一思想转成主要部分的加法构造。Hadamard 分解的额外责任是估计无零点商的实部，再用 Cauchy 估计证明其对数为多项式。正弦乘积展示了零点、增长与对称性共同确定特殊函数。
\end{chaptersummary}


\BookExercises

% 本题组自拟；各题新增训练动作见题意，未复制第三方原题。
\begin{exercise}\label{b3:ex:20-draft-01}
说明 \(\prod_{n\ge2}(1-z/n^2)\) 是否正规收敛，列出零点。
\end{exercise}
\begin{solution}
固定紧集半径 \(R\) 时 \(\sum R/n^2<\infty\)，正规乘积定理给整函数。有限零因子与非零尾部的分离说明零点恰为 \(n^2\)、\(n\ge2\)，均简单。
\end{solution}

\begin{exercise}\label{b3:ex:20-draft-02}
对零点 \(a_n=\sqrt n\) 求收敛指数与最小可用属。
\end{exercise}
\begin{solution}
幂和为 \(\sum n^{-s/2}\)，收敛当且仅当 \(s>2\)，所以指数二。\(p+1>2\) 才收敛，最小整数 \(p=2\)。
\end{solution}

\begin{exercise}\label{b3:ex:20-draft-03}
证明具有相同零点和重数的两个非恒零整函数之商为整指数。
\end{exercise}
\begin{solution}
商在共同零点处可去且无零点。对数导数整，在整平面有原函数 \(g\)；商乘 \(e^{-g}\) 导数零为非零常数，再把常数选一复对数并吸收进 \(g\)。
\end{solution}

\begin{exercise}\label{b3:ex:20-draft-04}
构造在每个整数 \(n\) 留数为一、没有其他极点的亚纯函数。
\end{exercise}
\begin{solution}
\(\pi\cot\pi z\) 在整数附近为 \(1/(z-n)+\)全纯项，所以符合。它由正弦、余弦整函数之商定义，其他点正弦非零，故无额外极点。
\end{solution}

\begin{exercise}\label{b3:ex:20-draft-05}
解释为何任何给定复数列 \(b_n\) 都能作为某整函数在正整数处的值。
\end{exercise}
\begin{solution}
正整数离散。取简单零点恰为正整数的整函数 \(P\)，用 Mittag--Leffler 构造主要部分 \(b_n/(P'(n)(z-n))\) 的 \(G\)。乘积 \(PG\) 的各极点可去，值为 \(b_n\)，没有增长限制。
\end{solution}

\begin{exercise}\label{b3:ex:20-draft-06}
证明无零点有限阶整函数形为 \(e^Q\)，且阶正数时必为整数。
\end{exercise}
\begin{solution}
Hadamard 乘积为空，且无原点零点，故为指数多项式。非恒定 \(Q\) 的次数 \(d\) 给最大模对数 \(|a_d|r^d+O(r^{d-1})\)，阶为整数 \(d\)。常数对应阶零。
\end{solution}

\begin{exercise}\label{b3:ex:20-draft-07}
由正弦乘积计算 \(\sum_{n\ge1}n^{-2}\)。
\end{exercise}
\begin{solution}
除以 \(\pi z\)，原点二阶系数由正弦级数为 \(-\pi^2/6\)。正规乘积展开的二阶系数为 \(-\sum n^{-2}\)，系数比较得 \(\pi^2/6\)。逐项取系数由局部一致收敛及 Cauchy 公式合法。
\end{solution}

\begin{exercise}\label{b3:ex:20-draft-08}
说明 \(e^{e^z}\) 有无零点，并判断其是否能写成指数多项式。
\end{exercise}
\begin{solution}
指数无零点，但函数无穷阶。若为 \(e^Q\)，其阶将为多项式次数，有限，矛盾。这显示 Hadamard 的有限阶条件不能删。
\end{solution}

% !TeX root = ../../Book3.tex
\chapter{\texorpdfstring{\(\Gamma\)}{\AnalysisPdfGamma} 函数与 Mellin 方法}
\label{b3:ch:21}
\BookChapterNumber{b3:ch:21}
\begin{chapterabstract}
Gamma 函数把阶乘递推、参数积分与整函数乘积联系起来。本章从 Euler 积分的两端控制出发，证明全纯性、递推和严格对数凸性，用 Bohr--Mollerup 定理说明正实轴上的唯一性，再通过递推和 Taylor 扣除完成亚纯延拓。Euler 极限确定倒数乘积的规范化，进而给出无零点性、对数导数、反射与乘法公式；正实轴 Stirling 公式则从积分峰值附近的尺度得到增长常数。最后转入 Mellin 变换，借对数坐标明确权重、Fourier 归一化与反演条件，并以 Gamma 的留数展开说明端点渐近如何成为变换的极点。
\end{chapterabstract}

\begin{learninggoals}
\begin{itemize}
\item 分别控制 Euler 积分在零与无穷远的行为，验证参数微分，并说明递推、正性、规范化与对数凸性各自承担的作用。
\item 用递推粘合亚纯延拓，计算极点留数及有限部分，区分收敛积分、Taylor 正则化与延拓后的函数值。
\item 从 Euler 极限得到倒数乘积，先证明非零性再取对数导数，核对 Euler 常数和局部幂和系数的来源。
\item 保留 Beta 换元的 Jacobi 因子以及反射、倍增、乘法公式的规范化常数，区分正实轴渐近和特殊竖线估计。
\item 从两端可积性确定 Mellin 收敛带，用对数坐标解释缩放、权重、乘法卷积与 Fourier 变换之间的关系。
\item 分开验证绝对反演、有限移线和余项极限的条件，追踪幂次、对数因子、极点位置与极点阶数的对应。
\end{itemize}
\end{learninggoals}

若只要求 \(F(n+1)=n!\) 和 \(F(z+1)=zF(z)\)，阶乘在复平面中的延拓并不唯一：乘上适当的一周期函数，这些条件仍然成立。Euler 积分提供了一个自然候选，但要理解它为何成为通常的 Gamma 函数，还须辨明积分带来的正性、凸性和增长信息。另一方面，同一个积分中的幂 \(t^{z-1}\) 也能用来分析一般函数；把 \(e^{-t}\) 换成 \(f(t)\)，就得到本章后半部的 Mellin 变换。

\ProseParagraphBreak
\secref{b3:sec:2101}先处理原积分。原点附近由实部的正下界控制，无穷远由指数下降控制；加入对数幂后仍可积，才允许逐次对参数求导。递推来自边界项消失的分部积分，Gaussian 积分确定半整数处的常数。阅读 Bohr--Mollerup 证明时，应注意凸性比较的是哪些割线，以及整数点已知值怎样夹住非整数处的值。固定实平移的 Gamma 比值可以直接从这些凸性估计得到，不必先使用 Stirling 公式。

\ProseParagraphBreak
\secref{b3:sec:2102}把递推转为向左延拓的工具。不同移动次数必须在重叠区域给出同一个函数；分母零点处则要检查分子是否也消失，才能确认极点阶数。Taylor 扣除提供另一种具体表示：先减去造成发散的低次项，再把它们的积分写成明确的有理奇部。这个过程解释了延拓值从何而来，却不使原先发散的 Euler 积分自动收敛。

\ProseParagraphBreak
\secref{b3:sec:2103}从有限积分推得 Euler 极限，再整理出 \(1/\Gamma\) 的正规乘积。一次指数修正保证尾项可和，Euler 常数记录有限乘积规范化的极限。这里先比较两个有限表达式及其极限，证明 \(\Gamma\) 无零点，然后才使用对数导数。由此得到的 \(\psi_0\)、高阶导数与局部幂和，也能回到参数积分中解释对数矩和曲率。

\ProseParagraphBreak
\secref{b3:sec:2104}将这些表示互相连接。相反参数的乘积给反射公式，Beta 积分的二维换元给倍增关系，递推与对数凸性则固定一般乘法公式中的函数自由度。Stirling 证明另从积分的峰值出发：在宽度为 \(\sqrt{x}\) 的中部出现 Gaussian 极限，两侧尾部需要各自估计。该证明只建立正实轴版本；整数和半整数竖线上的模值由反射与递推独立求出，不以未证明的复参数渐近代替。

\ProseParagraphBreak
\secref{b3:sec:2105}沿每条实部为 \(c\) 的竖线，把 Mellin 变换写成 \(e^{cu}f(e^u)\) 的 Fourier 变换。改变竖线也就改变权重，因此收敛带属于变换资料的一部分。绝对反演要求竖线上的变换可积；围道移动还要控制水平边；若继续让新竖线趋向左方，又需一条随移动次数变化的余项界。指数函数的例子将这三步分别完成，最后再把端点的幂与对数展开解释为亚纯变换的主部。

\ProseParagraphBreak
本章使用参数积分、支配收敛、全纯函数的局部一致极限与留数定理，以及\chapref{b3:ch:20}的正规乘积和正弦乘积。Mellin 反演所需的 Fourier 归一化会在使用处重新写明。Stein、Shakarchi 的《Complex Analysis》\cite{Stein2003Complex}可供继续阅读特殊函数与后续数论问题的联系；本章所用的 Gamma 恒等式、反演条件和移线估计均在正文中建立。

% !TeX root = ../../Book3.tex
\section{Euler \texorpdfstring{\(\Gamma\)}{\AnalysisPdfGamma} 函数}
\label{b3:sec:2101}

阶乘满足乘法递推。若要把整数变量换成实变量或复变量，需要先找到一个在整数处给出阶乘、同时又能对参数作分析运算的表达。Euler 积分提供了这样的函数，但递推和整数点取值并不足以唯一确定它。本节先从积分得到全纯性、递推和半整数值，再用正实轴上的对数凸性建立唯一性与平移比值估计。

\subsection{Euler 积分}
\label{b3:subsec:210101}

积分 \(\int_0^\infty e^{-t}t^n\,\mathrm{d}t\) 中，分部积分会把幂次降低一，并产生因子 \(n\)。为了使递推写成变量本身乘以函数，定义时把幂次取为 \(z-1\)。

\Needspace{6\baselineskip}
\begin{definition}\label{b3:def:21-gamma}
对 \(z\in\mathbb C\)、\(\operatorname{Re}z>0\)，定义\term[gamma-hanshu]{Gamma 函数}[gamma function]
\[
 \Gamma(z)\coloneqq\int_0^\infty e^{-t}t^{z-1}\,\mathrm{d}t,
 \quad t^{z-1}=\exp\bigl((z-1)\log t\bigr),\quad t>0,
\]
其中 \(\log t\) 是正实轴上的实对数。
\end{definition}
\symidx{\(\Gamma(z)\)：Gamma 函数}
积分变量只取正实数，所以复幂在这里没有绕过原点选择分支的问题。若 \(z=\sigma+\mathrm{i}y\)，则 \(|t^{z-1}|=t^{\sigma-1}\)，虚部只产生模为一的振荡因子。

\ProseParagraphBreak

两端的收敛机制不同。在 \((0,1)\) 上，\(e^{-1}\le e^{-t}\le1\)，所以绝对可积恰好要求 \(\sigma>0\)。在 \([1,\infty)\) 上，指数衰减压过任意固定实次幂，对所有实数 \(\sigma\) 都可积。因此 \(\operatorname{Re}z>0\) 是这个积分的绝对收敛域。要在积分号下对 \(z\) 求导，还须控制每次求导带来的 \(\log t\)。

\Needspace{6\baselineskip}
\begin{proposition}[全纯性与参数导数]\label{b3:prop:21-euler-derivatives}
Euler 积分定义的 \(\Gamma\) 在半平面 \(\operatorname{Re}z>0\) 全纯。对每个非负整数 \(k\)，有
\[
 \Gamma^{(k)}(z)
 =\int_0^\infty e^{-t}t^{z-1}(\log t)^k\,\mathrm{d}t,
 \quad \operatorname{Re}z>0,
\]
且这些积分在每个紧参数集上绝对一致收敛。此外，\(\Gamma(x)>0\) 对 \(x>0\) 成立，并且
\[
 |\Gamma(\sigma+\mathrm{i}y)|\le\Gamma(\sigma),
 \quad \sigma>0,\quad y\in\mathbb R.
\]
\end{proposition}
\begin{proof}
固定紧集 \(K\subset\{\operatorname{Re}z>0\}\)，选 \(0<a\le\operatorname{Re}z\le b\) 对所有 \(z\in K\) 成立。在 \(0<t<1\) 上，
\[
 |e^{-t}t^{z-1}(\log t)^k|
 \le t^{a-1}|\log t|^k.
\]
代换 \(u=-\log t\)，再作有限次分部积分，得到
\[
 \int_0^1t^{a-1}|\log t|^k\,\mathrm{d}t
 =\int_0^\infty e^{-au}u^k\,\mathrm{d}u
 =\frac{k!}{a^{k+1}}<\infty.
\]
在 \(t\ge1\) 上，则有
\[
 |e^{-t}t^{z-1}(\log t)^k|
 \le e^{-t}t^{b-1}(\log t)^k
 \le C_{b,k}e^{-t/2}.
\]
最后一个常数存在，因为 \(t^{b-1}(\log t)^k e^{-t/2}\) 在 \([1,\infty)\) 有界。这些控制函数不依赖 \(z\in K\)。

截断积分 \(G_N(z)=\int_{1/N}^{N}e^{-t}t^{z-1}\,\mathrm{d}t\) 是整函数，在有限积分区间可以逐次对参数求导。上述共同控制表明，\(G_N\) 及其各个固定阶导数的积分表达在 \(K\) 上一致收敛。由\cref{b3:prop:04-derivative-convergence}，极限全纯，导数正是陈述中的积分。正实参数时，被积函数处处为正，故 \(\Gamma(x)>0\)。最后，对 Euler 积分使用三角不等式即得竖线上的模估计。
\end{proof}
这里的一致性针对紧参数集；当实部趋近零时，共同控制中的 \(a\) 也会趋近零，不能保留同一个可积上界。右半平面上的竖线估计只给出有界性，尚未给出随 \(|y|\) 增大的衰减速度。

\subsection{函数方程}
\label{b3:subsec:210102}

积分定义使变量平移一步恰好对应一次分部积分。端点项是否消失，仍由原点的幂次和无穷远的指数分别控制。

\Needspace{6\baselineskip}
\begin{proposition}[递推与归一化]\label{b3:prop:21-recurrence}
对所有 \(z\in\mathbb C\)、\(\operatorname{Re}z>0\)，有
\[
 \Gamma(z+1)=z\Gamma(z),\quad \Gamma(1)=1.
\]
因此，对每个非负整数 \(n\)，\(\Gamma(n+1)=n!\)。
\end{proposition}
\begin{proof}
在 \([\varepsilon,R]\) 上分部积分，得到
\[
 \int_\varepsilon^R e^{-t}t^z\,\mathrm{d}t
 =\bigl[-e^{-t}t^z\bigr]_\varepsilon^R
  +z\int_\varepsilon^R e^{-t}t^{z-1}\,\mathrm{d}t.
\]
设 \(\sigma=\operatorname{Re}z>0\)。无穷远端的边界项模为 \(e^{-R}R^\sigma\to0\)，原点端的模为 \(e^{-\varepsilon}\varepsilon^\sigma\to0\)。两个积分都绝对收敛，所以可分别令 \(R\to\infty\)、\(\varepsilon\downarrow0\)，得到递推式。直接计算 \(\int_0^\infty e^{-t}\,\mathrm{d}t=1\) 给出归一化，再作整数次递推得到阶乘。
\end{proof}
归一化选定了乘法常数，却没有排除一周期因子。例如，固定 \(0<\eta<1\)，函数
\[
 F(z)=\Gamma(z)\bigl(1+\eta\sin(2\pi z)\bigr)
\]
在右半平面全纯，满足同一递推、同一整数点取值，而且在正实轴仍为正。它与 \(\Gamma\) 不同，因此即使同时要求全纯性和正实轴上的正性，递推仍不能给出唯一性。

\ProseParagraphBreak

递推还把半整数值归约到一个起始值。这个值由 Gaussian 积分确定，适合在讨论渐近公式之前独立算出。

\Needspace{6\baselineskip}
\begin{proposition}[Gaussian 积分与半整数值]\label{b3:prop:21-half-integers}
对每个非负整数 \(n\)，有
\[
 \Gamma\left(\frac12\right)=\sqrt\pi,
 \quad
 \Gamma\left(n+\frac12\right)
 =\frac{(2n)!}{4^n n!}\sqrt\pi.
\]
\end{proposition}
\begin{proof}
令 \(I=\int_{\mathbb R}e^{-u^2}\,\mathrm{d}u\)。这个积分有限且为正：在 \(|u|\ge1\) 上可由 \(e^{-|u|}\) 控制，在有限区间被积函数有界。被积函数非负，Tonelli 定理给
\[
 I^2=\int_{\mathbb R^2}e^{-(u^2+v^2)}\,\mathrm{d}u\,\mathrm{d}v.
\]
在圆环上使用极坐标换元，Jacobi 因子为 \(r\)，再让内半径趋零、外半径趋于无穷。非负性允许用单调收敛取极限，于是
\[
 I^2=\int_0^{2\pi}\int_0^\infty e^{-r^2}r\,\mathrm{d}r\,\mathrm{d}\theta
 =2\pi\cdot\frac12=\pi.
\]
由 \(I>0\) 确定 \(I=\sqrt\pi\)。在 Euler 积分中令 \(t=u^2\)，得到
\[
 \Gamma\left(\frac12\right)
 =2\int_0^\infty e^{-u^2}\,\mathrm{d}u=I.
\]
最后，重复使用递推式，
\[
 \Gamma\left(n+\frac12\right)
 =\left(\prod_{j=0}^{n-1}\left(j+\frac12\right)\right)\sqrt\pi
 =\frac{1\cdot3\cdots(2n-1)}{2^n}\sqrt\pi
 =\frac{(2n)!}{4^n n!}\sqrt\pi.
\]
\(n=0\) 时有限乘积为空，公式仍成立。
\end{proof}
同一换元给出 Gaussian 偶数矩
\[
 \int_{\mathbb R}u^{2n}e^{-u^2}\,\mathrm{d}u
 =\Gamma\left(n+\frac12\right),\quad n\ge0.
\]
奇数矩因被积函数为奇函数而为零；这些积分均绝对收敛。阶乘和半整数值由此分别对应指数权重与 Gaussian 权重下的整数矩。

\subsection{对数凸性}
\label{b3:subsec:210103}

Euler 积分中的参数出现在指数 \(t^{x-1}\) 上，因此参数的凸组合会变成两个正被积函数的几何平均。Hölder 不等式正好比较这种几何平均的积分。

\Needspace{6\baselineskip}
\begin{proposition}[严格对数凸性]\label{b3:prop:21-logconvex}
Euler 积分定义的实函数 \(x\mapsto\log\Gamma(x)\) 在 \((0,\infty)\) 上严格凸。即对 \(x,y>0\)、\(x\ne y\) 和 \(0<\lambda<1\)，
\[
 \Gamma(\lambda x+(1-\lambda)y)
 <\Gamma(x)^\lambda\Gamma(y)^{1-\lambda}.
\]
\end{proposition}
\begin{proof}
令 \(A(t)=e^{-t}t^{x-1}\)、\(B(t)=e^{-t}t^{y-1}\)。二者严格为正且可积，并且
\[
 e^{-t}t^{\lambda x+(1-\lambda)y-1}
 =A(t)^\lambda B(t)^{1-\lambda}.
\]
对右边使用指数 \(1/\lambda\) 与 \(1/(1-\lambda)\) 的 Hölder 不等式，得
\[
 \int_0^\infty A^\lambda B^{1-\lambda}\,\mathrm{d}t
 \le\left(\int_0^\infty A\,\mathrm{d}t\right)^\lambda
      \left(\int_0^\infty B\,\mathrm{d}t\right)^{1-\lambda}.
\]
在两个非零正函数之间，等号要求 \(A\) 与 \(B\) 几乎处处成比例，即 \(t^{x-y}\) 几乎处处为常数。但 \(x\ne y\) 时该函数在 \((0,\infty)\) 严格单调，不可能满足这个条件。因此不等式严格，取对数即得结论。
\end{proof}
还可以把严格性读成一个二阶导数公式。固定 \(x>0\)，以
\[
 w_x(t)=\frac{e^{-t}t^{x-1}}{\Gamma(x)},\quad t>0,
\]
作为概率密度，并令 \(m_x=\int_0^\infty(\log t)w_x(t)\,\mathrm{d}t\)。参数导数公式保证这些对数矩有限，且
\[
 \frac{\mathrm{d}^2}{\mathrm{d}x^2}\log\Gamma(x)
 =\frac{\Gamma''(x)}{\Gamma(x)}
  -\left(\frac{\Gamma'(x)}{\Gamma(x)}\right)^2
 =\int_0^\infty(\log t-m_x)^2w_x(t)\,\mathrm{d}t>0.
\]
最后一项是 \(\log t\) 的方差；密度在整个正半轴严格为正，而 \(\log t\) 不恒定，所以方差严格为正。归一化的分母 \(\Gamma(x)\) 是这一步所必需的，单独观察 \(\Gamma''(x)>0\) 只能得到普通凸性。

\subsection{Bohr–Mollerup 定理}
\label{b3:subsec:210104}

递推固定了整数点的对数差值，而凸性控制整数之间的割线斜率。把一个未知的 \(F(x)\) 向右递推到 \(F(n+x)\)，再利用相邻整数之间的斜率夹逼，就能让未知量落入相对宽度趋于零的区间。这给出\term[bohr-mollerup]{Bohr--Mollerup 定理}[Bohr--Mollerup theorem]。

\Needspace{6\baselineskip}
\begin{theorem}[Bohr--Mollerup]\label{b3:thm:21-bohr-mollerup}
设 \(F:(0,\infty)\to(0,\infty)\) 满足
\[
 F(1)=1,\quad F(x+1)=xF(x)\quad(x>0),
\]
并且 \(x\mapsto\log F(x)\) 凸。则对每个 \(x>0\)，\(F(x)=\Gamma(x)\)。
\end{theorem}
\begin{proof}
令 \(h=\log F\)。递推与归一化先给 \(F(n)=(n-1)!\)，因而对整数 \(n\ge2\)，
\[
 h(n)-h(n-1)=\log(n-1),\quad
 h(n+1)-h(n)=\log n.
\]
固定 \(0<x\le1\)。凸函数的割线斜率随区间向右移动而非减，所以
\[
 h(n)-h(n-1)
 \le\frac{h(n+x)-h(n)}x
 \le h(n+1)-h(n).
\]
右侧亦可由 \(n+x=(1-x)n+x(n+1)\) 的凸性直接得到。于是
\[
 x\log(n-1)\le h(n+x)-h(n)\le x\log n.
\]
指数化后，再用 \(F(n+x)=F(x)x(x+1)\cdots(x+n-1)\)，可得
\[
 A_n(x)\le F(x)\le B_n(x),
\]
其中
\[
 A_n(x)=\frac{(n-1)!(n-1)^x}{x(x+1)\cdots(x+n-1)},\quad
 B_n(x)=\frac{(n-1)!n^x}{x(x+1)\cdots(x+n-1)}.
\]
这两个端点都严格为正，且 \(B_n(x)/A_n(x)=(n/(n-1))^x\to1\)。

已经由积分证明的 \(\Gamma\) 同样满足归一化、递推和对数凸性，所以也落在同一个区间 \([A_n(x),B_n(x)]\) 中。两次夹逼相除，得到
\[
 \left(\frac{n-1}{n}\right)^x
 \le\frac{F(x)}{\Gamma(x)}
 \le\left(\frac n{n-1}\right)^x.
\]
令 \(n\to\infty\)，即得 \(F(x)=\Gamma(x)\) 于 \((0,1]\)。任意正实数可写成 \(x+m\)，其中 \(0<x\le1\)、\(m\ge0\) 为整数，再用有限次递推便扩展到全部正实轴。
\end{proof}
证明只要求 \(\log F\) 凸，不预先要求 \(F\) 可微。若另外给定一个在右半平面全纯的候选函数，且它在正实轴满足定理的条件，则先由本定理识别实轴上的值，再由恒等定理识别整个右半平面内的函数。实变量唯一性与复变量唯一性分别由这两步承担。

\ProseParagraphBreak

对数凸性也能给出精确到相对误差的平移比值。这里的自变量始终沿正实轴趋于无穷，平移参数则保持固定。

\Needspace{6\baselineskip}
\begin{proposition}[固定实平移的比值]\label{b3:prop:21-shift-ratio}
设 \(x>0\)、\(0\le a\le1\)。则
\[
 \left(\frac{x}{x+a}\right)^{1-a}
 \le\frac{\Gamma(x+a)}{x^a\Gamma(x)}\le1.
\]
特别地，对任意固定的 \(a,b\in\mathbb R\)，当 \(x\to+\infty\) 时有
\[
 \frac{\Gamma(x+a)}{\Gamma(x+b)}
 =x^{a-b}\bigl(1+O_{a,b}(x^{-1})\bigr).
\]
后一公式在所有参数 \(x+a,x+b\) 均为正的充分大半轴上使用。
\end{proposition}
\begin{proof}
\(a=0,1\) 时由归一化的比值和递推直接成立。对 \(0<a<1\)，将 \(x+a=(1-a)x+a(x+1)\) 代入对数凸性，得到
\[
 \Gamma(x+a)\le\Gamma(x)^{1-a}\Gamma(x+1)^a
 =x^a\Gamma(x).
\]
另一方面，\(x+1=a(x+a)+(1-a)(x+a+1)\)，故
\[
 \Gamma(x+1)
 \le\Gamma(x+a)^a\Gamma(x+a+1)^{1-a}
 =(x+a)^{1-a}\Gamma(x+a).
\]
使用 \(\Gamma(x+1)=x\Gamma(x)\) 并整理，便得下界。固定 \(a\in[0,1]\) 时，下界为 \((1+a/x)^{-(1-a)}=1+O_a(x^{-1})\)，所以
\[
 \frac{\Gamma(x+a)}{\Gamma(x)}=x^a\bigl(1+O_a(x^{-1})\bigr).
\]

任意固定实数 \(c\) 可写成 \(c=k+\theta\)，其中 \(k\in\mathbb Z\)、\(0\le\theta<1\)。若 \(k\ge0\)，有限次递推给
\[
 \Gamma(x+c)=\Gamma(x+\theta)
             \prod_{j=0}^{k-1}(x+\theta+j).
\]
若 \(k=-q<0\)，则
\[
 \Gamma(x+c)=\frac{\Gamma(x+\theta)}
 {\prod_{j=0}^{q-1}(x+c+j)}.
\]
对充分大的 \(x\)，这些因子均为正；每个因子提出一个 \(x\) 后，有限乘积或其倒数均为 \(1+O_c(x^{-1})\)。结合已经证明的 \(\theta\) 情形，得到 \(\Gamma(x+c)/\Gamma(x)=x^c(1+O_c(x^{-1}))\)。分别取 \(c=a,b\)，再相除，即得一般比值。
\end{proof}
误差常数依赖固定的平移参数；这条结论不能直接用于随 \(x\) 增大的平移量。它也没有给出复变量在任意方向趋于无穷的比值渐近。

\ProseParagraphBreak

半整数值与比值估计合起来，已经能够确定中心二项式系数的主项。对正整数 \(n\)，
\[
 \binom{2n}{n}
 =\frac{4^n}{\sqrt\pi}
   \frac{\Gamma(n+1/2)}{\Gamma(n+1)}
 =\frac{4^n}{\sqrt{\pi n}}\bigl(1+O(n^{-1})\bigr).
\]
这里常数 \(\sqrt\pi\) 来自 Gaussian 积分，比值的幂次来自对数凸性与递推。两种信息共同确定了这个离散系数的渐近尺度。

% !TeX root = ../../Book3.tex
\section{亚纯延拓}
\label{b3:sec:2102}

递推式既能向右计算，也能向左延拓。后一用法必须检查不同移动次数给出的值相同，并辨明分母的零点是真极点还是可去点。

\subsection{递推关系}
\label{b3:subsec:210201}

对任意正整数 \(N\)，在 \(\operatorname{Re}z>-N\) 定义
\[
\Gamma(z)=\frac{\Gamma(z+N)}{z(z+1)\cdots(z+N-1)}.
\]
当分母不为零时这是全纯函数，其他点至多为极点。用一次递推消去 \(z+N\) 可知，采用 \(N+1\) 与 \(N\) 的定义在重叠区域一致。因此这些局部表达式粘合成整个复平面上的亚纯函数，并继续满足递推式。以后 \(\Gamma\) 均指这一延拓。

\subsection{极点}
\label{b3:subsec:210202}

\begin{proposition}\label{b3:prop:21-poles}
Gamma 函数的极点恰为 \(0,-1,-2,\ldots\)，每个都是一阶极点。
\end{proposition}
\begin{proof}
若 \(z\) 不是非正整数，选足够大的 \(N\) 可使分子处于右半平面，且分母各因子非零，故该点全纯。在 \(z=-m\) 附近取 \(N=m+1\)。分子在该点等于 \(\Gamma(1)=1\)，分母只有因子 \(z+m\) 为零，其余因子都非零，所以是恰好一阶的极点。
\end{proof}
延拓并不赋予原积分在 \(\operatorname{Re}z\le0\) 上收敛的含义。例如 \(\Gamma(-1/2)=-2\sqrt\pi\)，而形式积分 \(\int_0^\infty e^{-t}t^{-3/2}\,\mathrm{d}t\) 发散且被积函数为正；两者不能混为一谈。

\subsection{留数}
\label{b3:subsec:210203}

上述局部表示直接给出
\[
\operatorname{Res}_{z=-m}\Gamma(z)=\frac{(-1)^m}{m!},\quad m=0,1,2,\ldots.
\]
事实上，提出 \(z+m\) 后，其他因子在 \(-m\) 的乘积为 \((-1)^m m!\)。在反演积分中移动围道时，这串留数将产生幂级数的系数，而极点位置给出相应的幂次。

\ProseParagraphBreak

复共轭关系 \(\Gamma(\overline z)=\overline{\Gamma(z)}\) 先由右半平面的积分得到，再经恒等定理延拓。负实轴相邻极点之间的符号因此可用递推明确计算：在 \((-m-1,-m)\) 上符号为 \((-1)^{m+1}\)。此符号变化来自越过一阶极点，不意味着存在实零点；事实上下一节将证明复零点也不存在。

\ProseParagraphBreak

极点附近还能得到有限部分，而不仅是留数。先由乘积的首项展开得到 \(\Gamma(z)=z^{-1}-\gamma+O(z)\)。在 \(z=-m+\varepsilon\) 使用递推，分母非零部分为
\[
 \prod_{j=0}^{m-1}(-m+\varepsilon+j)
 =(-1)^m m!\left(1-H_m\varepsilon+O(\varepsilon^2)\right).
\]
分子 \(\Gamma(1+\varepsilon)=1-\gamma\varepsilon+O(\varepsilon^2)\)。因此
\[
 \Gamma(-m+\varepsilon)=\frac{(-1)^m}{m!}
       \left(\frac1\varepsilon+H_m-\gamma+O(\varepsilon)\right),
 \qquad H_0=0.
\]
当把 Gamma 与另一个含零点的函数相乘时，留数决定可去后的值，有限部分则决定导数。只知道极点阶数不足以完成这些局部计算。
\ProseParagraphBreak
参数正则化提供一个与递推不同的延拓图像。在 \(-1<\operatorname{Re}z<0\) 中，积分
\[
 \int_0^1(e^{-t}-1)t^{z-1}\,\mathrm{d}t
 +\int_1^\infty e^{-t}t^{z-1}\,\mathrm{d}t+\frac1z
\]
绝对收敛并等于 Gamma 延拓。因为 \(e^{-t}-1=O(t)\)，第一项比原 Euler 积分多出一个可积幂次；最后的 \(1/z\) 恰把减掉的常数项补回。继续减去更多 Taylor 项，就能向更左的半平面推进。所谓解析延拓在这里表现为一项明确的奇部加一个普通收敛积分，并没有把发散积分本身赋予任意的数值。

% !TeX root = ../../Book3.tex
\section{乘积表示}
\label{b3:sec:2103}

上一节确定了全部极点，但还没有排除复零点。有限乘积的极限既给出零点信息，也解释 Euler 常数为什么出现在 Gamma 函数中。

\subsection{Weierstrass 乘积}
\label{b3:subsec:210301}

\begin{lemma}[Euler 极限]\label{b3:lem:21-euler-limit}
在 \(\operatorname{Re}z>0\) 上局部一致地有
\[
\Gamma(z)=\lim_{N\to\infty}\frac{N!N^z}{z(z+1)\cdots(z+N)}.
\]
\end{lemma}
\begin{proof}
连续分部积分给
\(\int_0^1 t^{z-1}(1-t)^N\,\mathrm{d}t=N!/[z(z+1)\cdots(z+N)]\)。乘以 \(N^z\)，换元 \(u=Nt\)，所得积分为
\(\int_0^N u^{z-1}(1-u/N)^N\,\mathrm{d}u\)。在 \(0<u<N\) 上有 \((1-u/N)^N\le e^{-u}\)，而它逐点趋于 \(e^{-u}\)。对紧集 \(a\le\operatorname{Re}z\le b\)，用 \((u^{a-1}+u^{b-1})e^{-u}\) 控制。将积分分成固定紧区间与共同小尾部，即得对 \(z\) 的局部一致收敛。
\end{proof}
这个证明同时保留复参数与一致性，避免只在实数处建立公式后未经说明就逐项微分。有限分部积分使用的端点条件来自 \(\operatorname{Re}z>0\)。

\subsection{Euler 常数}
\label{b3:subsec:210302}

记 \(H_N=\sum_{n=1}^N1/n\)。序列 \(H_N-\log N\) 单调递减，因为
\(1/(N+1)\le\log(1+1/N)\)，并由积分比较有下界零，故极限存在，记为\term[euler-changshu]{Euler 常数}[Euler constant] \(\gamma\)。把有限乘积写为
\[
\frac{z(z+1)\cdots(z+N)}{N!N^z}
=z\exp\bigl(z(H_N-\log N)\bigr)
\prod_{n=1}^N(1+z/n)e^{-z/n}.
\]
紧集上，尾部对数 \(\log(1+z/n)-z/n=O_K(n^{-2})\) 一致可和；前面有限个因子无需取对数。于是右侧在全平面局部一致收敛为整函数。

\subsection{\texorpdfstring{\(1/\Gamma\)}{1/\AnalysisPdfGamma} 作为整函数}
\label{b3:subsec:210303}

\begin{theorem}[Gamma 函数的 Weierstrass 乘积]\label{b3:thm:21-product}
Gamma 函数没有零点，并且
\[
\frac1{\Gamma(z)}=ze^{\gamma z}\prod_{n=1}^\infty(1+z/n)e^{-z/n}.
\]
右端整函数的零点恰为非正整数，且均为单零点。
\end{theorem}
\begin{proof}
设右端为 \(P(z)\)。Euler 极限与相应的有限倒数乘积在右半平面都局部一致收敛，它们的乘积恒为一，取极限得 \(P\Gamma=1\)，不必事先假设 \(\Gamma\ne0\)。尾部指数的和在避开有限个头部零点的邻域全纯，因此尾乘积不为零；全部零点正是所列的线性因子零点，重数为一。亚纯恒等定理把 \(P\Gamma=1\) 延拓到全平面，给出结论。
\end{proof}
对数导数在避开极点的紧集可逐项求导。若记 \(\psi_0(z)=\Gamma'(z)/\Gamma(z)\)，则
\[
\psi_0(z)=-\gamma-\frac1z+\sum_{n=1}^\infty\left(\frac1n-\frac1{n+z}\right).
\]
这里用下标避免与后文的素数计数函数 \(\psi(x)\) 混淆。代入 \(z=1\) 得 \(\Gamma'(1)=-\gamma\)，再用递推得 \(\psi_0(n)=H_{n-1}-\gamma\)。整数点的导数遂由同一个调和和描述。

\ProseParagraphBreak

在零点附近展开对数乘积，可以得到 Gamma 的局部 Taylor 系数。对 \(|z|<1\)，\(\Gamma(1+z)\) 无零点且在0点等于一，取在那里为零的对数。利用 \(\Gamma(1+z)=z\Gamma(z)\)，有
\[
 \log\Gamma(1+z)=-\gamma z+
   \sum_{k=2}^\infty\frac{(-1)^k}{k}
       \left(\sum_{n=1}^\infty\frac1{n^k}\right)z^k.
\]
证明中的换序来自 \(|z|\le r<1\) 时
\(\sum_n\sum_{k\ge2}|z/n|^k/k\le C_r\sum_n n^{-2}\)。内层数列以后记为 \(\zeta(k)\)，当前只需把它看成收敛的正项数列。常数项与一次项已经分别由 \(\Gamma(1)=1\) 和 Euler 常数决定，二次项才开始包含这些幂和。
\ProseParagraphBreak
特别地，\(\Gamma''(1)=\gamma^2+\sum n^{-2}\)。与积分微分公式比较，就得到
\[
 \int_0^\infty e^{-t}\log t\,\mathrm{d}t=-\gamma,
 \qquad
 \int_0^\infty e^{-t}(\log t)^2\,\mathrm{d}t
 =\gamma^2+\sum_{n\ge1}n^{-2}.
\]
第二式减去第一式的平方，正好是前面概率解释中 \(x=1\) 的方差。乘积的系数、参数积分的矩和凸性的二阶导数在同一恒等式中相遇。
\ProseParagraphBreak
对正实 \(x\)，对数导数还给出
\[
 \psi_0'(x)=\sum_{n=0}^\infty\frac1{(n+x)^2},\qquad
 \psi_0^{(j)}(x)=(-1)^{j+1}j!\sum_{n=0}^\infty(n+x)^{-j-1}\quad(j\ge1).
\]
这些级数在任何 \(x\ge a>0\) 的紧区间上一致可微，因而 \(\psi_0'\) 正，后续导数交替符号。对数凸性的严格性由此再次可见。用递减函数的积分比较还有
\(1/x\le\psi_0'(x)\le1/x+1/x^2\)，它量化了大参数时曲率逐渐减小的速度。
\ProseParagraphBreak
有限乘积中的指数修正也有必要性。若直接写 \(\prod_{n\ge1}(1+z/n)\)，对正实 \(z\) 其对数含发散的 \(z\sum1/n\)，因而不会趋于有限非零值。乘 \(e^{-z/n}\) 消去这一不可和的一次项，剩下二次误差才可和；前面的 \(e^{\gamma z}\) 则记录有限乘积正规化的极限。把该因子删去会保留零点，却破坏与 Euler 积分的对应及递推常数。

% !TeX root = ../../Book3.tex
\section{经典公式}
\label{b3:sec:2104}

乘积能把相反参数联系起来；凸性则能把相差有理数的参数联系起来。增长估计另需考察积分质量集中的位置，不能只由函数方程猜测。

\subsection{反射公式}
\label{b3:subsec:210401}

\begin{theorem}[Euler 反射公式]\label{b3:thm:21-reflection}
对 \(z\notin\mathbb Z\)，有 \(\Gamma(z)\Gamma(1-z)=\pi/\sin\pi z\)。
\end{theorem}
\begin{proof}
对 \(1/\Gamma(z)\) 与 \(1/\Gamma(-z)\) 的正规乘积相乘，指数因子相消，得
\[
\frac1{\Gamma(z)\Gamma(-z)}=-z^2\prod_{n\ge1}(1-z^2/n^2)
=-\frac{z\sin\pi z}{\pi}.
\]
最后一步使用\secref{b3:sec:2005}已证的正弦乘积。由 \(\Gamma(1-z)=-z\Gamma(-z)\) 即得公式。
\end{proof}
这称为\term[euler-fanshe-gongshi]{Euler 反射公式}[Euler reflection formula]。取 \(z=1/2\)，积分的正性确定 \(\Gamma(1/2)=\sqrt\pi\)，排除平方根符号的不确定性。

\subsection{倍增公式}
\label{b3:subsec:210402}

\begin{theorem}[Legendre 倍增公式]\label{b3:thm:21-duplication}
作为亚纯恒等式，有
\[
\Gamma(z)\Gamma(z+1/2)=2^{1-2z}\sqrt\pi\,\Gamma(2z).
\]
\end{theorem}
\begin{proof}
对 \(x>0\) 定义 \(F(x)=2^{x-1}\Gamma(x/2)\Gamma((x+1)/2)/\sqrt\pi\)。直接递推得 \(F(x+1)=xF(x)\)，且 \(F(1)=1\)。函数 \(\log F\) 为两个凸函数与线性函数之和，因此凸。Bohr--Mollerup 定理给 \(F(x)=\Gamma(x)\)；令 \(x=2z\)，先在正实轴成立，再由亚纯恒等定理扩展。
\end{proof}
该式称为\term[legendre-beizeng-gongshi]{Legendre 倍增公式}[Legendre duplication formula]。例如它把半整数阶乘写成
\(\Gamma(n+1/2)=(2n)!\sqrt\pi/(4^nn!)\)，便于精确计算高斯积分的偶数矩。

\subsection{乘法公式}
\label{b3:subsec:210403}

对整数 \(m\ge1\)，\term[gauss-chengfa-gongshi]{Gauss 乘法公式}[Gauss multiplication formula]为
\begin{equation}\label{b3:eq:21-multiplication}
\prod_{k=0}^{m-1}\Gamma(z+k/m)
=(2\pi)^{(m-1)/2}m^{1/2-mz}\Gamma(mz).
\end{equation}
先用与倍增公式相同的方法，令
\(F(x)=C m^x\prod_{k=0}^{m-1}\Gamma((x+k)/m)\)。相邻乘积约去后得到 \(F(x+1)=xF(x)\)，且对数凸，选 \(C\) 使 \(F(1)=1\) 即知它等于 \(\Gamma(x)\)。求常数时，将反射公式对 \(k=1,\ldots,m-1\) 相乘，使用
\[
\prod_{k=1}^{m-1}\sin(\pi k/m)=m/2^{m-1}.
\]
这个有限恒等式来自 \((X^m-1)/(X-1)\) 在 \(X=1\) 的值：各非平凡单位根与一的距离为 \(2\sin(\pi k/m)\)，取模即可。故
\(\prod_{k=1}^{m-1}\Gamma(k/m)=(2\pi)^{(m-1)/2}/\sqrt m\)，平方根取正值。代回 \(C\)，得到所示常数。

\subsection{Stirling 公式}
\label{b3:subsec:210404}

\begin{theorem}[Stirling]\label{b3:thm:21-stirling}
当实数 \(x\to+\infty\) 时，
\[
\Gamma(x)=\sqrt{2\pi}\,x^{x-1/2}e^{-x}(1+o(1)).
\]
\end{theorem}
\begin{proof}
在 \(\Gamma(x+1)\) 的积分中令 \(t=x+u\sqrt x\)，提出 \(x^xe^{-x}\sqrt x\)。剩余积分为
\[
I_x=\int_{-\sqrt x}^{\infty}\exp\left[x\left(\log(1+u/\sqrt x)-u/\sqrt x\right)\right]\,\mathrm{d}u.
\]
固定 \(u\) 时指数趋于 \(-u^2/2\)。对 \(|v|\le1/2\)，Taylor 定理及二阶导数给 \(\log(1+v)-v\le-cv^2\)，故中部由 \(e^{-cu^2}\) 控制。设 \(h(v)=v-\log(1+v)\)。对 \(v\ge1/2\)，有 \(h'(v)=v/(1+v)\ge1/3\)，于是右尾积分不超过
\(3x^{-1/2}e^{-xh(1/2)}\)。对 \(-1<v\le-1/2\)，\(h'(v)\le-1\)，左尾同样不超过 \(x^{-1/2}e^{-xh(-1/2)}\)。先用支配收敛处理扩大的中部，再用这两个指数小界删去尾部，得 \(I_x\to\int_{\mathbb R}e^{-u^2/2}\,\mathrm{d}u=\sqrt{2\pi}\)。最后除以 \(x\) 即得结论。
\end{proof}
这给出\term[stirling-gongshi]{Stirling 公式}[Stirling formula]的正实轴版本。证明解释了平方根因子：积分主要来自峰值 \(t=x\) 周围宽度为 \(\sqrt x\) 的区间。它不声称在穿过负实轴的复扇形上一致成立；复参数渐近展开还必须指定对数分支及与负实轴的角距离。NIST 的《Digital Library of Mathematical Functions》\cite[Sections 5.5 and 5.11]{NationalInstituteofStandardsNIST}分别列出上述函数关系和复参数渐近版本，使用时应保留其区域条件。


\ProseParagraphBreak

反射公式能给虚方向上的精确模值，而不需要复 Stirling 公式。对实 \(t\ne0\)，利用共轭关系与递推，有
\[
 |\Gamma(\mathrm{i}t)|^2=\frac{\pi}{t\sinh\pi t},\qquad
 |\Gamma(1+\mathrm{i}t)|^2=\frac{\pi t}{\sinh\pi t}.
\]
确实，反射式在 \(z=\mathrm{i}t\) 给分母 \(\sin\pi\mathrm{i}t=\mathrm{i}\sinh\pi t\)，而 \(\Gamma(1-\mathrm{i}t)=-\mathrm{i}t\Gamma(-\mathrm{i}t)\)；移项后得到第一式，第二式再乘 \(t^2\)。当 \(t\to0\)，第二式趋于一，恰与 \(\Gamma(1)=1\) 相符；当 \(|t|\to\infty\)，它给 \(|\Gamma(1+\mathrm{i}t)|\sim\sqrt{2\pi|t|}e^{-\pi|t|/2}\)。这是一条竖线上完整的衰减结论，不被误写为全扇形的统一渐近。
\ProseParagraphBreak
在半整数竖线上，同理可直接将 \(z=1/2+\mathrm{i}t\) 代入反射公式，得到
\[
 |\Gamma(1/2+\mathrm{i}t)|^2=\frac\pi{\cosh\pi t}.
\]
再对实部作整数递推，便能得到所有整数或半整数竖线上的精确乘积表达式。这些特殊线已经足以检验很多 Mellin 反演中的绝对收敛。若需要任意实部的统一估计，则必须另作插值或扇形分析，当前几条特殊线的公式并不自动包含这一推广。

% !TeX root = ../../Book3.tex
\section{Mellin 变换}
\label{b3:sec:2105}

平移问题自然使用 Fourier 变换，伸缩问题则更适合先取对数坐标。正半轴上的乘法 \(x\mapsto ax\) 在 \(u=\log x\) 中变成平移 \(u\mapsto u+\log a\)；测度 \(\mathrm dx/x\) 也恰好变成 \(\mathrm du\)。Mellin 变换由这个换元产生，它的复参数实部选择权重，虚部记录对数变量的振荡。两者承担不同作用，反演时必须同时保留。

\subsection{定义}
\label{b3:subsec:210501}

\Needspace{6\baselineskip}
\begin{definition}[Mellin 变换与绝对收敛带]\label{b3:def:21-mellin}
设 \(f:(0,\infty)\to\mathbb C\) 可测。对每个使积分绝对收敛的 \(s\in\mathbb C\)，定义其\term[mellin-bianhuan]{Mellin 变换}[Mellin transform]为
\[
 \mathcal Mf(s)\coloneqq\int_0^\infty f(x)x^{s-1}\,\mathrm dx,
 \quad x^{s-1}=\exp((s-1)\log x),
\]
其中 \(\log x\) 为正实轴上的实对数。对实数 \(c\)，记
\[
 \|f\|_{1,c}\coloneqq\int_0^\infty|f(x)|x^{c-1}\,\mathrm dx,
 \quad L^1_c=\{\,f:\|f\|_{1,c}<\infty\,\},
\]
函数按几乎处处相等识别。若集合 \(I_f=\{\,c\in\mathbb R:\|f\|_{1,c}<\infty\,\}\) 的内部为非空区间 \((a,b)\)，则称 \(a<\operatorname{Re}s<b\) 为 \(f\) 的绝对收敛带；允许 \(a=-\infty\) 或 \(b=+\infty\)。
\end{definition}
\symidx{\(\mathcal Mf(s)\)：Mellin 变换}
\symidx{\(\lVert f\rVert_{1,c}\)：Mellin 权重的绝对可积范数}
绝对值满足 \(|x^{s-1}|=x^{\operatorname{Re}s-1}\)，所以是否绝对收敛只取决于实部。若在一个实部收敛，在同一条竖线的每一点都收敛；这与函数在一个开区域全纯是两件不同的事。

\Needspace{6\baselineskip}
\begin{proposition}[收敛区间与带内全纯性]\label{b3:prop:21-mellin-strip}
设 \(f:(0,\infty)\to\mathbb C\) 可测。其绝对可积权重集合 \(I_f\) 是一个区间，允许为空、单点及具有任意端点开闭性。若实数 \(a<b\) 满足 \(f\in L^1_a\cap L^1_b\)，则 \(\mathcal Mf\) 在闭带 \(a\le\operatorname{Re}s\le b\) 有界且连续，在开带内全纯，并且对每个非负整数 \(k\)，
\[
 (\mathcal Mf)^{(k)}(s)
 =\int_0^\infty f(x)x^{s-1}(\log x)^k\,\mathrm dx,
 \quad a<\operatorname{Re}s<b.
\]
这些导数积分在开带的紧集上绝对一致收敛。
\end{proposition}
\begin{proof}
若 \(a\le\sigma\le b\)，则在 \(0<x\le1\) 上 \(x^{\sigma-1}\le x^{a-1}\)，在 \(x\ge1\) 上 \(x^{\sigma-1}\le x^{b-1}\)。因此
\[
 |f(x)x^{s-1}|
 \le |f(x)|\bigl(x^{a-1}\cdot\mathbf1_{(0,1]}(x)
                  +x^{b-1}\cdot\mathbf1_{[1,\infty)}(x)\bigr).
\]
右边可积且与 \(s\) 无关。这既证明中间权重均可积，也由支配收敛给出闭带连续性及上界 \(|\mathcal Mf(s)|\le\|f\|_{1,a}+\|f\|_{1,b}\)。所以 \(I_f\) 具有区间性质。

在开带内固定一条较小闭带 \(a+\delta\le\operatorname{Re}s\le b-\delta\)，其中 \(0<2\delta<b-a\)。对任意固定 \(k\)，函数 \(x^\delta|\log x|^k\) 在 \((0,1]\) 有界，\(x^{-\delta}(\log x)^k\) 在 \([1,\infty)\) 有界。因此加入 \((\log x)^k\) 后，仍由上述两个端点权重的常数倍控制。截断到 \([\epsilon,R]\) 的积分是 \(s\) 的整函数，其各阶导数恰由插入 \((\log x)^k\) 得到。可积控制使这些积分及导数在较小闭带上一致收敛，故可取全纯极限并逐次微分。
\end{proof}

端点不应从带内结论猜测。给定实数 \(c_0\)，函数
\[
 f(x)=\frac{x^{-c_0}}{(1+|\log x|)^2}
\]
满足 \(I_f=\{c_0\}\)：令 \(u=\log x\)，其加权绝对积分为
\[
 \int_{\mathbb R}\frac{e^{(c-c_0)u}}{(1+|u|)^2}\,\mathrm du,
\]
仅在 \(c=c_0\) 时有限。它在一条竖线上有变换，却没有非空的绝对收敛开带。相反，若 \(f\) 可积且支集包含于某个 \([\delta,R]\subset(0,\infty)\)，则所有实权重都可积，变换为整函数。端点行为决定参数域，不能只看被积函数中出现的复幂。

\subsection{缩放}
\label{b3:subsec:210502}

Mellin 变换的换元规则应连同权重一起记忆。正半轴上的伸缩不改变复参数，乘上一个幂则移动参数所在的直线；取倒数坐标会交换原点与无穷远。

\Needspace{6\baselineskip}
\begin{proposition}[缩放、幂乘与倒数换元]\label{b3:prop:21-mellin-rules}
设 \(f:(0,\infty)\to\mathbb C\) 可测，\(a>0\)、\(\lambda\in\mathbb C\)。凡右侧积分绝对收敛，下列左侧积分也绝对收敛，且有
\[
 \begin{aligned}
 \mathcal M[f(a\,\cdot)](s)&=a^{-s}\mathcal Mf(s),\\
 \mathcal M[x^\lambda f(x)](s)&=\mathcal Mf(s+\lambda),\\
 \mathcal M[f(1/x)](s)&=\mathcal Mf(-s).
 \end{aligned}
\]
全部正实数的复幂均用实对数定义。另设实数 \(c\) 和函数 \(f\in C^1((0,\infty))\) 满足
\[
 f,\;xf'\in L^1_c,
 \quad \lim_{x\downarrow0}x^cf(x)=\lim_{x\to\infty}x^cf(x)=0.
\]
则在 \(\operatorname{Re}s=c\) 上有
\begin{equation}\label{b3:eq:21-mellin-euler-operator}
 \mathcal M[xf'](s)=-s\mathcal Mf(s).
\end{equation}
\end{proposition}
\begin{proof}
第一式令 \(y=ax\)，有 \(x^{s-1}\,\mathrm dx=a^{-s}y^{s-1}\,\mathrm dy\)；对绝对值作同一换元也验证收敛性。第二式直接合并指数，绝对权重由 \(\operatorname{Re}s\) 移到 \(\operatorname{Re}s+\operatorname{Re}\lambda\)。第三式令 \(y=1/x\)，反转端点后得到 \(x^{s-1}\,\mathrm dx=-y^{-s-1}\,\mathrm dy\)，正是参数 \(-s\) 的积分。

最后在 \([\epsilon,R]\) 分部积分：
\[
 \int_\epsilon^R f'(x)x^s\,\mathrm dx
 =[x^sf(x)]_\epsilon^R-s\int_\epsilon^R f(x)x^{s-1}\,\mathrm dx.
\]
因 \(|x^{s-c}|=1\)，端点假设使边界项趋零；两项绝对可积假设允许令 \(\epsilon\downarrow0\)、\(R\to\infty\)，得到\eqref{b3:eq:21-mellin-euler-operator}。
\end{proof}

算子 \(x\,\mathrm d/\mathrm dx\) 是伸缩在参数原点的导数，因为 \(\left.\mathrm d f(e^u x)/\mathrm du\right|_{u=0}=xf'(x)\)。它对应乘子 \(-s\)。重复使用此式时，每一级函数都必须满足相应加权可积性和边界条件；形式地把微分方程换成代数方程，并不会自动消去端点贡献。

\subsection{与 Gamma 函数的关系}
\label{b3:subsec:210503}

Euler 积分就是最基本的 Mellin 变换对：
\[
 \mathcal M[e^{-x}](s)=\Gamma(s),\quad \operatorname{Re}s>0.
\]
原点处 \(e^{-x}\) 趋于一，要求正实部；无穷远的指数衰减则容许任意固定幂权重。再令 \(t=ax^p\)，得到
\begin{equation}\label{b3:eq:21-mellin-exponential}
 \mathcal M[e^{-ax^p}](s)=\frac1p a^{-s/p}\Gamma(s/p),
 \quad a,p>0,\quad \operatorname{Re}s>0.
\end{equation}
其中 \(1/p\) 与 \(a^{-s/p}\) 都由微分元和幂权重共同产生，不能只改变 Gamma 的参数而漏掉常数。

\ProseParagraphBreak

一个同时检验两端的例子是 \(f(x)=(1+x)^{-a}\)，其中 \(a>0\) 为实数。原点要求 \(\operatorname{Re}s>0\)，无穷远要求 \(\operatorname{Re}s<a\)。在这条带内作 \(u=x/(1+x)\)，由\cref{b3:prop:21-beta-integral}，
\begin{equation}\label{b3:eq:21-mellin-beta}
 \int_0^\infty\frac{x^{s-1}}{(1+x)^a}\,\mathrm dx
 =B(s,a-s)=\frac{\Gamma(s)\Gamma(a-s)}{\Gamma(a)}.
\end{equation}
取 \(a=1\)，反射公式将右边化成 \(\pi/\sin\pi s\)。它延拓为亚纯函数，但原积分依然只在 \(0<\operatorname{Re}s<1\) 绝对收敛；例如在实数 \(s>1\) 时，原积分于无穷远发散，不等于延拓的有限值。

\ProseParagraphBreak

一般的端点扣除也可在此统一说明。局部幂项 \(x^\alpha\) 在 \((0,1)\) 上给 \(1/(s+\alpha)\)，它的对数幂则给更高阶极点。为把这项计算变成延拓结论，必须同时控制扣除后的余项与未改变的无穷远积分。

\Needspace{6\baselineskip}
\begin{proposition}[端点幂与对数项的 Mellin 延拓]\label{b3:prop:21-mellin-endpoint}
设 \(f:(0,\infty)\to\mathbb C\) 可测且局部可积。设 \(A,b\in\mathbb R\)，有限个复数 \(\alpha_j\) 满足 \(\operatorname{Re}\alpha_j<A\)，并给定非负整数 \(q_j\) 及复数 \(c_{j,\ell}\)。假设当 \(x\downarrow0\) 时，
\[
 f(x)=\sum_{j=1}^J\sum_{\ell=0}^{q_j}
            c_{j,\ell}x^{\alpha_j}(\log x)^\ell+O(x^A),
\]
且对每个实数 \(\sigma<b\)，有
\[
 \int_1^\infty|f(x)|x^{\sigma-1}\,\mathrm dx<\infty.
\]
当 \(J\ge1\) 时，假设 \(\max_j(-\operatorname{Re}\alpha_j)<b\)；当 \(J=0\) 时，假设 \(-A<b\)。则原变换在 \(J\ge1\) 时于 \(\max_j(-\operatorname{Re}\alpha_j)<\operatorname{Re}s<b\) 绝对收敛，并在 \(-A<\operatorname{Re}s<b\) 有亚纯延拓
\begin{equation}\label{b3:eq:21-mellin-endpoint-extension}
 \mathcal Mf(s)=H(s)+\sum_{j=1}^J\sum_{\ell=0}^{q_j}
       \frac{(-1)^\ell\ell!c_{j,\ell}}{(s+\alpha_j)^{\ell+1}},
\end{equation}
其中 \(H\) 在这条较大开带全纯。没有主项时，原变换已经在该较大开带全纯。列出的极点只表示可能的极点；相同位置的系数须先合并，可能发生抵消。
\end{proposition}
\begin{proof}
在 \((0,1)\) 上令
\[
 r(x)=f(x)-\sum_{j=1}^J\sum_{\ell=0}^{q_j}
                   c_{j,\ell}x^{\alpha_j}(\log x)^\ell.
\]
原点附近有 \(|r(x)|\le Cx^A\)，而在远离零的紧区间，\(r\) 仍可积。所以
\[
 H(s)=\int_0^1r(x)x^{s-1}\,\mathrm dx
      +\int_1^\infty f(x)x^{s-1}\,\mathrm dx
\]
在 \(-A<\operatorname{Re}s<b\) 绝对收敛。每条内闭带都能在左右各留少量幂次余量，从而吸收任意固定次 \(|\log x|\)；与\cref{b3:prop:21-mellin-strip}相同的支配论证给出全纯性。

在原来各主项都可积的开带内，对 \(\operatorname{Re}(s+\alpha_j)>0\)，先积分 \(x^{s+\alpha_j-1}\)，再在参数上微分 \(\ell\) 次，得到
\[
 \int_0^1x^{s+\alpha_j-1}(\log x)^\ell\,\mathrm dx
 =\frac{(-1)^\ell\ell!}{(s+\alpha_j)^{\ell+1}}.
\]
参数微分仍由少量幂次余量控制。将有限个主项补回，便在非空的原收敛带得到\eqref{b3:eq:21-mellin-endpoint-extension}，其右端遂给出唯一的亚纯延拓。
\end{proof}

对 \(e^{-x}\) 扣除 \(N\) 个 Taylor 项，\(N\ge1\)，就是
\[
 \Gamma(s)=H_N(s)+\sum_{k=0}^{N-1}\frac{(-1)^k}{k!(s+k)},
 \quad \operatorname{Re}s>-N,
\]
其中 \(H_N\) 全纯。这从 Mellin 角度再次解释了 Gamma 的极点与留数：幂次 \(k\) 产生位置 \(-k\)，Taylor 系数 \((-1)^k/k!\) 变成留数。该例在无穷远有任意幂权重的可积性，因此可将命题分别用于任意大的有限 \(b\)。

\ProseParagraphBreak

这里只证明了“有受控端点展开，因此有这些极点”的方向。反向从极点读取函数渐近时，还须有可用的反演和围道移动估计。一个亚纯表达式本身没有记载积分在高处的增长，不能仅列出留数就宣布已恢复原函数。

\subsection{Mellin 反演}
\label{b3:subsec:210504}

固定一条绝对收敛竖线 \(\operatorname{Re}s=c\)。令
\[
 g_c(u)=e^{cu}f(e^u),\quad u\in\mathbb R.
\]
对数换元给出精确的范数关系与变换关系：
\begin{equation}\label{b3:eq:21-mellin-fourier}
 \begin{aligned}
 \|g_c\|_{L^1(\mathbb R)}&=\|f\|_{1,c},\\
 \mathcal Mf(c+\mathrm it)&=\int_{\mathbb R}g_c(u)e^{\mathrm itu}\,\mathrm du
                         =\widehat g_c(-t/(2\pi)).
 \end{aligned}
\end{equation}
这里 \(\widehat g(\xi)=\int g(u)e^{-2\pi\mathrm iu\xi}\,\mathrm du\) 是本书的 Fourier 归一化。正号 \(e^{\mathrm itu}\) 和因子 \(2\pi\) 决定反演中的符号与常数。

\Needspace{6\baselineskip}
\begin{theorem}[Mellin 反演]\label{b3:thm:21-mellin-inversion}
设 \(f:(0,\infty)\to\mathbb C\) 可测，\(c\in\mathbb R\)，且
\[
 \int_0^\infty|f(x)|x^{c-1}\,\mathrm dx<\infty,
 \quad \int_{\mathbb R}|\mathcal Mf(c+\mathrm it)|\,\mathrm dt<\infty,
\]
其中 \(\mathcal Mf(s)=\int_0^\infty f(x)e^{(s-1)\log x}\,\mathrm dx\)。则
\begin{equation}\label{b3:eq:21-mellin-inversion}
 f(x)=\frac1{2\pi}\int_{\mathbb R}
               \mathcal Mf(c+\mathrm it)x^{-c-\mathrm it}\,\mathrm dt
\end{equation}
对几乎处处的 \(x>0\) 成立，右端定义 \((0,\infty)\) 上的连续函数。若 \(f\) 连续，等式对每个 \(x>0\) 成立。
\end{theorem}
\begin{proof}
令 \(g(u)=e^{cu}f(e^u)\) 以及 \(F(t)=\mathcal Mf(c+\mathrm it)\)。由\eqref{b3:eq:21-mellin-fourier}，\(g\in L^1(\mathbb R)\)，且 \(F(t)=\int g(v)e^{\mathrm itv}\,\mathrm dv\)。对 \(\epsilon>0\)，在反演积分中暂时插入 Gaussian 因子，定义
\[
 J_\epsilon(u)=\frac1{2\pi}\int_{\mathbb R}
                   F(t)e^{-\epsilon t^2}e^{-\mathrm itu}\,\mathrm dt.
\]
代入 \(F\) 后，绝对二重积分由
\[
 \|g\|_1\int_{\mathbb R}e^{-\epsilon t^2}\,\mathrm dt<\infty
\]
控制，故 Fubini 定理适用。

为核对内积分的常数，记 \(A_\epsilon(w)=\int e^{-\epsilon t^2}e^{-\mathrm itw}\,\mathrm dt\)。参数微分及一次分部积分给
\[
 A_\epsilon'(w)=-\frac{w}{2\epsilon}A_\epsilon(w),
 \quad A_\epsilon(0)=\sqrt{\pi/\epsilon}.
\]
分部积分的端点项因 Gaussian 衰减而消失；参数微分由 \(|t|e^{-\epsilon t^2}\) 控制。解这个实参数微分方程，得到
\[
 \frac1{2\pi}A_\epsilon(w)
 =K_\epsilon(w)\coloneqq\frac{e^{-w^2/(4\epsilon)}}{\sqrt{4\pi\epsilon}},
 \quad J_\epsilon=K_\epsilon*g.
\]
核 \(K_\epsilon\) 非负、积分为一，是固定的径向非增可积核在尺度 \(\sqrt\epsilon\) 下的缩放。第一卷命题14.4.6说明：这样的核作用于全空间 \(L^1\) 函数，在每个 Lebesgue 点收敛到函数值。这里 \(g\in L^1(\mathbb R)\)，核的全部条件也已核对，所以
\[
 J_\epsilon(u)\longrightarrow g(u)
\]
在 \(g\) 的几乎所有点成立。采用无限支集的 Gaussian 核时，全空间可积性仍不可省去；只知道核的质量向原点集中，不足以直接套用这一逐点结论。

另一方面，假设 \(F\in L^1\) 允许以 \(|F(t)|\) 作支配函数，得到对每个 \(u\) 都成立的极限
\[
 J_\epsilon(u)\longrightarrow
 J(u)=\frac1{2\pi}\int_{\mathbb R}F(t)e^{-\mathrm itu}\,\mathrm dt.
\]
同一个支配收敛论证说明 \(J\) 连续。于是 \(g=J\) 几乎处处；令 \(u=\log x\)，再乘 \(x^{-c}\)，便得\eqref{b3:eq:21-mellin-inversion}。这一换元保持零测集结论，因为指数映射在每个有界区间 Lipschitz，而实轴可被可数个有界区间覆盖。若 \(f\) 连续，两个连续函数几乎处处相等就必处处相等。
\end{proof}

反演也可写为
\[
 f(x)=\frac1{2\pi\mathrm i}
       \int_{c-\mathrm i\infty}^{c+\mathrm i\infty}\mathcal Mf(s)x^{-s}\,\mathrm ds.
\]
直线由下到上，\(\mathrm ds=\mathrm i\,\mathrm dt\)，所以两种记法的常数一致。它是绝对收敛的竖线积分；尚未允许改变直线的位置。

\Needspace{6\baselineskip}
\begin{corollary}[一条竖线上的唯一性]\label{b3:cor:21-mellin-uniqueness}
设 \(c\in\mathbb R\)，可测函数 \(f,h\in L^1_c\)。若
\[
 \mathcal Mf(c+\mathrm it)=\mathcal Mh(c+\mathrm it),\quad t\in\mathbb R,
\]
则 \(f=h\) 几乎处处；若二者连续，则处处相等。
\end{corollary}
\begin{proof}
差 \(f-h\) 属于 \(L^1_c\)，其竖线变换恒为零，因而自动属于 \(L^1(\mathbb R,\mathrm dt)\)。对差应用反演定理，就得到零函数。这里不需要分别假设两个原变换在竖线上可积。
\end{proof}

竖线上的绝对可积性不是加权可积性的自动后果。例如 \(f=\mathbf1_{(0,1)}\) 在每个 \(c>0\) 属于 \(L^1_c\)，且 \(\mathcal Mf(s)=1/s\)；但 \(1/(c+\mathrm it)\) 的模只按 \(|t|^{-1}\) 衰减，不能直接使用本节的绝对反演定理。对称截断或其他正则化仍可能给出反演，需要另行陈述其收敛方式。

\Needspace{6\baselineskip}
\begin{proposition}[绝对反演的一个充分条件]\label{b3:prop:21-mellin-inversion-test}
设 \(c\in\mathbb R\)，\(f:(0,\infty)\to\mathbb C\)，且 \(g(u)=e^{cu}f(e^u)\) 属于 \(C^2(\mathbb R)\)。若 \(g,g',g''\in L^1(\mathbb R)\)，且 \(g(u),g'(u)\to0\) 当 \(u\to\pm\infty\)，则
\[
 |\mathcal Mf(c+\mathrm it)|
 \le\min\{\|g\|_1,\,|t|^{-2}\|g''\|_1\}\quad(t\ne0).
\]
因此 \(f\) 满足 Mellin 反演定理的全部条件，且反演处处成立。
\end{proposition}
\begin{proof}
范数关系给出 \(f\in L^1_c\)。在 \(\int g(u)e^{\mathrm itu}\,\mathrm du\) 中分部积分两次，端点项均因假设消失，得到模界 \(|t|^{-2}\|g''\|_1\)；不分部积分则得到 \(\|g\|_1\)。前一个界控制 \(|t|\ge1\)，后一个控制 \(|t|\le1\)，二者合并可积。\(g\) 连续也保证 \(f(x)=x^{-c}g(\log x)\) 连续，故反演处处成立。
\end{proof}

\subsection{对数坐标、乘法卷积与围道移动}
\label{b3:subsec:210505}

每条竖线对应不同的加权函数 \(g_c\)。改变 \(c\) 时，不只是给同一 Fourier 变换换了记号，而是改变了输入在两端的权重。对于伸缩结构，最自然的卷积同样应使用不变测度 \(\mathrm dt/t\)。

\Needspace{6\baselineskip}
\begin{proposition}[Mellin 卷积]\label{b3:prop:21-mellin-convolution}
设 \(c\in\mathbb R\)，可测函数 \(f,h\in L^1_c\)。则积分
\[
 (f*_M h)(x)\coloneqq\int_0^\infty f(t)h(x/t)\,\frac{\mathrm dt}{t}
\]
对几乎处处的 \(x>0\) 绝对收敛，定义 \(L^1_c\) 中的函数，并满足
\[
 \|f*_M h\|_{1,c}\le\|f\|_{1,c}\cdot\|h\|_{1,c}.
\]
在整条 \(\operatorname{Re}s=c\) 上，
\begin{equation}\label{b3:eq:21-mellin-convolution}
 \mathcal M(f*_M h)(s)=\mathcal Mf(s)\mathcal Mh(s).
\end{equation}
\end{proposition}
\begin{proof}
先对绝对值使用 Tonelli 定理，再在内部积分令 \(x=tu\)，得到
\[
 \begin{aligned}
 &\int_0^\infty\int_0^\infty
       |f(t)|\cdot|h(x/t)|\,\frac{\mathrm dt}{t}\,x^{c-1}\,\mathrm dx\\
 &=\int_0^\infty|f(t)|t^{c-1}\,\mathrm dt
        \int_0^\infty|h(u)|u^{c-1}\,\mathrm du<\infty.
 \end{aligned}
\]
这先给出卷积的几乎处处绝对存在性，再由三角不等式给出范数界。对 \(s=c+\mathrm iy\)，带复幂的相应二重积分具有同一个绝对值控制，故可以用 Fubini 交换积分。换元后 \(x^{s-1}\,\mathrm dx=t^su^{s-1}\,\mathrm du\)，外层的 \(\mathrm dt/t\) 正好将 \(t^s\) 变为 \(t^{s-1}\,\mathrm dt\)，便得到两个 Mellin 积分的乘积。
\end{proof}

这一卷积称为\term[mellin-juanji]{Mellin 卷积}[Mellin convolution]。若分别令 \(g_c(u)=e^{cu}f(e^u)\)、\(k_c(u)=e^{cu}h(e^u)\)，则直接换元给
\[
 e^{cu}(f*_Mh)(e^u)=(g_c*k_c)(u)
\]
几乎处处成立。乘法卷积就是对数坐标下的加法卷积；若把 \(\mathrm dt/t\) 错换成 \(\mathrm dt\)，权重和乘积公式都会随之改变。

\ProseParagraphBreak

移动 Gamma 反演的围道时，水平连接边扫过有限实部条带，因此需要在整个条带上统一控制高处的函数值。上一节特殊竖线的精确模值提示了衰减现象；从 Euler 积分出发的分部积分将给出移线所需的统一衰减。

\Needspace{6\baselineskip}
\begin{lemma}[Gamma 函数在有限条带的竖向衰减]\label{b3:lem:21-gamma-strip-decay}
设 \(a<b\) 为实数，\(k\ge0\) 为整数。存在常数 \(C_{a,b,k}\)，使 Gamma 函数的亚纯延拓满足
\[
 |\Gamma(\sigma+\mathrm it)|\le C_{a,b,k}(1+|t|)^{-k},
 \quad a\le\sigma\le b,\quad |t|\ge1.
\]
\end{lemma}
\begin{proof}
先设 \(a>0\)。Euler 积分令 \(x=e^u\)，得到
\[
 \Gamma(\sigma+\mathrm it)
 =\int_{\mathbb R}g_\sigma(u)e^{\mathrm itu}\,\mathrm du,
 \quad g_\sigma(u)=e^{\sigma u-e^u}.
\]
对每个固定 \(j\)，逐次求导可写成
\[
 g_\sigma^{(j)}(u)=g_\sigma(u)P_j(\sigma,e^u),
\]
其中 \(P_j\) 是多项式。事实上 \(P_0=1\)，递推为
\[
 P_{j+1}(\sigma,v)=(\sigma-v)P_j(\sigma,v)
                         +v\frac{\partial P_j}{\partial v}(\sigma,v).
\]
\(\sigma\in[a,b]\) 时系数统一有界。在 \(u\le0\) 上，\(e^u\le1\)，所以每个导数的模由 \(Ce^{au}\) 控制；在 \(u\ge0\) 上，可由 \(Ce^{(b+j)u-e^u}\) 控制。两端均可积，且各阶导数的边界值都为零。分部积分 \(k\) 次，便得到统一的 \(C|t|^{-k}\) 上界。

一般 \(a\) 下，选正整数 \(q\) 使 \(a+q>0\)。递推给
\[
 \Gamma(\sigma+\mathrm it)
 =\frac{\Gamma(\sigma+q+\mathrm it)}
         {\prod_{j=0}^{q-1}(\sigma+j+\mathrm it)}.
\]
当 \(|t|\ge1\) 时，每个分母因子的模至少为 \(|t|\)，分子已有正实部闭带上的任意固定幂次衰减。故所需估计也对整个给定条带成立。所有极点都在实轴上，当前的高处区域不遇到它们。
\end{proof}

\(e^{-x}\in L^1_c\) 对每个 \(c>0\) 成立，而引理使 \(t\mapsto\Gamma(c+\mathrm it)\) 绝对可积。因此由反演定理，
\begin{equation}\label{b3:eq:21-exponential-inversion}
 e^{-x}=\frac1{2\pi\mathrm i}
       \int_{c-\mathrm i\infty}^{c+\mathrm i\infty}\Gamma(s)x^{-s}\,\mathrm ds,
 \quad c>0,\quad x>0.
\end{equation}
这一步已经独立恢复原函数。下面的留数计算将进一步恢复它的 Taylor 级数，不能把两个步骤互相作为未经证明的依据。

\Needspace{6\baselineskip}
\begin{proposition}[指数函数反演的移线余项]\label{b3:prop:21-exponential-contour}
对每个整数 \(N\ge0\) 和实数 \(x>0\)，有
\[
 e^{-x}=\sum_{n=0}^N\frac{(-x)^n}{n!}+R_N(x),
\]
其中
\begin{equation}\label{b3:eq:21-exponential-remainder}
 R_N(x)=\frac1{2\pi\mathrm i}
       \int_{-N-1/2-\mathrm i\infty}^{-N-1/2+\mathrm i\infty}
                    \Gamma(s)x^{-s}\,\mathrm ds
\end{equation}
绝对收敛。存在与 \(N,x\) 无关的常数 \(C\)，使
\[
 |R_N(x)|\le C\frac{x^{N+1/2}}{\Gamma(N+3/2)}.
\]
因此余项在每个固定区间 \(0<x\le X<\infty\) 上一致趋于零。
\end{proposition}
\begin{proof}
在\eqref{b3:eq:21-exponential-inversion}中取 \(c=1\)。固定 \(N,x\)，考虑左右边实部分别为 \(-N-1/2\) 与 \(1\)、上下边虚部分别为 \(\pm T\) 的矩形，按正向绕行。其右边由下到上，左边由上到下。矩形内的极点恰为 \(0,-1,\ldots,-N\)，由\secref{b3:sec:2102}的留数计算，
\[
 \operatorname{Res}_{s=-n}\bigl(\Gamma(s)x^{-s}\bigr)
 =\frac{(-1)^nx^n}{n!}.
\]
对水平边使用\cref{b3:lem:21-gamma-strip-decay}，取 \(k=2\)。在这个固定实部区间，\(|x^{-s}|=x^{-\operatorname{Re}s}\) 有界，两条水平边的总长固定，所以其积分为 \(O_{N,x}(T^{-2})\)，趋于零。两个竖线积分都绝对收敛，故令 \(T\to\infty\)，留数定理给出所述有限和与余项，符号由上述边界方向确定。

为同时控制 \(N\)，将新竖线上的 Gamma 向右递推 \(N+1\) 次：
\[
 |\Gamma(-N-1/2+\mathrm it)|
 \le\frac{|\Gamma(1/2+\mathrm it)|}
              {\prod_{j=0}^N(j+1/2)}
 =\frac{\sqrt\pi}{\Gamma(N+3/2)}
                         |\Gamma(1/2+\mathrm it)|.
\]
上一节的精确模值 \(|\Gamma(1/2+\mathrm it)|^2=\pi/\cosh\pi t\) 表明右端的分子关于 \(t\) 可积，积分为一个固定常数。又新直线上 \(|x^{-s}|=x^{N+1/2}\)，所以得到所示统一余项界。

固定 \(X>0\)，令 \(a_N=X^{N+1/2}/\Gamma(N+3/2)\)。递推给
\[
 \frac{a_{N+1}}{a_N}=\frac{X}{N+3/2}\longrightarrow0,
\]
故 \(a_N\to0\)。对所有 \(0<x\le X\) 以 \(a_N\) 控制即可得到一致性。这一步甚至不需要再次使用 Stirling 公式。
\end{proof}

最后把移线估计中真正使用的条件抽出来。它们与函数是否叫作 Gamma 无关，也说明为什么仅有亚纯延拓还不够。

\Needspace{6\baselineskip}
\begin{proposition}[有限条带内的 Mellin 移线]\label{b3:prop:21-mellin-contour-shift}
设 \(a<b\) 为实数，\(F\) 在闭条带 \(a\le\operatorname{Re}s\le b\) 的某个开邻域亚纯，在该闭条带内的极点恰为有限个点 \(s_1,\ldots,s_m\)，且边界两条竖线上无极点。假设
\[
 \int_{\mathbb R}|F(a+\mathrm it)|\,\mathrm dt<\infty,
 \quad\int_{\mathbb R}|F(b+\mathrm it)|\,\mathrm dt<\infty,
\]
并存在 \(T_j\to\infty\)，使
\[
 \sup_{a\le\sigma\le b}|F(\sigma+\mathrm iT_j)|
 +\sup_{a\le\sigma\le b}|F(\sigma-\mathrm iT_j)|\longrightarrow0.
\]
则对每个 \(x>0\)，两个竖线积分绝对收敛，并满足
\begin{equation}\label{b3:eq:21-mellin-contour-shift}
 \frac1{2\pi\mathrm i}\int_{b-\mathrm i\infty}^{b+\mathrm i\infty}F(s)x^{-s}\,\mathrm ds
 =\sum_{\nu=1}^m\operatorname{Res}_{s=s_\nu}\bigl(F(s)x^{-s}\bigr)
  +\frac1{2\pi\mathrm i}\int_{a-\mathrm i\infty}^{a+\mathrm i\infty}F(s)x^{-s}\,\mathrm ds.
\end{equation}
若 \(F\) 在 \(s_0\) 的主要部分为 \(\sum_{k=1}^q A_k(s-s_0)^{-k}\)，则相应留数为
\[
 \operatorname{Res}_{s=s_0}\bigl(F(s)x^{-s}\bigr)
 =x^{-s_0}\sum_{k=1}^q\frac{A_k(-\log x)^{k-1}}{(k-1)!}.
\]
\end{proposition}
\begin{proof}
取上述两条竖线和高度 \(\pm T_j\) 的正向矩形。\(j\) 充分大时，它包含全部指定极点。水平边上 \(|x^{-s}|\) 至多为 \(\max\{x^{-a},x^{-b}\}\)，所以每条水平边的积分模至多为
\[
 (b-a)\max\{x^{-a},x^{-b}\}
           \sup_{a\le\sigma\le b}|F(\sigma\pm\mathrm iT_j)|,
\]
其极限为零。竖线上的绝对可积性使截断积分趋于整个直线积分。留数定理于是给出\eqref{b3:eq:21-mellin-contour-shift}。

在 \(s_0\) 附近展开
\[
 x^{-s}=x^{-s_0}\sum_{j=0}^{\infty}
                   \frac{(-\log x)^j(s-s_0)^j}{j!}.
\]
与有限主要部分相乘，选取 \((s-s_0)^{-1}\) 的系数，就得到最后的公式。
\end{proof}

例如 \(s=-\alpha\) 处的 \(q\) 阶极点贡献 \(x^\alpha\) 乘一个次数至多 \(q-1\) 的 \(\log x\) 多项式，与\cref{b3:prop:21-mellin-endpoint}中的正向计算相互核对。若把直线向左移到 \(a\)，新积分的模至多为
\[
 \frac{x^{-a}}{2\pi}\int_{\mathbb R}|F(a+\mathrm it)|\,\mathrm dt.
\]
这给 \(x\downarrow0\) 时可读出的余项尺度；向右移线则控制 \(x\to\infty\)，并应按\eqref{b3:eq:21-mellin-contour-shift}重新移项，保留留数符号。端点展开、亚纯极点和反演渐近由此组成一条有条件的双向联系，后续 theta 与 zeta 的积分方法会再次使用它。


\begin{chaptersummary}
Euler 积分先在右半平面定义 Gamma 函数，其全纯性、参数导数与递推都来自可以逐项核对的两端控制。正实轴上的严格对数凸性不仅排除非平凡周期因子，还给出固定平移的比值估计。Bohr--Mollerup 的条件共同确定函数，单独的阶乘插值或递推并不足够。

\ProseParagraphBreak
递推延拓与 Taylor 扣除描述同一个亚纯函数。全部非正整数都是单极点，留数决定奇部，有限部分进一步决定与零点相消后的导数。倒数乘积则说明 Gamma 没有复零点；其线性指数常数、对数导数与局部幂和都由 Euler 极限统一确定。

\ProseParagraphBreak
反射、Beta、倍增和乘法公式分别利用相反参数、积分换元与递推唯一性，规范化常数必须随这些步骤一起保留。Stirling 公式来自大参数下积分质量的集中；固定平移、增长平移与竖向参数是不同的极限问题，不能把一个区域中的渐近无条件移用到另一个区域。

\ProseParagraphBreak
Mellin 变换在对数坐标中成为带权 Fourier 变换，缩放成为乘子，乘法卷积成为普通乘积。收敛带记录原函数两端的可积性，亚纯延拓提供更多参数值，却不改变原积分的收敛范围。绝对反演、水平边消失和远处竖线余项趋零分别支持恢复函数、有限移线和最终留数展开；端点幂次决定极点位置，对数幂决定极点阶数。
\end{chaptersummary}

\BookExercises

% 自拟：由Gamma对数矩进入最优仿射逼近，利用协方差和正核表示求出严格剩余误差及其精确尺度。
\Needspace{6\baselineskip}
\begin{exercise}[对数函数的最优仿射逼近]\label{b3:ex:21-1}
对 \(x>0\)，用
\[
 \mathbb E_x[H(T)]=\frac1{\Gamma(x)}
 \int_0^\infty H(t)t^{x-1}e^{-t}\,\mathrm{d}t
\]
表示该正权重下的平均，并令 \(\psi_0=\Gamma'/\Gamma\)。在所有实常数 \(A,B\) 中，求使
\[
 \mathbb E_x\bigl[(\log T-AT-B)^2\bigr]
\]
最小的唯一一对 \((A,B)\)。证明最小值为
\(D(x)=\psi_0'(x)-1/x>0\)，并进一步证明
\[
 \lim_{x\to\infty}x^2D(x)=\frac12.
\]
这里允许使用本章的对数导数级数，但须说明最小值为何严格为正，以及最后一个极限的支配条件。
\end{exercise}
\begin{solution}
Euler 积分的递推与参数微分给
\[
 \mathbb E_x[T]=x,\quad
 \mathbb E_x[T^2]=x(x+1),\quad
 \mathbb E_x[\log T]=\psi_0(x),\quad
 \operatorname{Var}_x(\log T)=\psi_0'(x).
\]
这里方差表示 \(\mathbb E_x[(Y-\mathbb E_xY)^2]\)。所需混合矩同样绝对可积；原点附近的幂乘对数可积，无穷远仍有指数下降。于是
\[
 \mathbb E_x[T\log T]
 =\frac{\Gamma'(x+1)}{\Gamma(x)}
 =x\psi_0(x)+1,
\]
最后一步由 \(\Gamma'(x+1)=\Gamma(x)+x\Gamma'(x)\) 得到。因此 \(T\) 的方差为 \(x\)，而 \(T\) 与 \(\log T\) 的协方差为一。

把两变量中心化并配方，得到精确恒等式
\[
 \mathbb E_x\bigl[(\log T-AT-B)^2\bigr]
 =\psi_0'(x)-\frac1x
   +x\left(A-\frac1x\right)^2
   +\bigl(B+Ax-\psi_0(x)\bigr)^2.
\]
故唯一最优选择为
\(A=1/x\)、\(B=\psi_0(x)-1\)。若相应最小值为零，连续函数 \(\log t-t/x-B\) 必须在这个处处正的密度下几乎处处为零，继而在全部 \(t>0\) 为零。但其导数 \(1/t-1/x\) 不恒为零，矛盾，所以 \(D(x)>0\)。

为了求精确极限，由正项级数和 Tonelli 定理有
\[
 \psi_0'(x)=\sum_{n=0}^\infty\frac1{(x+n)^2}
 =\int_0^\infty\frac{u e^{-xu}}{1-e^{-u}}\,\mathrm{d}u.
\]
减去 \(1/x=\int_0^\infty e^{-xu}\,\mathrm{d}u\)，记
\(q(u)=u/(1-e^{-u})-1\)。对 \(u>0\)，不等式
\(1-e^{-u}<u\) 和 \(e^u>1+u\) 分别给 \(0<q(u)<u\)；Taylor 展开还给 \(q(u)=u/2+O(u^2)\)。令 \(v=xu\)，则
\[
 x^2D(x)=\int_0^\infty e^{-v}xq(v/x)\,\mathrm{d}v.
\]
被积函数逐点趋于 \(ve^{-v}/2\)，并由可积函数 \(ve^{-v}\) 控制。支配收敛遂给极限 \(\Gamma(2)/2=1/2\)。这个剩余误差衡量的是去掉最佳线性部分后的对数波动。
\end{solution}

% 自拟：把阻尼Gamma积分的计算与统一振荡尾估计相连，区分去阻尼极限和通常广义积分，并验证条件而非绝对收敛。
\Needspace{6\baselineskip}
\begin{exercise}[去阻尼与广义积分]\label{b3:ex:21-2}
设 \(b,c>0\)、\(\varepsilon>0\)，计算
\[
 I_\varepsilon(c,b)=\int_0^\infty
 x^{c-1}e^{-\varepsilon x}\cos(bx)\,\mathrm{d}x,
\]
并求 \(\varepsilon\downarrow0\) 的极限。证明当 \(0<c<1\) 时，通常的广义积分
\(\int_0^\infty x^{c-1}\cos(bx)\,\mathrm{d}x\) 条件收敛，且等于这一极限。再以 \(c=1\) 说明：去阻尼极限存在并不保证通常广义积分存在。
\end{exercise}
\begin{solution}
对 \(\operatorname{Re}w>0\)，函数
\[
 J(w)=\int_0^\infty x^{c-1}e^{-wx}\,\mathrm{d}x
\]
全纯；在每个紧参数集上，积分由 \(x^{c-1}e^{-\delta x}\) 控制，参数导数加入的幂也可积。正实 \(w\) 上换元给 \(J(w)=\Gamma(c)w^{-c}\)。在右半平面取正实轴上为实值的对数分支，恒等定理把该式推广到整个右半平面。令 \(w=\varepsilon-\mathrm{i}b\)，取实部，得
\[
 I_\varepsilon(c,b)
 =\Gamma(c)(\varepsilon^2+b^2)^{-c/2}
 \cos\left(c\arctan\frac b\varepsilon\right).
\]
所以对每个 \(c>0\)，去阻尼极限都为
\(\Gamma(c)b^{-c}\cos(\pi c/2)\)。

现在设 \(0<c<1\)。对 \(\varepsilon\ge0\)，振幅
\(a_\varepsilon(x)=x^{c-1}e^{-\varepsilon x}\) 在 \(x>0\) 单调递减并趋于零。分部积分给共同尾界
\[
 \left|\int_R^\infty a_\varepsilon(x)e^{\mathrm{i}bx}\,\mathrm{d}x\right|
 \le\frac{a_\varepsilon(R)+\int_R^\infty|a_\varepsilon'(x)|\,\mathrm{d}x}{b}
 =\frac{2a_\varepsilon(R)}b
 \le\frac{2R^{c-1}}b.
\]
当 \(\varepsilon=0\) 时，右侧同时证明振荡尾积分存在。在固定区间 \((0,R)\)，以 \(x^{c-1}\) 支配即可令 \(\varepsilon\downarrow0\)；再令 \(R\to\infty\)，上述统一尾界使两个极限相容。因此
\[
 \int_0^\infty x^{c-1}\cos(bx)\,\mathrm{d}x
 =\Gamma(c)b^{-c}\cos\frac{\pi c}2.
\]
这不是绝对收敛。围绕 \(2\pi k/b\) 取固定长度的小区间，使其中 \(|\cos(bx)|\ge1/2\)。第 \(k\) 个区间上的绝对积分不小于 \(Ck^{c-1}\)，而该正项级数发散。

最后取 \(c=1\)。已得公式给 \(I_\varepsilon(1,b)=\varepsilon/(\varepsilon^2+b^2)\to0\)，但
\(\int_0^R\cos(bx)\,\mathrm{d}x=\sin(bR)/b\) 没有极限。前面的统一尾界正是在 \(c=1\) 时失去趋零性质，因而不能把同一论证用到端点。
\end{solution}

% 自拟：将任意平移Gamma商的零极点分类与共同极点处的有限部分相连，辨明何时真正退化为有理函数。
\Needspace{6\baselineskip}
\begin{exercise}[Gamma 商的零极点与相消]\label{b3:ex:21-3}
设 \(a,b\in\mathbb C\)，把
\[
 Q(z)=\frac{\Gamma(z+a)}{\Gamma(z+b)}
\]
理解为亚纯函数。确定它的全部零点和极点，并分别判定它何时是有理函数、多项式或某个多项式的倒数。若某点 \(z_0\) 满足
\(z_0+a=-n\)、\(z_0+b=-r\)，其中 \(n,r\) 为非负整数，求 \(Q\) 在这个共同极点处可去延拓后的值与导数。
\end{exercise}
\begin{solution}
Gamma 没有零点，而其倒数在非正整数有且仅有单零点。因此若 \(b-a\notin\mathbb Z\)，两列点不相交，\(Q\) 的极点恰为 \(-a-n\)，零点恰为 \(-b-n\)，其中 \(n\ge0\)，且均为一阶。没有其他零极点。非零有理函数只有有限多个零极点，所以此时 \(Q\) 不是有理函数。

若 \(b-a=m\) 是正整数，有限次递推给
\[
 Q(z)=\frac1{\prod_{k=0}^{m-1}(z+a+k)}.
\]
故只剩 \(-a,-a-1,\ldots,-a-m+1\) 这些单极点，没有零点。若 \(a-b=m\) 是正整数，则
\[
 Q(z)=\prod_{k=0}^{m-1}(z+b+k),
\]
其零点为 \(-b,-b-1,\ldots,-b-m+1\)，均为单零点，没有极点。若 \(a=b\)，则 \(Q=1\)。这些恒等式先在避开极点的开集成立，再作为亚纯恒等式延拓，已经包含所有表观奇点的相消。

所以 \(Q\) 为有理函数当且仅当 \(a-b\in\mathbb Z\)；为多项式当且仅当 \(a-b\) 是非负整数；为非零多项式的倒数当且仅当 \(b-a\) 是非负整数。常函数一同时属于后两类。

最后令 \(w=z-z_0\)。记 \(H_0=0\)、\(H_j=\sum_{k=1}^j1/k\)。本章的极点有限部分给
\[
 \Gamma(-n+w)=\frac{(-1)^n}{n!w}
 \left(1+(H_n-\gamma)w+O(w^2)\right),
\]
分母有把 \(n\) 换为 \(r\) 的同类展开。两式相除得
\[
 Q(z_0+w)=(-1)^{n-r}\frac{r!}{n!}
 \left(1+(H_n-H_r)w+O(w^2)\right).
\]
因此
\[
 Q(z_0)=(-1)^{n-r}\frac{r!}{n!},\quad
 Q'(z_0)=(-1)^{n-r}\frac{r!}{n!}(H_n-H_r).
\]
值由留数之比确定，导数还需要两边有限部分的差；形式上的“两个无穷大相除”不能给出这些资料。
\end{solution}

% 自拟：构造保持递推、规范化与严格普通凸性的非平凡周期扰动，说明Bohr--Mollerup中的对数凸性不可直接弱化。
\Needspace{6\baselineskip}
\begin{exercise}[普通凸性不能代替对数凸性]\label{b3:ex:21-4}
对实数 \(\varepsilon\)，令
\[
 F_\varepsilon(x)=\Gamma(x)e^{\varepsilon\sin2\pi x},\quad x>0.
\]
证明存在 \(\varepsilon_0>0\)，使每个 \(0<|\varepsilon|\le\varepsilon_0\) 对应的 \(F_\varepsilon\) 都正且严格凸，满足
\(F_\varepsilon(1)=1\) 与 \(F_\varepsilon(x+1)=xF_\varepsilon(x)\)，却不等于 \(\Gamma\)。再证明只要 \(\varepsilon\ne0\)，\(\log F_\varepsilon\) 就不是凸函数。无需给出 \(\varepsilon_0\) 的数值，但须证明一个对所有 \(x>0\) 有效的选择确实存在。
\end{exercise}
\begin{solution}
正性、规范化与递推直接来自指数因子的一周期性。若 \(\varepsilon\ne0\)，取 \(x=1/4\) 即知这个因子不等于一，所以 \(F_\varepsilon\ne\Gamma\)。问题在于选择一个足够小而对全半轴适用的扰动。

置 \(\omega=2\pi\)，并在本解中记 \(\psi=\psi_0\)。对数导数级数给 \(\psi'(x)>0\)；原点展开给 \(\psi(x)=-1/x+O(1)\)，而
\(\psi(n)=H_{n-1}-\gamma\) 与单调性给 \(\psi(x)\to\infty\)。因此 \(|\psi(x)|\) 在 \(x\downarrow0\) 和 \(x\to\infty\) 两端都趋于无穷。

定义连续正函数
\[
 q(x)=\psi'(x)+\psi(x)^2,\quad
 r(x)=2\omega|\psi(x)|+\omega^2.
\]
有
\[
 \frac{q(x)}{r(x)}
 \ge\frac{\psi(x)^2}{2\omega|\psi(x)|+\omega^2}
 \longrightarrow\infty
\]
于两端。在中间紧区间上，这个连续正商取得正的最小值，故
\(\mu=\inf_{x>0}q(x)/r(x)>0\)。取 \(\varepsilon_0=\mu/2\)。逐次求导得
\[
 \frac{F_\varepsilon''}{F_\varepsilon}
 =q+2\varepsilon\omega\psi\cos(\omega x)
   -\varepsilon\omega^2\sin(\omega x)
   +\varepsilon^2\omega^2\cos^2(\omega x)
 \ge q-|\varepsilon|r\ge\frac q2>0.
\]
因此所有这样的 \(F_\varepsilon\) 都严格凸。

另一方面，
\[
 (\log F_\varepsilon)''(x)
 =\psi'(x)-\varepsilon\omega^2\sin(\omega x).
\]
本章的级数积分比较给 \(\psi'(x)\to0\)。当 \(\varepsilon>0\) 时取 \(x_n=n+1/4\)，当 \(\varepsilon<0\) 时取 \(x_n=n+3/4\)，两种情形都有
\[
 (\log F_\varepsilon)''(x_n)
 =\psi'(x_n)-|\varepsilon|\omega^2<0
\]
于充分大的 \(n\)。所以对数不凸。普通凸性可以由一阶对数导数的平方维持，而对数凸性只看二阶对数导数；这正是本例中两个条件出现不同结论的原因。
\end{solution}

% 自拟：给同一有理亚纯函数选取三个不同反演竖线，求出各自原函数与实际收敛带，显示移线留数的作用。
\Needspace{6\baselineskip}
\begin{exercise}[同一表达式的三个反演条带]\label{b3:ex:21-5}
设 \(\alpha<\beta\) 为实数，令
\[
 F(s)=\frac1{(s-\alpha)(\beta-s)}.
\]
对 \(x>0\) 及 \(c\in\mathbb R\setminus\{\alpha,\beta\}\)，证明
\[
 I_c(x)=\frac1{2\pi\mathrm{i}}
 \int_{c-\mathrm{i}\infty}^{c+\mathrm{i}\infty}F(s)x^{-s}\,\mathrm{d}s
\]
绝对收敛，并分别计算 \(c<\alpha\)、\(\alpha<c<\beta\) 和 \(c>\beta\) 时的函数。求这三个函数各自的最大绝对收敛带，用留数解释相邻两种答案的差，并说明为什么它们不违背 Mellin 变换的唯一性。
\end{exercise}
\begin{solution}
在任意不经过极点的竖线上，\(F(c+\mathrm{i}t)=O((1+|t|)^{-2})\)，而 \(|x^{-c-\mathrm{i}t}|=x^{-c}\)，故所示积分绝对收敛。记 \(d=\beta-\alpha>0\)，考虑三个连续函数
\[
 f_0(x)=\frac1d
 \begin{cases}
 x^{-\alpha},&0<x\le1,\\
 x^{-\beta},&x>1,
 \end{cases}
\]
以及
\[
 f_-(x)=\frac1d
 \begin{cases}
 0,&0<x\le1,\\
 x^{-\beta}-x^{-\alpha},&x>1,
 \end{cases}
 \quad
 f_+(x)=\frac1d
 \begin{cases}
 x^{-\alpha}-x^{-\beta},&0<x<1,\\
 0,&x\ge1.
 \end{cases}
\]
对 \(f_0\)，零端要求 \(\operatorname{Re}s>\alpha\)，无穷远要求 \(\operatorname{Re}s<\beta\)，且
\[
 \mathcal Mf_0(s)=\frac1d
 \left(\frac1{s-\alpha}+\frac1{\beta-s}\right)=F(s).
\]
对 \(f_-\)，只有无穷远端需要检查，那里
\(f_-(x)\sim-x^{-\alpha}/d\)。其最大绝对收敛域为 \(\operatorname{Re}s<\alpha\)，在该域逐项积分给
\[
 \mathcal Mf_-(s)=\frac1d
 \left(\frac1{\beta-s}-\frac1{\alpha-s}\right)=F(s).
\]
同样，\(f_+(x)\sim-x^{-\beta}/d\) 于零端，其最大绝对收敛域为 \(\operatorname{Re}s>\beta\)，并且
\[
 \mathcal Mf_+(s)=\frac1d
 \left(\frac1{s-\alpha}-\frac1{s-\beta}\right)=F(s).
\]
在各边界实部，绝对积分都含非零常数倍的 \(\int\mathrm{d}x/x\)，所以不能把边界加入。三个函数都在各自竖线的权重下可积，变换沿竖线也可积；Mellin 绝对反演因此处处给
\[
 I_c=
 \begin{cases}
 f_-,&c<\alpha,\\
 f_0,&\alpha<c<\beta,\\
 f_+,&c>\beta.
 \end{cases}
\]
在 \(x=1\) 时，中间答案为 \(1/d\)，两侧答案为零，这些值也由各函数的连续性确定。

反演被积函数在两个极点的留数为
\[
 \operatorname{Res}_{s=\alpha}\bigl(F(s)x^{-s}\bigr)=\frac{x^{-\alpha}}d,
 \quad
 \operatorname{Res}_{s=\beta}\bigl(F(s)x^{-s}\bigr)=-\frac{x^{-\beta}}d.
\]
由下往上的竖线向右移时，右线积分等于左线积分加所越过极点的留数。因此
\[
 f_0-f_-=\frac{x^{-\alpha}}d,\quad
 f_+-f_0=-\frac{x^{-\beta}}d.
\]
有限水平区间上 \(F(\sigma+\mathrm{i}T)=O(T^{-2})\)，且 \(x^{-\sigma}\) 一致有界，所以矩形的水平边确实趋于零。唯一性要求两个原函数在同一权重下可积，并在共同竖线上有相同变换；这里三个最大收敛带互不相交，不满足这个前提。只保留亚纯表达式而忘记收敛带，会丢失原函数信息。
\end{solution}

% 自拟：以高阶有限差分制造原点抵消，计算被消去的Gamma极点处的值，并用正积分证明对应零点确为单零点。
\Needspace{6\baselineskip}
\begin{exercise}[高阶抵消与可去值]\label{b3:ex:21-6}
设 \(a>0\)，\(m\ge1\) 为整数，令
\[
 f_m(x)=e^{-ax}(1-e^{-x})^m,\quad
 A_m(s)=\sum_{k=0}^m(-1)^k\binom mk(a+k)^{-s},
\]
其中幂使用正实数的实对数。求 \(f_m\) 的最大 Mellin 绝对收敛域，并证明其中
\[
 \mathcal Mf_m(s)=\Gamma(s)A_m(s)
\]
按可去延拓解释。计算 \(s=-j\)、\(0\le j<m\) 处的值，证明这些点都是 \(A_m\) 的恰好一重零点；最后求右端在 \(s=-m\) 的留数。
\end{exercise}
\begin{solution}
在原点附近，\(f_m(x)=x^m(1+O(x))\)；无穷远则由 \(e^{-ax}\) 控制。因此最大绝对收敛域为 \(\operatorname{Re}s>-m\)，边界实部已有 \(x^{-1}\) 型发散。在这个半平面上，Mellin 积分定义全纯函数，记为 \(G_m\)。

在较小的半平面 \(\operatorname{Re}s>0\)，有限二项展开后的每项都可以单独积分，所以
\[
 G_m(s)=\sum_{k=0}^m(-1)^k\binom mk
 \int_0^\infty e^{-(a+k)x}x^{s-1}\,\mathrm{d}x
 =\Gamma(s)A_m(s).
\]
\(A_m\) 是整函数。要核对前 \(m\) 个 Gamma 极点的抵消，使用有限恒等式
\[
 \sum_{k=0}^m(-1)^k\binom mk e^{(a+k)t}
 =e^{at}(1-e^t)^m=(-1)^mt^m+O(t^{m+1}).
\]
在原点逐次求导便得
\[
 \sum_{k=0}^m(-1)^k\binom mk(a+k)^j=0\quad(0\le j<m),
\]
而 \(j=m\) 时该和等于 \((-1)^m m!\)。因此
\(A_m(-j)=0\) 对 \(0\le j<m\) 成立，足以消去这些点的单极点。恒等定理遂把 \(G_m=\Gamma A_m\) 推广到整个 \(\operatorname{Re}s>-m\)。

在 \(s=-j\) 附近，Gamma 的留数为 \((-1)^j/j!\)，故可去值为
\[
 G_m(-j)=\frac{(-1)^j}{j!}A_m'(-j)
 =\frac{(-1)^{j+1}}{j!}
 \sum_{k=0}^m(-1)^k\binom mk(a+k)^j\log(a+k).
\]
这里 \(-j>-m\)，原来的正积分仍收敛，而且 \(f_m(x)>0\)，所以 \(G_m(-j)>0\)。于是 \(A_m'(-j)\ne0\)，证明这些零点恰为一重。零点阶数不能仅从前面的有限差分恒等式直接读出，正积分排除了更高阶相消。

在 \(s=-m\)，已经算得 \(A_m(-m)=(-1)^m m!\ne0\)，因此右端仍有单极点，留数为
\[
 \frac{(-1)^m}{m!}A_m(-m)=1.
\]
它与原点展开的首项系数一相符。有限差分消除了前 \(m\) 个极点，也恰好把原积分的左收敛边界推进了 \(m\) 个单位。
\end{solution}

% 自拟：从不同指数截断幂的乘法卷积导出部分分式公式，再通过紧单纯形积分解释参数合并时出现的对数因子。
\Needspace{6\baselineskip}
\begin{exercise}[不同指数的卷积与合并极限]\label{b3:ex:21-7}
设 \(m\ge2\)，\(a_1,\ldots,a_m\) 为实数，并定义
\[
 h_a(x)=
 \begin{cases}
 x^a,&0<x<1,\\
 0,&x\ge1,
 \end{cases}
 \quad H=h_{a_1}*_M\cdots*_Mh_{a_m},
\]
其中 \((f*_Mg)(x)=\int_0^\infty f(t)g(x/t)\,\mathrm{d}t/t\)。证明当参数两两不同时，
\[
 H(x)=
 \begin{cases}
 \displaystyle\sum_{j=1}^m
 \frac{x^{a_j}}{\prod_{k\ne j}(a_k-a_j)},&0<x<1,\\
 0,&x\ge1.
 \end{cases}
\]
再证明当所有 \(a_j\to a\) 时，这个表达式在 \((0,\infty)\) 的紧集上一致趋于
\[
 \frac{x^a(-\log x)^{m-1}}{(m-1)!}\cdot\mathbf1_{(0,1)}(x).
\]
求两种情形的 Mellin 变换与最大绝对收敛域，并说明系数分母趋零为什么不导致卷积发散。
\end{exercise}
\begin{solution}
先用对数坐标写出不含相消的积分。令 \(x=e^{-u}\)。当 \(u>0\) 时，多次卷积和换元给
\[
 H(e^{-u})=u^{m-1}\int_{\Delta}
 \exp\left[-u\left(\sum_{j=1}^{m-1}a_jw_j
              +a_m\left(1-\sum_{j=1}^{m-1}w_j\right)\right)\right]
 \,\mathrm{d}w_1\cdots\mathrm{d}w_{m-1},
\]
其中
\[
 \Delta=\left\{\,(w_1,\ldots,w_{m-1})\in[0,\infty)^{m-1}
          \ \middle|\ \sum_{j=1}^{m-1}w_j\le1\,\right\}.
\]
具体地，两因子卷积在此坐标变成 \(\int_0^u\) 上的加法卷积；归纳后令各积分变量为 \(uw_j\)，就得到上式。所有被积函数非负，Tonelli 定理保证迭代次序合法。若 \(x\ge1\)，支集条件没有正长度的积分区域，故 \(H(x)=0\)。这个紧单纯形表示还证明 \(H\) 连续，且因 \(m\ge2\)，在 \(x=1\) 的值为零。

设 \(a_* =\min_j a_j\)。在 \(\operatorname{Re}s>-a_*\)，各因子都在同一权重下可积，卷积定理给
\[
 \mathcal MH(s)=\prod_{j=1}^m\frac1{s+a_j}.
\]
当参数两两不同时，部分分式展开为
\[
 \prod_{j=1}^m\frac1{s+a_j}
 =\sum_{j=1}^m\frac{c_j}{s+a_j},\quad
 c_j=\frac1{\prod_{k\ne j}(a_k-a_j)}.
\]
系数由乘以 \(s+a_j\) 后令 \(s=-a_j\) 得到。每项正是 \(c_jh_{a_j}\) 的 Mellin 变换，因此题给候选式与 \(H\) 在共同竖线上有相同变换。该有理函数沿竖线为 \(O(|t|^{-m})\)，属于 \(L^1\)。此外 \(\sum_jc_j=0\)，因为在部分分式两侧令 \(s\to+\infty\) 后比较 \(1/s\) 项即可；故候选式在 \(x=1\) 也连续。对两个连续函数使用绝对反演，便知它们处处相等。

若最小参数为 \(a_*=a_{j_*}\)，其系数
\(c_{j_*}=\prod_{k\ne j_*}(a_k-a_*)^{-1}>0\)，其余幂次都较大。因此
\(H(x)\sim c_{j_*}x^{a_*}\) 于零端，最大绝对收敛域确为 \(\operatorname{Re}s>-a_*\)；无穷远没有贡献。

现在让所有 \(a_j\to a\)。紧单纯形上的被积函数在每个有界 \(u\ge0\) 区间上一致趋于 \(e^{-au}\)。单纯形的体积为 \(1/(m-1)!\)，可由逐次积分归纳得到，因而
\[
 H(e^{-u})\longrightarrow
 \frac{e^{-au}u^{m-1}}{(m-1)!}
\]
在这些区间上一致成立。任意 \(x>0\) 的紧区间都只对应有界的 \(u\) 区间；\(x\ge1\) 时又恒为零，所以得到题中的局部一致极限。在参数全部等于 \(a\) 时，上式本身就是精确公式，且
\[
 \mathcal MH(s)=\frac1{(s+a)^m},\quad
 \operatorname{Re}s>-a.
\]
零端的 \(x^a(-\log x)^{m-1}\) 说明这个收敛域同样最大。分母 \(\prod_{k\ne j}(a_k-a_j)\) 的趋零只是部分分式表达中的相消变得显著；原卷积的紧区域积分始终有界并连续依赖参数。简单极点合并成高阶极点，原函数中相应出现对数幂。
\end{solution}

% 自拟：把固定平移渐近推广到增长参数的一致区间，计算二次指数修正的有效尺度，并在临界尺度保留三次项。
\Needspace{6\baselineskip}
\begin{exercise}[增长平移的有效尺度]\label{b3:ex:21-8}
对实数 \(x>0\)、\(h\) 满足 \(x+h>0\)，令
\[
 R(x,h)=\frac{\Gamma(x+h)}{x^h\Gamma(x)}.
\]
设 \(L(x)\ge0\) 且 \(L(x)=o(x^{2/3})\)。证明
\[
 \sup_{|h|\le L(x)}
 \left|\log R(x,h)-\frac{h^2}{2x}\right|
 \longrightarrow0\quad(x\to\infty).
\]
由此分别讨论 \(h=o(\sqrt x)\)、\(h/\sqrt x\to\lambda\in\mathbb R\)，以及
\(\sqrt x=o(|h|)\)、\(|h|=o(x^{2/3})\) 时的 \(R(x,h)\)。最后固定 \(c\in\mathbb R\setminus\{0\}\)，令 \(h=cx^{2/3}\)，计算
\[
 \lim_{x\to\infty}e^{-h^2/(2x)}R(x,h),
\]
说明上面的有效范围为什么不能直接扩大为 \(|h|=O(x^{2/3})\)。
\end{exercise}
\begin{solution}
写正实轴 Stirling 公式的对数形式
\[
 \log\Gamma(t)=\left(t-\frac12\right)\log t-t
                 +\frac12\log(2\pi)+r(t),\quad r(t)\to0.
\]
对 \(|h|\le L(x)\)，充分大的 \(x\) 满足 \(|h|\le x/2\)，于是 \(x+h\ge x/2>0\)。两次应用此式并约去共同项，得到
\[
 \log R(x,h)=\left(x+h-\frac12\right)
             \log\left(1+\frac hx\right)-h
             +r(x+h)-r(x).
\]
由于
\(\sup_{t\ge x/2}|r(t)|\to0\)，最后两项之差对整个 \(h\) 区间一致趋于零。这一步仅使用 Stirling 余项趋零，不对余项作未经证明的微分。

在 \(|u|\le1/2\) 上，Taylor 余项统一满足
\(\log(1+u)=u-u^2/2+O(|u|^3)\)。代入 \(u=h/x\)，整理得到
\[
 \left(x+h-\frac12\right)\log\left(1+\frac hx\right)-h
 =\frac{h^2}{2x}
   +O\left(\frac{|h|}x+\frac{|h|^3}{x^2}\right),
\]
其中常数不依赖于 \(x,h\)。所以所求上确界不超过
\[
 C\left(\frac{L(x)}x+\frac{L(x)^3}{x^2}\right)
 +2\sup_{t\ge x/2}|r(t)|,
\]
这在 \(L=o(x^{2/3})\) 时趋于零。

因而若 \(h=o(\sqrt x)\)，则 \(R(x,h)\to1\)；若 \(h/\sqrt x\to\lambda\)，则
\(R(x,h)\to e^{\lambda^2/2}\)。特别地，\(h=\sqrt x\) 时极限为 \(e^{1/2}\)。若
\(\sqrt x=o(|h|)\) 且 \(|h|=o(x^{2/3})\)，则 \(h^2/(2x)\to\infty\)，故 \(R(x,h)\to\infty\)。这些结论对正、负平移都适用，只要保持 \(x+h>0\)。

最后取 \(h=cx^{2/3}\)。这时必须把对数保留到三次：
\[
 \log(1+u)=u-\frac{u^2}2+\frac{u^3}3+O(u^4).
\]
代入同一恒等式，有
\[
 \log R(x,h)-\frac{h^2}{2x}
 =-\frac{h}{2x}-\frac{h^3}{6x^2}
   +O\left(\frac{h^2}{x^2}+\frac{|h|^4}{x^3}\right)+o(1)
 \longrightarrow-\frac{c^3}6.
\]
于是所求极限为 \(e^{-c^3/6}\ne1\)。在临界尺度，三次项保留了非零贡献；固定平移公式中可以忽略的项，在增长参数下必须重新估计。
\end{solution}

